% arXiv version

\documentclass[11pt,a4paper]{article}
\usepackage{amsmath,amssymb,amsthm,enumerate,mathrsfs,bbm,xcolor}
\usepackage[english]{babel}

\usepackage[top=2.5cm, bottom=3.5cm, left=2.5cm, right=2.5cm]{geometry}

\usepackage[hidelinks]{hyperref}

%%%%%%%%%% Start TeXmacs macros
\newcommand{\TeXmacs}{T\kern-.1667em\lower.5ex\hbox{E}\kern-.125emX\kern-.1em\lower.5ex\hbox{\textsc{m\kern-.05ema\kern-.125emc\kern-.05ems}}}
\newcommand{\cdummy}{\cdot}
\newcommand{\comma}{{,}}
\newcommand{\mathD}{\mathrm{D}}
\newcommand{\mathd}{\mathrm{d}}
\newcommand{\nobracket}{}
\newcommand{\nocomma}{}
\newcommand{\nosymbol}{}
\newcommand{\precprec}{\prec\!\!\!\prec}
\newcommand{\tmSep}{; }
\newcommand{\tmaffiliation}[1]{\thanks{\textit{Affiliation:} #1}}
\newcommand{\tmem}[1]{{\em #1\/}}
\newcommand{\tmemail}[1]{\thanks{\textit{Email:} \texttt{#1}}}
\newcommand{\tmnote}[1]{\thanks{\textit{Note:} #1}}
\newcommand{\tmop}[1]{\ensuremath{\operatorname{#1}}}
\newcommand{\tmsep}{, }
\newcommand{\tmtextit}[1]{{\itshape{#1}}}
\newcommand{\tmkeywords}{\textbf{Keywords:} }
%/newcommand{\upl}{+}
\newcommand{\withTeXmacstext}{This document has been produced using the GNU {\TeXmacs} text editor (see \url{http://www.texmacs.org})}
\newenvironment{enumeratenumeric}{\begin{enumerate}[1.] }{\end{enumerate}}
\newenvironment{enumerateroman}{\begin{enumerate}[i.] }{\end{enumerate}}
\newenvironment{itemizeminus}{\begin{itemize} \renewcommand{\labelitemi}{$-$}\renewcommand{\labelitemii}{$-$}\renewcommand{\labelitemiii}{$-$}\renewcommand{\labelitemiv}{$-$}}{\end{itemize}}
\newtheorem{theorem}{Theorem}
\newtheorem{corollary}[theorem]{Corollary}
\newtheorem{lemma}[theorem]{Lemma}
\theoremstyle{remark}
\newtheorem{remark}[theorem]{Remark}
%{\theorembodyfont{\rmfamily}\newtheorem{remark}{Remark}}
%%%%%%%%%% End TeXmacs macros



\newcommand{\R}{\mathbb{R}}
\newcommand{\N}{\mathbb{N}}
\newcommand{\Z}{\mathbb{Z}}
\newcommand{\E}{\mathbb{E}}

\newcommand{\Lcal}{\mathcal{L}}


%\renewcommand{\vartheta}{\theta}





\newcommand{\CA}{\ensuremath{\mathscr{A}}}
\newcommand{\CB}{\ensuremath{\mathscr{B}}}
\newcommand{\CC}{\ensuremath{\mathscr{C}}}
\newcommand{\CD}{\ensuremath{\mathscr{D}}}
\newcommand{\CE}{\ensuremath{\mathscr{E}}}
\newcommand{\CF}{\ensuremath{\mathscr{F}}}
\newcommand{\CG}{\ensuremath{\mathscr{G}}}
\newcommand{\CH}{\ensuremath{\mathscr{H}}}
\newcommand{\CI}{\ensuremath{\mathscr{I}}}
\newcommand{\CJ}{\ensuremath{\mathscr{J}}}
\newcommand{\CK}{\ensuremath{\mathscr{K}}}
\newcommand{\CL}{\ensuremath{\mathscr{L}}}
\newcommand{\CM}{\ensuremath{\mathscr{M}}}
\newcommand{\CN}{\ensuremath{\mathscr{N}}}
\newcommand{\CO}{\ensuremath{\mathscr{O}}}
\newcommand{\CP}{\ensuremath{\mathscr{P}}}
\newcommand{\CQ}{\ensuremath{\mathscr{Q}}}
\newcommand{\CR}{\ensuremath{\mathscr{R}}}
\newcommand{\CS}{\ensuremath{\mathscr{S}}}
\newcommand{\CT}{\ensuremath{\mathscr{T}}}
\newcommand{\CU}{\ensuremath{\mathscr{U}}}
\newcommand{\CV}{\ensuremath{\mathscr{V}}}
\newcommand{\CZ}{\ensuremath{\mathscr{Z}}}


\begin{document}



\title{Paracontrolled distributions and singular PDEs \tmnote{{\withTeXmacstext}}}

\author{
  Massimiliano Gubinelli \\ CEREMADE \& CNRS UMR 7534 \\ Universit{\'e} Paris-Dauphine\\
%  \tmaffiliation{CEREMADE \& CNRS UMR 7534{\tmSep}
%  Universit{\'e} Paris-Dauphine}
  \texttt{gubinelli@ceremade.dauphine.fr}
  \and
  Peter Imkeller \\ Institut f\"ur Mathematik \\ Humboldt-Universit\"at zu Berlin\\
%  \tmaffiliation{Institut f{\"u}r Mathematik{\tmSep}
%  Humboldt-Universit{\"a}t zu Berlin}
  \texttt{imkeller@math.hu-berlin.de}
  \and
  Nicolas Perkowski \\  CEREMADE \& CNRS UMR 7534  \\ Universit{\'e} Paris-Dauphine\\
%  \tmaffiliation{Institut f{\"u}r Mathematik{\tmSep}
%  Humboldt-Universit{\"a}t zu Berlin}
  \texttt{perkowski@ceremade.dauphine.fr}
}

%\date{\tmdate.}

\maketitle

\begin{abstract}
   We introduce an approach to study certain singular PDEs which is based on  techniques from paradifferential calculus and on ideas from the theory of controlled rough paths. We illustrate its applicability on some model problems like differential equations driven by fractional Brownian motion, a fractional Burgers type SPDE driven by space-time white noise, and a non-linear version of the parabolic Anderson model with a white noise potential.
\end{abstract}

\tmkeywords{Rough paths}{\tmsep}{Paraproducts}{\tmsep}{Besov spaces}{\tmsep}{Stochastic partial differential equations}{\tmsep}{Renormalization}

%\tableofcontents

\section{Introduction}

In this paper we introduce the notion of \emph{paracontrolled distribution} and show how to use it to give a meaning to and solve partial differential equations involving non-linear operations on generalized functions.  More precisely, we combine the idea of \emph{controlled paths}, introduced  in~\cite{Gubinelli2004}, with the \emph{paraproduct} introduced by Bony~\cite{Bony1981} and the related paradifferential calculus, in order to develop a non-linear theory for a certain class of distributions.

The approach presented here works for generalized functions defined on an index set of arbitrary dimension and constitutes a flexible and lightweight generalization of Lyons' rough path theory~\cite{Lyons1998}. In particular it allows to handle problems involving singular stochastic PDEs which were substantially out of reach with previously known methods.

In order to set the stage for our analysis let us list some of the problems which are amenable to be analyzed in the paracontrolled framework:
\begin{enumerate}%[1.]
  \item The \emph{rough} differential equation (\textsc{rde}) driven by a $d$--dimensional Gaussian process $X$: 
  \[
     \partial_t u(t) = F(u(t)) \partial_t X(t),%\qquad u(0)=u_0,
  \]
  where $F\colon\R^n \rightarrow \Lcal(\R^d, \R^n)$ is a smooth vector-field . Typically, $X$ will be a Brownian motion or a fractional Brownian motion with Hurst exponent $H\in(0,1)$. The paracontrolled analysis works up to $H>1/3$.  While we do not have any substantial new results for this problem, it is a useful pedagogical example on which we can easily describe our approach.
  
  \item Generalizations of Hairer's Burgers-like SPDE (\textsc{burgers}):
  \[
     L u = G(u) \partial_x u + \xi.
  \]
  Here $u\colon \R_+ \times \mathbb{T} \rightarrow \R^n$, where $\mathbb{T} = (\R/2\pi \Z)$ denotes the torus, $L = \partial_t + (-\Delta)^{\sigma}$, where $-(-\Delta)^\sigma$ is the fractional Laplacian with periodic boundary conditions and we will take $\sigma \in > 5/6$, and $\xi$ is a space-time white noise with values in $\R^n$. Moreover, $G\colon \R^n \rightarrow \Lcal (\R^n, \R^n)$ is a smooth field of linear transformations.
  
  \item A non-linear generalization of the parabolic Anderson model (\textsc{pam}):
  \[
     L u = F(u) \diamond \xi,
  \]
  where $u \colon \R_+ \times \mathbb{T}^2 \rightarrow \R$, $L = \partial_t - \Delta$ is the parabolic operator corresponding to the heat equation, and where $\xi$ is a random potential which is sampled according to the law of the white noise on  $\mathbb{T}^2$ and is therefore independent of the time variable. We allow for a general smooth function $F\colon \R \rightarrow \R$, the linear case $F(u)=u$ corresponding to the standard parabolic Anderson model. The symbol $\diamond$ stands for a renormalized product which is necessary to have a well defined problem. 
  
  \item The one-dimensional periodic Kardar--Parisi--Zhang equation (\textsc{kpz}):
  \[
   L h = ``(\partial_x h)^2" + \xi,
  \]
    where $u \colon \R_+ \times \mathbb{T} \rightarrow \R$, $L = \partial_t - \Delta$, and where $\xi$ is a space-time white noise. Here $``(\partial_x h)^2"$ denotes the necessity of an additive renormalization in the definition of the square of the distribution $\partial_x h$.

  \item The three-dimensional, periodic,  stochastic quantization equation for the $(\phi)^4_3$ euclidean quantum field (\textsc{sq}):
  \[
   L \phi = ``\frac{\lambda}{4!} (\phi)^3" + \xi,
  \]
    where $\phi \colon \R_+ \times \mathbb{T}^3 \rightarrow \R$, $L = \partial_t - \Delta$, $\xi$ is a space-time white noise, and where $``(\phi)^3"$ denotes a suitable renormalization of a cubic polynomial of $\phi$ and $\lambda$ is the coupling constant of the scalar theory.
  
\end{enumerate}

In this paper we will consider in detail the three cases~\textsc{rde}, \textsc{burgers}, \textsc{pam}. In all cases we will exhibit a space of paracontrolled distributions where the equations are well posed (in a suitable sense), and admit at least a local solution which is unique.
%
The three-dimensional stochastic quantization equation \textsc{sq} is studied by R.~Catellier and K.~Chouk in~\cite{Catellier2013} by applying the paracontrolled technique. The paracontrolled analysis of \textsc{kpz}  will be presented elsewhere~\cite{Gubinelli-KPZ-2013}.


\bigskip

The kind of results which will be obtained below can be exemplified by the following statement for \textsc{rde}s. Below $\CC^\alpha = B^\alpha_{\infty, \infty}$ stands for the H\"older-Besov space of index $\alpha$ on $\mathbb{R}$. Given two distributions $f\in \CC^\alpha$ and $g\in \CC^\beta$ with $\alpha+\beta >0$ we can always consider a certain distribution $f \circ g$ which is obtained via a bilinear operation of $f,g$ and which belongs to $\CC^{\alpha+\beta}$.
\begin{theorem}
Let $\xi:[0,1]\to \mathbb{R}^n$ be a continuous function and $F:\mathbb{R}^d\to \mathcal{L}(\mathbb{R}^n, \mathbb{R}^d)$ be a family of  smooth vector-fields. Let $u:[0,1]\to \mathbb{R}^d$ be a 
 solution of the Cauchy problem
\[
\partial_t u(t) = F(u(t)) \xi(t),\qquad
u(0)=u_0,
\]
where $u_0 \in \R^d$. Let $\vartheta$ be a solution to $\partial_t \vartheta = \xi$ and let $R\xi=(\xi,\vartheta\circ \xi)$. Then for all $\alpha \in (1/3,1)$ there exists a continuous map $\Psi:\mathbb{R}^d\times\CC^{\alpha-1}\times\CC^{2\alpha-1}\to \CC^{\alpha}$ such that $u=\Psi(u_0,R\xi)$ for all $\xi\in C([0,1];\mathbb{R}^d)$.
\end{theorem}

In particular, this theorem provides a natural way of extending the solution map to data $\xi$ which are merely distributions in $\CC^{\alpha-1}$. It suffices to approximate $\xi$ by a sequence of smooth functions $(\xi^n)$ converging to $\xi$ in $\CC^{\alpha-1}$, and to prove that the ``lifted" sequence $(R\xi^n)$ converges to some limit in $ \CC^{\alpha-1}\times\CC^{2\alpha-1}$. The uniqueness of this limit is not guaranteed however, and each possible limit will give rise to a different notion of solution to the \textsc{rde}, just like in standard rough path theory. 

The space $\mathcal{X}$ obtained by taking the closure in  $\CC^{\alpha-1}\times\CC^{2\alpha-1}$ of the set of all elements of the form $R\xi$ for smooth $\xi$ replaces the space of (geometric) rough paths, and the above theorem is a partial restatement of Lyons' continuity result: namely that the (It\^o) solution map $\Psi$, going from data to solution of the differential equation, is a continuous map from the rough path space $\mathcal{X}$ to $\CC^\alpha$. The space $\mathcal{X}$ is fibered over $\CC^{\alpha-1}$. It allows us to equip the driving distribution with enough information to control the continuity of the solution map to our \textsc{red} problem -- and as we will see below, also the continuity of the solution maps to suitable PDEs. In various contexts the space $\mathcal{X}$ can take different forms, and in general it does not seem to have the rich geometrical and algebraic structure of standard rough paths. 

The verification that suitable approximations $(\xi^n)$ are such that their lifts $(R\xi^n)$  converge in $ \CC^{\alpha-1}\times\CC^{2\alpha-1}$ depends on the particular form of $\xi$. In the case of $\xi$ being a Gaussian stochastic process (like in all our examples above), this verification is the result of almost sure convergence results for elements in a fixed chaos of an underlying Gaussian process, and the proofs rely on elementary arguments on Gaussian random variables.   

Even in the case of \textsc{rde}s, the paracontrolled analysis leads to some interesting results. For example, we have that a more general equation of the form
$$
\partial_t u(t) = F(u(t)) \xi(t) + F'(u(t))F(u(t))\eta(t) ,\qquad
u(0)=u_0,
$$
where $\eta \in C([0,1];\mathbb{R}^n \times \mathbb{R}^n)$,
has a solution map which depends continuously on $(\xi, \vartheta\circ \xi +\eta)\in\CC^{\alpha-1}\times \CC^{2\alpha-1}$. The remarkable fact here is that the solution map depends only on the combination $\vartheta\circ \xi +\eta$ and not on each term separately. Such structural features of the solution map, which can be easily seen using the paracontrolled analysis, are very important in situations where renomalizations are needed, as for example in the \textsc{pam} model. In the \textsc{rde} context we can simply remark that setting $\eta = -\vartheta\circ \xi$, the solution map becomes a continuous function of $\xi \in \CC^{\alpha-1}$, without any further requirement on the bilinear object $\vartheta\circ \xi$. Thus, the equation
$$
\partial_t u(t) = F(u(t)) \xi(t) - F'(u(t))F(u(t))(\vartheta\circ \xi)(t) ,\qquad
u(0)=u_0,
$$
can be readily extended to any $\xi\in\CC^{\alpha-1}$ by continuity. In that sense, this equation should be interpreted as a deterministic version of an It\^o stochastic differential equation where, at the price of a modification of standard rules of calculus, we are able to solve more general problems than in the Stratonovich setting.


\bigskip


We remark that, even if only quite implicitly, paraproducts  have been already exploited in the rough path context in the work of Unterberger on the renormalization of rough paths~\cite{Unterberger2010a, Unterberger2010}, where it is referred to as ``Fourier normal-ordering'', and in the related work of Nualart and Tindel~{\cite{Nualart2011}}.

In this paper we construct weak solutions for the SPDEs under consideration. For an approach using mild solutions see~\cite{Perkowski2013Thesis}. See also~\cite{Gubinelli2013}, where we use the decomposition of continuous functions in a certain Fourier series and similar ideas as developed below, in order to give a new and relatively elementary approach to rough path integration.






\paragraph{Relevant literature.}


Before going into the details, let us describe the context of our study. Consider for example the \textsc{rde} problem above. 
Schwartz' theory of distributions gives a robust framework for defining linear operations on irregular generalized functions. But when trying to handle non-linear operations, we quickly run into problems. For example, in Schwartz' theory, it is not possible to define the product $F(u) \partial_t X(t)$ in the case where $X$ is the sample path of a Brownian motion. The standard analysis of this difficulty goes as follows: $X$ is an $\alpha$--H\"older continuous process for any $\alpha<1/2$, but not better. The solution $u$ has to have the same regularity, which is transferred to $F(u)$ if $F$ is smooth. In this situation, the product $F(u)\partial_t X$ corresponds to the product of an $\alpha$--H\"older continuous function with the distribution $\partial_t X$ which is of order $\alpha-1$. A well known result of analysis (see Section~\ref{sec:bony} below) tells us that a necessary condition for this product to be well defined is that the sum of the orders is positive, that is $2\alpha-1>0$, which is barely violated in the Brownian setting.
This is the classical problem which motivated It\^o's theory of stochastic integrals.

%One of the main objectives of stochastic analysis has indeed been that of studying local properties of stochastic processes.  In  the theory of continuous diffusion processes, the time-derivative $\partial_{t}X_{t}$ of a Brownian motion $X$ plays the role of a universal model of small scale characteristics. It\^o's approach to encode a large class of diffusions $u$ has been  to write them as a local perturbation of a Brownian path $X$ via a \emph{stochastic differential equation} that prescribes the variations of $u$ in terms of those of $X$ via a formal \textsc{rde}:
%\[
%   \partial_{t} u_{t} = F( u_{t} ) \partial_{t} X_{t}.
%\]
%Using It\^{o}'s stochastic integral this equation can ideed be interpreted as a stochastic differential equation and the problematic product can be stipulated to correspond to the time derivative of the stochastic integral:
%\[
%   F(X_t) \partial_t B_t = \partial_t \int_0^t F(X_s) \mathd B_s.
%\]
It\^{o}'s integral has however quite stringent structural requirements: an ``arrow of time'' (i.e. a filtration and adapted integrands), a probability measure (it is defined as $L^2$--limit), and $L^2$--orthogonal increments of the integrator (the integrator needs to be a (semi-) martingale).

If one or several of these assumptions are violated, then Lyons' rough path integral~\cite{Lyons1998, Lyons2002, Lyons2007, Friz2010} can be an effective alternative. For example, it allows to construct pathwise  integrals for, among other processes, fractional Brownian motion, which is not a semimartingale. 



%In one dimension, stochastic analysis usually relies on integrals in order to define products; for a one-dimensional $x$, defining the product $G(u) \partial_x u$ is essentially equivalent to defining the integral $\int G(u) \mathd_x u$. But in the multidimensional case the link between integrals and products is not so compelling.

In the last years, several other works applied rough path techniques to SPDEs. 
%attempted to figure out applications of rough path theory to SPDEs. 
But they all relied on special features of the problem at hand in order to apply the integration theory provided by the rough path machinery. 

A first series of works attempts to deal with ``time"-like irregularities by adapting the standard rough path approach:

\begin{itemize}
\item[--]
Deya, Gubinelli, Lejay, and Tindel~\cite{Gubinelli2006,Gubinelli2012a,Deya2012} deal with SPDEs of the form
\begin{align*}
   L u (t, x) = \sigma ( u ( t, x ) ) \eta ( t,x),
\end{align*}
where $x \in \mathbb{T}$, $L = \partial_t - \Delta$, the noise $\eta$ is a space-time Gaussian distribution (for example white in time and colored in space), and $\sigma$ is some non-linear coefficient. They interpret this as an evolution equation (in time), taking values in a space of functions (with respect to the space variable). They extend the rough path machinery to handle the convolution integrals that appear when applying the heat flow to the noise.

\item[--]
Friz, Caruana, Diehl, and Oberhauser~\cite{Caruana2009,Caruana2011,Friz2011b,Diehl2012} deal with fully non-linear stochastic PDEs with a special structure. Among others, of the form
\begin{align*}
   \partial_t u ( t, x ) = F ( u, \partial_{x} u, \partial_{x}^2 u) + \sigma( t, x) \partial_{x} u( t, x) \eta(t),
\end{align*}
where the spatial index $x$ can be multidimensional, but the noise $\eta$ only depends on time. Such an SPDE can be reinterpreted as a standard PDE with random coefficients via a change of variables involving the flow of the stochastic characteristics associated to $\sigma$. This flow is handled using usual rough path results for \textsc{red}s.

\item[--]
Teichmann~\cite{Teichmann2011} studies semilinear SPDEs of the form
\begin{align*}
   (\partial_t - A) u(t,x) = \sigma(u)(t,x) \eta(t,x),
\end{align*}
where $A$ is a suitable linear operator, in general unbounded, and $\sigma$ is a general non-linear operation on the unknown $u$ which however should satisfy some restrictive conditions. The SPDE is transformed into an SDE with bounded coefficients by applying a  transformation based on the group generated by $A$ on a suitable space.
\end{itemize}

The ``arrow of time'' condition of It\^o's integral is typically violated if the index is a spatial variable and not a temporal variable. Another series of works applied rough path integrals to deal with situations involving irregularities in the ``space" directions:

\begin{itemize}
\item[--]
Bessaih, Gubinelli, and Russo~\cite{Bessaih2005} and Brzezniak, Gubinelli, and Neklyudov~\cite{Brzezniak2010} consider the vortex filament equation which describes the (approximate) motion of a closed vortex line $x(t,\cdot)\in C(\mathbb{T},\mathbb{R}^3)$ in an incompressible three-dimensional fluid:
$$
\frac{\mathd}{\mathd t} x(t,\sigma) = u^{x(t,\cdot)}(x(t,\sigma)),\qquad u^{x(t,\cdot)}(y) = \int_{\mathbb{T}} K(y-x(t,\sigma)) \partial_\sigma x(t,\sigma) \mathd \sigma,
$$
where $K:\mathbb{R}^3 \to \mathcal{L}(\mathbb{R}^3,\mathbb{R}^3)$ is a smooth antisymmetric field of linear transformations of $\mathbb{R}^3$. In the modeling of turbulence it is interesting to study this equation with initial condition $x(0,\cdot)$ sampled according to the law of the three-dimensional Brownian bridge. In this case, the regularity of $x(t,\sigma)$ with respect to $\sigma$ is no better than Brownian for any positive time, and thus the integral in the definition of the velocity field $u^{x(t,\cdot)}$ is not well defined. Rough path theory allows to make sense of this integral and then of the equation.% Here, the irregularities appear along the direction of the (one-dimensional) variable parameterizing the vortex line.

\item[--]
Hairer, Maas, and Weber~\cite{Hairer2011, Hairer2013, Hairer2013b, Hairer2012} build on the insight of Hairer that rough path theory allows to make sense of SPDEs which are ill-defined in standard function spaces due to spatial irregularities. Hairer and Weber~\cite{Hairer2013} extend the \textsc{burgers} type SPDE that we presented above to the case of multiplicative noise. Hairer, Maas, and Weber~\cite{Hairer2012} study approximations to this equation, where they discretize the spatial derivative as $\partial_x u(t,x) \simeq 1/\varepsilon (u(t,x+\varepsilon) - u(t,x))$. They show that in the limit $\varepsilon \rightarrow 0$, this introduces a Stratonovich type correction term to the equation. Finally, Hairer~\cite{Hairer2013b} uses this approach to  define and solve for the first time the Kardar--Parisi--Zhang equation, an SPDE of one spatial index variable that describes the random growth of an interface. This equation was introduced by Kardar, Parisi, and Zhang~\cite{Kardar1986}, and prior to Hairer's work it could only be solved by applying a spatial transform (the Cole-Hopf transform) which had the effect of linearizing the equation.
\end{itemize}

%
%The ``arrow of time'' condition is typically violated if the index is a spatial variable and not a temporal variable. This can be nicely illustrated in two models:
%\begin{enumerate}
%\item The vortex filament equation which describes the (approximate) motion of a closed vortex line $x(t,\cdot)\in C(\mathbb{T},\mathbb{R}^3)$ in an incompressible three-dimensional fluid:
%$$
%\frac{\mathd}{\mathd t} x(t,\sigma) = u^{x(t,\cdot)}(x(t,\sigma)),\qquad u^{x(t,\cdot)}(y) = \int_{\mathbb{T}} K(y-x(t,\sigma)) \partial_\sigma x(t,\sigma) \mathd \sigma,
%$$
%where $K:\mathbb{R}^3 \to \mathcal{L}(\mathbb{R}^3,\mathbb{R}^3)$ is a smooth antisymmetric field of linear transformations of $\mathbb{R}^3$. In the modeling of turbulence it is interesting to study this equation with initial condition $x(0,\cdot)$ sampled according to the law of the three-dimensional Brownian bridge. In this case, the regularity of $x(t,\sigma)$ with respect to $\sigma$ is no better than Brownian for any positive time, and thus the integral in the definition of the velocity field $u^{x(t,\cdot)}$ is not well defined.
%
%\item Hairer's Burgers type SPDE:
%\[
%   L u = G(u) \partial_x u + \xi,
%\]
%on $[0,T] \times \mathbb{T}$, where $\xi$ is a space-time white noise, and $L = \partial_t - \Delta$. This problem is motivated by insights from path sampling, where one can formally derive the equation as an SPDE whose invariant measure describes the law of a certain conditioned diffusion. For every fixed $t > 0$, the solution $\vartheta$ to $L \vartheta = \xi$, $L \vartheta(0) = 0$, is $\alpha$--H\"older continuous in space for every $\alpha < 1/2$. We would expect $u$ to have the same regularity as $v$. But then the product $G(u) \partial_x u$ is ill-defined: we have $G(u) \in \CC^\alpha$ and $\partial_x u \in \CC^{\alpha-1}$, and therefore we are in a situation where $2\alpha -1 <0$.
%
%\end{enumerate}
%
%In both examples, the irregularity is of Brownian kind but it shows up in a ``spatial" variable and it cannot be treated using It\^{o}'s theory. Both problems can however be solved using rough path techniques. The vortex model has been analyzed by Bessaih, Gubinelli, and Russo in~\cite{Bessaih2005}, and the Burgers type equation in a remarkable paper of Hairer~\cite{Hairer2011}.
%It is then natural to ask about extensions of Hairer's approach to higher-dimensional Burgers equations, and in general about applications of rough path ideas in the context of stochastic (or deterministic but very irregular) partial differential equations. 



\paragraph{Alternative approaches.}

In all the papers cited above, the intrinsic one-dimensional nature of rough path theory severely limits possible improvements or applications to other contexts. To the best of our knowledge, the first attempt to remove these limitations is the still unpublished work by Chouk and Gubinelli~\cite{Chouk2013}, extending rough path theory to handle (fractional) Brownian sheets (Gaussian two-parameter stochastic processes akin to (fractional) Brownian motion).

In the recent paper~\cite{Hairer2013a}, Hairer has introduced a theory of \emph{regularity structures} with the aim of giving a more general and versatile notion of regularity. Hairer's theory is also inspired by the theory of controlled rough paths, and it can also be considered a generalization of it to functions of a multidimensional index variable. The crucial insight is that the regularity of the solution to an equation driven by -- say -- Gaussian space-time white noise should not be described in the classical way. Usually we say that a function is smooth if it can be approximated around every point by a polynomial of a given degree (the Taylor polynomial). Since the solution to an SPDE does not look like a polynomial at all, this is not the correct way of describing its regularity. We rather expect that the solution locally looks like the driving noise (more precisely like the noise convoluted with the Green kernel of the linear part of the equation; in the case of RDEs this is the time integral of the white noise, i.e. the Brownian motion). Therefore, in Hairer's theory a function is called smooth if it can locally be well approximated by this convolution (and higher order terms depending on the noise). This notion of smoothness induces a natural topology in which the solutions to semilinear SPDEs depend continuously on the driving signal. Hairer's approach is very general, and allows to handle more complicated problems than the ones we are currently able to treat in the paracontrolled approach. If there is a merit in our approach, then its relative simplicity, the fact that it seems to be very adaptable so that it can be easily modified to treat problems with a different structure, and that we make the connection between harmonic analysis and rough paths.



\paragraph{Plan of the paper.}

Section~\ref{sec:paracontrolled} develops the calculus of paracontrolled distributions. In Section~\ref{sec:ODE} we solve ordinary differential equations driven by suitable Gaussian processes such as the fractional Brownian motion with Hurst index $H>1/3$. In Section~\ref{sec:burgers} we solve a fractional Burgers type equation driven by white noise, and in Section~\ref{sec:PAM} we study a non-linear version of the parabolic Anderson model. %In Appendix~\ref{sec:notation} we fix some conventions and notations, and 
In Appendix~\ref{sec:para} we recall the main concepts of Littlewood-Paley theory and of Bony's paraproduct, and Appendix~\ref{sec:convolution commutators} contains a commutator estimate between paraproduct and time integral. We stress the fact that this paper is mostly self-contained, and in particular we will not need any results from rough path theory and just basic elements of the theory of Besov spaces.



\paragraph{Acknowledgments.} This research has been partly supported by the ANR Project
ECRU (ANR-09-BLAN-0114-01), by a Junior fellowship of the Institut Universitaire de France (IUF), and by a Postdoctoral grant of the Fondation Sciences Math\'ematiques de Paris. The main part of the research was carried out while N.~P. was employed by Humboldt-Universit\"at zu Berlin.
During an Oberwolfach workshop in the summer of 2012, M. Hairer discussed with one of us (M.~G.) his approach to extend rough
path theory and  M.~G. would like
to thank M. Hairer for suggesting the application to the two-dimensional non-linear
parabolic Anderson model discussed in this paper.



\paragraph{Notation and conventions.}%\label{sec:notation}

Throughout the paper, we use the notation $a \lesssim b$ if there exists a constant $c>0$, independent of the variables under consideration, such that $a \leqslant c \cdot b$, and we write $a \simeq b$ if $a \lesssim b$ and $b \lesssim a$. If we want to emphasize the dependence of $c$ on the variable $x$, then we write $a(x) \lesssim_{x} b(x)$. For index variables $i$ and $j$ of Littlewood-Paley decompositions (see below) we write $i \lesssim j$ if there exists $N \in \mathbb{N}$, independent of $j$, such that $i \leqslant j +N$, and we write $i \sim j$ if $i \lesssim j$ and $j \lesssim i$.

An \tmtextit{annulus} is a set of the form $\CA = \{ x \in \mathbb{R}^{d} :a \leqslant | x | \leqslant b \}$ for some $0<a<b$. A {\tmem{ball}} is a set of the form $\CB = \{ x \in \mathbb{R}^{d} : | x | \leqslant b \}$. $\mathbb{T}= \R / (2\pi \Z)$ denotes the torus.

The H\"older-Besov space $B^\alpha_{\infty, \infty}(\R^d, \R^n)$ for $\alpha \in \R$ will be denoted by $\CC^\alpha$, equipped with the norm $\lVert \cdot \rVert_{\alpha}= \lVert \cdot \rVert_{B^\alpha_{\infty,\infty}}$. The local space $\CC^\alpha_\mathrm{loc}$ consists of all $u$ which satisfy $\varphi u \in \CC^\alpha$ for every infinitely differentiable $\varphi$ of compact support. Given two Banach spaces $X,Y$ we denote by $\Lcal ( X,Y )$ the Banach space of linear maps from $X$ to $Y$, endowed with the operator norm $\lVert \cdummy \rVert_{\Lcal(X,Y)}$. More generally, given $k \in \mathbb{N}$ and Banach spaces $X_{1} , \ldots ,X_{k}$, we write $\Lcal^{k} ( X_{1}\times \ldots \times X_{k} ,Y )$ for the space of $k$-linear maps from $X_{1} \times \ldots \times X_{k}$ to $Y$, and $\lVert \cdummy \rVert_{\Lcal^{k} (X_{1} \times \ldots \times X_{k} ,Y )}$ for the operator norm. We denote by $C( X,Y )$ the Banach space of continuous maps from $X$ to $Y$, endowed with the supremum norm $\lVert \cdummy \rVert_{C ( X,Y )}$. We write $C_{T} Y=C ( [ 0,T ] ,Y )$ for the space of continuous maps from $[ 0,T ]$ to $Y$, equipped with the supremum norm $\lVert \cdummy \rVert_{C_{T} Y}$. If $\alpha \in ( 0,1 )$, then we also define $C^{\alpha}_{T} Y$ as the space of $\alpha$-H{\"o}lder continuous functions from $[ 0,T ]$ to $Y$, endowed with the seminorm
\[
   \| f \|_{C^{\alpha}_{T} Y} = \sup_{0 \leqslant s<t \leqslant T} \frac{\| f( t ) -f ( s ) \|_{Y}}{| t-s |^{\alpha}}.
\]
If $f$ is a map from $A \subset \R$ to the linear space $Y$, then we write $f_{s,t} = f(t) - f(s)$, so that $\| f \|_{C^\alpha_T Y} = \sup_{0 \leqslant s<t \leqslant T} \| f_{s,t} \|_{Y}/ | t-s |^{\alpha}$. For  $f \in L^p(\mathbb{T})$ we write $\Vert  f(x) \Vert ^p_{L^p_x (\mathbb{T} )} := \int_{\mathbb{T}} | f ( x )|^p \mathd x$.

For a multi-index $\mu = ( \mu_{1} , \ldots , \mu_{d} ) \in \mathbb{N}^{d}$ we write $| \mu | = \mu_{1} + \ldots + \mu_{d}$ and $\partial^{\mu} = \partial^{| \mu |} / \partial_{x_{1}}^{\mu_{1}} \cdots \partial_{x_{d}}^{\mu_{d}}$. %{\color{red} For $x= ( x_{1} , \ldots ,x_{d} ) \in \mathbb{R}^{d}$ we write $x^{\mu} =x_{1}^{\mu_{1}} \cdots x_{d}^{\mu_{d}}$.}
$\mathD F$ or $F'$ denote the total derivative of $F$. For $k \in \mathbb{N}$ we denote by $\mathD^{k} F$ the $k$-th order derivative of $F$, and by $C^{k}_{b} =C^{k}_{b} (\mathbb{R}^{d} ,\mathbb{R}^{n} )$ the space of $k$ times continuously differentiable functions whose partial derivatives up to order $k$ are bounded, equipped with is usual norm $\lVert \cdot \rVert_{C^k_b}$. %$\|F\|_{C^{k}_{b}} = \sum_{|\mu | \leqslant k} \| \partial^{\mu} F\|_{L^{\infty}} \simeq \sum_{j \leqslant k} \| \mathD^{j} F \|_{L^{\infty}}$.
We also write $\partial_{x}$ for the partial derivative in direction $x$, and if $F\colon \mathbb{R} \times\mathbb{R}^{d} \rightarrow \mathbb{R}^{n}$, then we write $\mathD_{x} F(t,x)$ for its spatial derivative in the point $( t,x ) \in \mathbb{R} \times \mathbb{R}^{d}$.

The space of real valued infinitely differentiable functions of compact support is denoted by $\CD ( \mathbb{R}^{d} )$ or $\CD$. The space of Schwartz functions is denoted by $\CS ( \mathbb{R}^{d} )$ or $\CS$. Its dual, the space of tempered distributions, is $\CS' ( \mathbb{R}^{d} )$ or $\CS'$. If $u$ is a vector of $n$ tempered distributions on $\mathbb{R}^{d}$, then we write $u \in \CS' (\mathbb{R}^{d} ,\mathbb{R}^{n} )$. The Fourier transform is defined with the normalization
\[
   \CF u ( z ) = \hat{u} ( z ) = \int_{\mathbb{R}^{d}} e^{- \iota \langle z,x \rangle} u ( x ) \mathd x,
\]
so that the inverse Fourier transform is given by $\CF^{-1} u ( z ) = ( 2 \pi)^{-d} \CF u ( -z )$. If $\varphi$ is a smooth function, such that $\varphi$ and all its partial derivatives are at most of polynomial growth at infinity, then we define the Fourier multiplier $\varphi ( \mathD )$ by $\varphi (\mathD ) u= \CF^{-1} ( \varphi \CF u)$ for any $u \in \CS'$. More generally, we define $\varphi ( \mathD ) u$ by this formula whenever the right hand side makes sense. The scaling operator $\Lambda$ on $\CS'$ is defined for $\lambda >0$ by $\Lambda_\lambda u = u(\lambda\cdot)$.

Throughout the paper, $(\chi, \rho)$ will denote a dyadic partition of unity, and $(\Delta_j)_{j \ge -1}$ will denote the Littlewood-Paley blocks associated to this partition of unity, i.e. $\Delta_{-1}= \chi(\mathD)$ and $\Delta_j = \rho(2^{-j} \mathD)$ for $j \ge 0$. We will often write $\rho_j$, by which we mean $\chi$ if $j=-1$, and we mean $\rho(2^{-j} \cdot)$ if $j \ge 0$. We also use the notation $S_j = \sum_{i < j} \Delta_i$.


\section{Paracontrolled calculus}\label{sec:paracontrolled}


\subsection{Bony's paraproduct}\label{sec:bony}

Paraproducts are bilinear operations introduced by Bony~\cite{Bony1981} in order to linearize a class of non-linear PDE problems. In this section we will introduce paraproducts to the extent of our needs. We will be using the Littlewood-Paley theory of Besov spaces. The reader can peruse Appendix~\ref{sec:para}, where we summarize the basic elements of Besov space theory and Littlewood-Paley decompositions which will be needed in the remainder of the paper.

 One of the simplest situations where paraproducts appear naturally is in the analysis of the product of two Besov distributions. 
In general, the product $f g$ of two distributions $f \in \CC^{\alpha}$ and $g \in \CC^{\beta}$ is not well defined unless $\alpha + \beta >0$. In terms of Littlewood--Paley blocks, the product $f g$ can be (at least formally) decomposed as
\[
   fg= \sum_{j \geqslant -1} \sum_{i \geqslant -1} \Delta_{i} f \Delta_{j} g=f \prec g+f \succ g+f \circ g.
\]
Here $f \prec g$ is the part of the double sum with $i<j-1$, and $f \succ g$ is the part with $i>j+1$, and $f \circ g$ is the ``diagonal'' part, where $|i-j | \leqslant 1$. More precisely, we define
\[
   f \prec g= g\succ f = \sum_{j \geqslant -1} \sum_{i=-1}^{j-2} \Delta_{i} f \Delta_{j} g \qquad \text{and}\qquad
   f \circ g= \sum_{| i-j | \leqslant 1} \Delta_{i} f \Delta_{j} g.
\]
We also introduce the notation
\[
   f \succcurlyeq g=f \succ g+f \circ g.
\]
This decomposition  behaves nicely with respect to Littlewood--Paley theory. Of course, the decomposition depends on the dyadic partition of unity used to define the blocks $\Delta_{j}$, and also on the particular choice of the pairs $( i,j )$ in the diagonal part. Our choice of taking all $(i,j)$ with $| i-j| \leqslant 1$ into the diagonal part corresponds to property \ref{def:dyadic partition3}. in the definition of dyadic partition of unity in Appendix~\ref{sec:para},  where we assumed that $\tmop{supp} ( \rho ( 2^{-i} \cdummy ) ) \cap \tmop{supp} ( \rho ( 2^{-j} \cdummy ) ) = \emptyset$ for $| i-j | >1$. This means that every term in the series
\[
   f \prec g = \sum_{j \geqslant -1} \sum_{i=-1}^{j-2} \Delta_{i} f \Delta_{j} g= \sum_{j \geqslant -1} S_{j-1} f \Delta_{j} g
\]
has a Fourier transform which is supported in a suitable annulus, and of course the same holds true for $f \succ g$. On the other side, the terms in the diagonal part $f \circ g$ have a Fourier transform which is supported in a ball. We call $f \prec g$ and $f \succ g$ \tmtextit{paraproducts}, and $f \circ g$ the \tmtextit{resonant} term.

Bony's crucial observation is that $f \prec g$ (and thus $f \succ g$) is always a well-defined distribution. In particular, if $\alpha >0$ and $\beta \in \mathbb{R}$, then $( f,g ) \mapsto f \prec g$ is a bounded bilinear operator from $\CC^{\alpha} \times \CC^{\beta}$ to $\CC^{\beta}$. Heuristically, $f \prec g$ behaves at large frequencies like $g$ (and thus retains the same regularity), and $f$ provides only a modulation of $g$ at larger scales. The only difficulty in defining $fg$ for arbitrary distributions lies in handling the diagonal term $f \circ g$. The basic result about these bilinear operations is given by the following estimates.

\begin{lemma}[Paraproduct estimates, \cite{Bony1981}]\label{thm:paraproduct}
   For any $\beta \in \mathbb{R}$ we have
   \begin{equation} \label{eq:para-1}
       \| f \prec g \|_{\beta} \lesssim_{\beta} \| f \|_{L^{\infty}} \| g \|_{\beta} ,
   \end{equation}
   and for $\alpha <0$ furthermore
   \begin{equation}\label{eq:para-2}
      \| f \prec g \|_{\alpha + \beta} \lesssim_{\alpha , \beta} \| f \|_{\alpha} \| g \|_{\beta}.
   \end{equation}
   For $\alpha + \beta >0$ we have
   \begin{equation}\label{eq:para-3}
      \| f \circ g \|_{\alpha + \beta} \lesssim_{\alpha , \beta} \| f \|_{\alpha} \| g \|_{\beta}.
    \end{equation}
\end{lemma}

\begin{proof}
  Observe that there exists an annulus $\CA$ such that $S_{j-1} f \Delta_{j} g$ has Fourier transform supported in $2^{j} \CA$, and that for $f \in L^{\infty}$ we have
  \[
     \| S_{j-1} f \Delta_{j} g \|_{L^{\infty}} \leqslant \| S_{j-1} f  \|_{L^{\infty}} \| \Delta_{j} g \|_{L^{\infty}} \leqslant \| f \|_{L^{\infty}} 2^{-j \beta} \| g \|_{\beta}.
  \]
  On the other side, if $\alpha <0$ and $f \in \CC^{\alpha}$, then
  \[
     \| S_{j-1} f \Delta_{j} g \|_{L^{\infty}} \leqslant \sum_{i \leqslant j-2} \| \Delta_{i} f \|_{L^{\infty}} \| \Delta_{j} g \|_{L^{\infty}} \lesssim \| f \|_{\alpha} \| g \|_{\beta} \sum_{i \leqslant j-2} 2^{-i\alpha -j \beta} \lesssim \| f \|_{\alpha} \| g \|_{\beta} 2^{-j ( \alpha + \beta )} .
  \]
  By Lemma~\ref{lem: Besov regularity of series}, we thus obtain~{\eqref{eq:para-1}} and~{\eqref{eq:para-2}}. To estimate $f \circ g$, observe that the term $u_{j} = \Delta_{j} f \sum_{i: | i-j | \leqslant 1} \Delta_{i} g$ has Fourier transform supported in a ball $2^{j} \CB$, and that
  \[
     \| u_{j} \|_{L^{\infty}} \lesssim \| \Delta_{j} f \|_{L^{\infty}} \sum_{i: | i-j | \leqslant 1} \| \Delta_{i} g \|_{L^{\infty}} \lesssim \| f \|_{\alpha} \| g \|_{\beta} 2^{- ( \alpha + \beta ) j}.
  \]
  So if $\alpha + \beta >0$, then we can apply the second part of Lemma~\ref{lem: Besov regularity of series} to obtain that $f \circ g= \sum_{j \geqslant -1} u_{j}$ is an element of $\CC^{\alpha + \beta}$ and that equation~{\eqref{eq:para-3}} holds.
\end{proof}

A natural corollary is that the product $fg$ of two elements $f \in \CC^{\alpha}$ and $g \in \CC^{\beta}$ is well defined as soon as $\alpha + \beta >0$, and that it belongs to $\CC^{\gamma}$, where $\gamma = \min \{\alpha , \beta , \alpha + \beta \}$.


\subsection{Paracontrolled distributions and RDEs}\label{sec:paracontrolled ode}

Consider the \textsc{rde}
\begin{equation}\label{eq:rde-intro} 
  \partial_{t} u=F ( u ) \xi, \qquad u(0) = u_0,
\end{equation}
where $u_0 \in \R^d$, $u\colon \mathbb{R} \rightarrow \mathbb{R}^{d}$ is a continuous vector
valued function, $\partial_{t}$ is the time derivative, $\xi\colon \mathbb{R} \rightarrow \mathbb{R}^{n}$ is a vector valued distribution with values in
$\CC^{\alpha -1}$ for some $\alpha \in ( 1/3,1 )$, and
$F\colon \mathbb{R}^{d} \rightarrow \Lcal ( \mathbb{R}^{n} ,\mathbb{R}^{d})$ is a family of vector fields on $\mathbb{R}^{d}$. A natural approach is to
understand this equation as limit of the classical ODEs
\begin{equation}\label{eq:rde-reg-intro}
  \partial_{t} u^{\varepsilon} = F ( u^{\varepsilon} ) \xi^{\varepsilon}, \qquad u^\varepsilon(0) = u_0,
\end{equation}
for a family of smooth approximations $(\xi^{\varepsilon})$ of $\xi$ such that
$\xi^{\varepsilon} \rightarrow \xi$ in $\CC^{\alpha -1}$ as $\varepsilon \rightarrow 0$. In order to pass to the limit, we are looking for a
priori estimates on $u^{\varepsilon}$ which require only a control on the
$\CC^{\alpha -1}$ norm of $\xi$.

To avoid cumbersome notation, we will work at the level of
equation~{\eqref{eq:rde-intro}} for smooth $\xi$, where it should be understood that
our aim is to obtain a priori estimates for the solution, in order to safely pass
to the limit and extend the solution map to a larger class of data. The
natural regularity of $u$ is $\CC^{\alpha}$, since $u$ should gain one
derivative with respect to $F ( u ) \xi$, which will not behave better than
$\xi$, and will therefore be in $\CC^{\alpha -1}$.

We use the paraproduct decomposition to write the right hand side of {\eqref{eq:rde-intro}} as a
sum of the three terms
%\begin{equation}
%  F ( u ) \prec \xi + F ( u ) \circ  \xi + F ( u ) \succ \xi,
%\end{equation}
%where $F(u) \prec \xi$ is the part of the product where $\xi$ has higher frequencies, $F(u) \succ \xi$ is the part where $F(u)$ has higher frequencies, and $F(u) \circ \xi$ is the part where the frequencies of $F(u)$ and $\xi$ are comparable. The point of this decomposition is that the components $F(u) \prec \xi$ and $F(u) \succ \xi$ are always well defined tempered distributions irrespectively of the regularities of $F(u)$ and of $\xi$. The resonant term $F(u) \circ \xi$ is well defined  \emph{only if} the regularities of $F(u)$ and $\xi$ add up to a strictly positive number. The expected regularities of the single components are as follows:
\begin{equation}\label{eq:para-rhs-intro}
  \underbrace{F ( u ) \prec \xi}_{\alpha -1} + \underbrace{F ( u ) \circ  \xi}_{2 \alpha -1} + \underbrace{F ( u ) \succ \xi}_{2 \alpha -1}
\end{equation}
(where the quantity indicated by the underbrace corresponds to the expected regularity of each term). Note however that unless $2 \alpha -1>0$, the resonant term
$F ( u ) \circ   \xi$ cannot be controlled using only the $\CC^{\alpha}$--norm
of $u$ and the $\CC^{\alpha -1}$--norm of $\xi$. If $F$ is at least in $C^2$, we can use a paralinearization result (see Lemma~\ref{lemma:para-taylor} below) to rewrite this term as
\begin{equation}\label{eq:rde expand Fu}
   F ( u ) \circ \xi =F' ( u )  ( u \circ \xi ) + \Pi_{F} ( u, \xi ),
\end{equation}
where the remainder $\Pi_{F} ( u, \xi )$ is well defined under the condition
$3 \alpha -1>0$, provided that $u \in \CC^{\alpha}$ and $\xi \in \CC^{\alpha -1}$. In this case it belongs to $\CC^{3 \alpha -1}$. The
difficulty is now localized in the linearized resonant product $u \circ \xi$. In order to control this term, we would like to exploit the fact that
the function $u$ is not a generic element of $\CC^{\alpha}$ but that it has a
specific structure, since its derivative $\partial_{t} u$ has to match the
paraproduct decomposition given in~\eqref{eq:para-rhs-intro}. Thus, we postulate
that the solution $u$ is given by the following \tmtextit{paracontrolled} ansatz:
\[
   u=u^{\vartheta} \prec \vartheta +u^{\sharp},
\]
where $u^{\vartheta} , \vartheta \in \CC^{\alpha}$ and the remainder
$u^{\sharp}$ is in $\CC^{2 \alpha}$. This decomposition allows for a finer analysis
of the resonant term $u \circ \xi$. Indeed, we have
\begin{equation}\label{eq:rde expand resonant term}
   u \circ   \xi = ( u^{\vartheta} \prec \vartheta ) \circ   \xi +u^{\sharp} \circ   \xi =u^{\vartheta} ( \vartheta \circ \xi ) +C ( u^{\vartheta} , \vartheta , \xi ) +u^{\sharp} \circ   \xi ,
\end{equation}
where the commutator is defined by $C ( u^{\vartheta} ,\vartheta , \xi ) = ( u^{\vartheta} \prec \vartheta ) \circ \xi -u^{\vartheta}( \vartheta \circ \xi )$. %Below we will show that under the condition $3\alpha-1 >0$, the trilinear operator $C$ is bounded from $\CC^\alpha \times \CC^\alpha \times \CC^{\alpha - 1}$ to $\CC^{3\alpha - 1}$.
Observe now that the term $u^{\sharp} \circ   \xi$
does not pose any further problem, as it is bounded in $\CC^{3 \alpha -1}$.
Moreover, we will show that the commutator is a bounded multilinear function of its arguments as
long as the sum of their regularities is strictly positive, see
Lemma~\ref{lemma:commutator} below. By assumption, we have $3 \alpha -1>0$,
and therefore $C ( u^{\vartheta} , \vartheta , \xi ) \in \CC^{3 \alpha -1}$.
The only problematic term which remains to be handled is thus $\vartheta \circ   \xi$. Here we need to make the assumption that $\vartheta \circ \xi \in \CC^{2 \alpha -1}$ in order for the product $u^{\vartheta} (\vartheta \circ   \xi )$ to be well defined. That assumption is not guaranteed
by the analytical estimates at hand, and it has to be added as a
further requirement. Granting this, we have obtained that the right hand side
of equation~{\eqref{eq:rde-intro}} is well defined and a continuous function of $(u,u^{\vartheta} ,u^{\sharp} , \vartheta , \xi , \vartheta \circ \xi )$.

The paracontrolled ansatz and the Leibniz rule for the paraproduct now imply
that~{\eqref{eq:rde-intro}} can be rewritten as
\[
   \partial_{t} u= \partial_{t} ( u^{\vartheta} \prec \vartheta +u^{\sharp} ) = \partial_{t} u^{\vartheta} \prec \vartheta +u^{\vartheta} \prec \partial_{t} \vartheta + \partial_{t} u^{\sharp} =F ( u ) \prec \xi +F ( u ) \circ   \xi +F ( u ) \succ \xi.
\]
If we choose $\vartheta$ such that $\partial_{t} \vartheta = \xi$ %and $\vartheta ( 0 ) =0$
and we set $u^{\vartheta} =F ( u )$, then we can
use~{\eqref{eq:rde expand Fu}} and~{\eqref{eq:rde expand resonant term}} to
obtain the following equation for the remainder $u^{\sharp}$:
\begin{align*}%\label{eq:Phi definition} 
  \partial_{t} u^{\sharp} & = F' ( u ) F ( u ) (\vartheta \circ \xi ) +F ( u ) \succ \xi - ( \partial_{t} F ( u ) \prec \vartheta ) \\
  &\qquad +F' ( u ) C ( F ( u ) , \vartheta , \xi ) +F' ( u )  ( u^{\sharp} \circ \xi) + \Pi_{F} ( u, \xi ).
\end{align*}
Together with the equation $u=F ( u ) \prec \vartheta +u^{\sharp}$, this completely describes the solution and allows us to obtain an a priori estimate on $u$ in terms of $(u_0, \|\xi\|_{\alpha-1}, \|\vartheta \circ \xi\|_{2\alpha-1})$. With this estimate at hand, it is now easy to show that if $F \in C^3_b$, then $u$ depends continuously on the data $(u_0, \xi, \vartheta \circ \xi)$, so that we can pass to the limit in~\eqref{eq:rde-reg-intro} and make sense of the solution to~\eqref{eq:rde-intro} also for irregular $\xi \in \CC^{\alpha-1}$ as long as $\alpha > 1/3$.


\subsection{Commutator estimates and paralinearization}

In this section we prove some lemmas which will allow us to perform algebraic computations with the paraproduct and the resonant term, thus justifying the analysis of the previous section.

\begin{lemma}[see also Lemma 2.97 of {\cite{Bahouri2011}}]\label{l:bahouri commutator}
  Let $f \in \CC^{\alpha}$ for $\alpha \in ( 0,1 )$, and let $g \in L^{\infty}$. For any $j \geqslant -1$ we have
  \[
     \| [ \Delta_{j} ,f ] g \|_{L^{\infty}} = \| \Delta_{j} ( f g ) -f \Delta_{j} g \|_{L^{\infty}} \lesssim 2^{- \alpha j} \|  f \|_{\alpha} \| g \|_{L^{\infty}} .
  \]
\end{lemma}

This commutator lemma is easily proven by writing $\Delta_{j} = \rho_{j} (\mathD )$ as a convolution operator, and using the embedding of $\CC^{\alpha}$ in the space of H{\"o}lder continuous functions.

\begin{lemma}\label{lemma:paraprojection}
  Assume that $\alpha \in ( 0,1 )$ and $\beta \in \mathbb{R}$, and let $f \in \CC^{\alpha}$ and $g \in \CC^{\beta}$. Then
  \[
     \Delta_{j} ( f \prec g ) =f \Delta_{j} g+R_{j} ( f,g ) ,
  \]
  for all $j \geqslant -1$, with a remainder $R_{j} ( f,g )$ which satisfies $\| R_{j} ( f,g ) \|_{L^{\infty}} \lesssim 2^{-j ( \alpha + \beta )} \| f \|_{\alpha} \| g  \|_{\beta}$.
\end{lemma}

\begin{proof}
  Note that $f \prec g= \sum_{i} f \prec \Delta_{i} g$, and that the Fourier  transform of $f \prec \Delta_{i} g$ is supported in an annulus of the form $2^{i} \CA$. Hence, we have $\Delta_{j} ( f \prec \Delta_{i} g ) \neq 0$ only if $j \sim i$, which leads to
  \begin{align*}
     \Delta_{j} ( f \prec g ) &= \sum_{i:i \sim j} \Delta_{j} ( f \prec \Delta_{i} g ) = \sum_{i:i \sim j} \Delta_{j} ( f \Delta_{i} g ) -  \sum_{i:i \sim j} \Delta_{j} ( f \succcurlyeq \Delta_{i} g ) \\
     & = \sum_{i:i \sim j} f \Delta_{j} \Delta_{i} g- \sum_{i:i \sim j} [\Delta_{j} ,f ] \Delta_{i} g- \sum_{i:i \sim j} \Delta_{j} ( f \succcurlyeq \Delta_{i} g ),
  \end{align*}
  where we recall that $[ \Delta_{j} ,f ] \Delta_{i} g= \Delta_{j} ( f\Delta_{i} g ) -f \Delta_{j} \Delta_{i} g$ denotes the commutator. The sum over $i$ with $i \sim j$ can be chosen to encompass enough terms so that $\Delta_{j} g= \sum_{i:i \sim j} \Delta_{j} \Delta_{i} g$, and therefore we conclude that
  \[
     \| \Delta_{j} ( f \prec g ) -f \Delta_{j} g \|_{L^{\infty}} \leqslant  \sum_{i:i \sim j} \| [ \Delta_{j} ,f ] \Delta_{i} g \|_{L^{\infty}} -  \sum_{i:i \sim j} \| \Delta_{j} ( f \succcurlyeq \Delta_{i} g )
     \|_{L^{\infty}}.
  \]
  We apply Lemma~\ref{l:bahouri commutator} to each term of the first sum, and the paraproduct estimates to each term of the second sum, to obtain
  \[
     \| \Delta_{j} ( f \prec g ) -f \Delta_{j} g \|_{L^{\infty}} \lesssim  2^{-j ( \alpha + \beta )} \| f \|_{\alpha} \| g \|_{\beta} .
  \]
\end{proof}

Using this result, it is easy to prove our basic commutator lemma.

\begin{lemma}\label{lemma:commutator}
   Assume that $\alpha \in ( 0,1 )$ and $\beta , \gamma \in \mathbb{R}$ are such that $\alpha + \beta + \gamma >0$ and $\beta + \gamma <0$. Then for smooth $f,g,h$, the trilinear operator
  \[
     C ( f,g,h ) = ( ( f \prec g ) \circ h ) -f ( g \circ h )
  \]
  allows for the bound
  \[
     \| C ( f,g,h ) \|_{\alpha + \beta + \gamma} \lesssim \| f \|_{\alpha} \|  g \|_{\beta} \| h \|_{\gamma} .
  \]
  Thus, $C$ can be uniquely extended to a bounded trilinear operator in $\Lcal^{3} \left( \CC^{\alpha}\! \times\! \CC^{\beta}\! \times\! \CC^{\alpha}, \CC^{\alpha + \beta + \gamma} \right)$.
\end{lemma}

\begin{proof}
  Let $f,g,h$ be smooth functions and write
  \[
     C ( f,g,h ) = ( ( f \prec g ) \circ h ) -f ( g \circ h ) = \sum_{j,k  \geqslant -1} \sum_{i: | i-j | \leqslant 1} [ \Delta_{i} ( \Delta_{k} f  \prec g ) \Delta_{j} h- \Delta_{k} f \Delta_{i} g \Delta_{j} h ].
  \]
  Observe that for fixed $k$, the term $\Delta_{k} f \prec g$ has a Fourier transform supported outside of a ball $2^{k} \CB$. Thus, we have $\Delta_{i}( \Delta_{k} f \prec g ) = \mathbf{1}_{i \gtrsim k} \Delta_{i} ( \Delta_{k}f \prec g )$. We can therefore apply Lemma~\ref{lemma:paraprojection} to obtain
  \begin{align}\label{eq:commutator pr1}\nonumber
     C ( f,g,h ) &= \sum_{j,k \geqslant -1} \sum_{i: | i-j | \leqslant 1} [\mathbf{1}_{i \gtrsim k} ( \Delta_{k} f \Delta_{i} g+R_{i} ( \Delta_{k} f,g ) ) \Delta_{j} h- \Delta_{k} f \Delta_{i} g \Delta_{j} h] \\
    &= \sum_{j,k \geqslant -1} \sum_{i: | i-j | \leqslant 1} [\mathbf{1}_{i \gtrsim k} R_{i} ( \Delta_{k} f,g ) \Delta_{j} h -  \mathbf{1}_{i \leqslant k-N} \Delta_{k} f \Delta_{i} g \Delta_{j} h]
  \end{align}
  for some fixed $N \in \mathbb{N}$. We treat the two sums separately. First
  observe that for fixed $k$, the term $\sum_{i: | i-j | \leqslant 1} \mathbf{1}_{i \leqslant k-N} \Delta_{k} f \Delta_{i} g \Delta_{j} h$ has a
  Fourier transform which is supported in a ball of the form $2^{k} \CB$.
  Moreover,
  \begin{align*}
     \Bigg\| \sum_{i: | i-j | \leqslant 1} \mathbf{1}_{i \leqslant k-N} \Delta_{k} f \Delta_{i} g \Delta_{j} h \Bigg\|_{L^{\infty}} & \lesssim 2^{-k \alpha} \| f \|_{\alpha} \sum_{i=-1}^{k-N} 2^{-i ( \beta + \gamma)} \| g \|_{\beta} \| h \|_{\gamma} \\
     & \simeq 2^{-k ( \alpha + \beta + \gamma )} \| f \|_{\alpha} \| g \|_{\beta} \| h \|_{\gamma},
  \end{align*}
  where in the second step we used that $\beta + \gamma <0$. Since $\alpha +  \beta + \gamma >0$, the estimate for the second series
  in~{\eqref{eq:commutator pr1}} therefore follows from Lemma~\ref{lem: Besov regularity of series}.
  
  For the first series, recall that $R_{i} ( \Delta_{k} f,g ) = \Delta_{i} ( \Delta_{k} f \prec g ) - \Delta_{k} f \Delta_{i} g$. So for fixed $j$, the
  Fourier transform of $\sum_{k \geqslant -1} \sum_{i: | i-j | \leqslant 1} \mathbf{1}_{i \gtrsim k} R_{i} ( \Delta_{k} f,g ) \Delta_{j} h$ is
  supported in ball $2^{j} \CB$. Furthermore,
  Lemma~\ref{lemma:paraprojection} yields
  \begin{align*}
     &\Bigg\| \sum_{k \geqslant -1} \sum_{i: | i-j | \leqslant 1} \mathbf{1}_{i \gtrsim k} R_{i} ( \Delta_{k} f,g ) \Delta_{j} h \Bigg\|_{L^{\infty}} = \Bigg\| \sum_{i: | i-j | \leqslant 1} R_{i}  \Bigg( \sum_{k \lesssim i} \Delta_{k} f,g \Bigg) \Delta_{j} h  \Bigg\|_{L^{\infty}} \\
     &\hspace{50pt} \lesssim \sum_{i: | i-j | \leqslant 1} 2^{-i ( \alpha + \beta )} \bigg\| \sum_{k \lesssim i} \Delta_{k} f \bigg\|_{\alpha} \| g \|_{\beta} 2^{-j \gamma} \| h \|_{\gamma} \lesssim 2^{-j ( \alpha + \beta + \gamma )} \| f \|_{\alpha} \| g \|_{\beta} \| h \|_{\gamma},
  \end{align*}
  so that the claimed bound for $\| C ( f,g,h ) \|_{\alpha + \beta + \gamma}$
  follows from another application of Lemma~\ref{lem: Besov regularity of series}.
  
  Now we can uniquely extend $C$ to a bounded trilinear operator on the
  closure of the smooth functions in $\CC^{\alpha} \times \CC^{\beta} \times \CC^{\gamma}$. Unfortunately, the closure is a strict subset of
  $\CC^{\alpha} \times \CC^{\beta} \times \CC^{\gamma}$. But this can be
  easily resolved: We obtain similar bounds for $C$ acting on $\CC^{\alpha'} \times \CC^{\beta'} \times \CC^{\gamma'}$ for $\alpha' \in ( 0,1 )$ and
  $\beta' , \gamma' \in \mathbb{R}$, such that $\alpha' < \alpha$, $\beta' < \beta$, $\gamma' < \gamma$, and $\alpha' + \beta'+ \gamma' >0$. Since
  $\CC^{\alpha} \times \CC^{\beta} \times \CC^{\gamma}$ is contained in the
  closure of the smooth functions in $\CC^{\alpha'} \times \CC^{\beta'} \times \CC^{\gamma'}$, the extension of $C$ to $\CC^{\alpha} \times \CC^{\beta} \times \CC^{\gamma}$ is unique.
\end{proof}

\begin{remark}
  The restriction $\beta + \gamma <0$ is not problematic. If $\beta + \gamma >0$, then $( f \prec g ) \circ h$ can be treated with the usual paraproduct estimates, without the need of introducing a commutator. If $\beta + \gamma =0$, then we can apply the commutator estimate with $\gamma' < \gamma$
  sufficiently close to $\gamma$ such that $\alpha + \beta + \gamma' >0$.
\end{remark}

Let us define the local $C^{k}$ norm of a $k$ times continuously
differentiable function $F$ by setting
\[
   \mathcal{M}_{F,k} ( u ) = \sum_{\ell \leqslant k} \sup_{| x | \leqslant \| u \|_{L^{\infty}}} | \mathD^{\ell} F ( x ) |
\]
for every $u \in L^{\infty}$. Our next result is a simple paralinearization lemma for non-linear operators.

\begin{lemma}[see also~\cite{Bahouri2011}, Theorem 2.92]\label{lemma:paralinearization}
  Let $\alpha \in ( 0,1 )$ and let $F \in C^{2}$. There exists a locally bounded map $R_{F} \colon  \CC^{\alpha} \rightarrow \CC^{2 \alpha}$ such that
  \begin{equation}\label{eq:para-linearization}
    F ( f ) =F' ( f ) \prec f+R_{F} ( f )
  \end{equation}
  for all $f \in \CC^{\alpha}$. More precisely, we have
  \[
     \| R_{F} ( f ) \|_{2 \alpha} \lesssim \mathcal{M}_{F,0} ( f ) +\mathcal{M}_{F' ,1} ( f ) \| f \|_{\alpha}^{2}.
  \]
  If $F \in C^{3}$, then $R_{F}$ is locally Lipschitz continuous:
  \[
     \| R_{F} ( f ) -R_{F} ( g ) \|_{2 \alpha} \lesssim ( \mathcal{M}_{F' ,2}( f ) +\mathcal{M}_{F' ,2} ( g ) ) ( \| f-g \|_{\alpha} + \| f-g \|_{\alpha}^{2} ).
  \]
\end{lemma}

\begin{proof}
  Since $f$ is a function, we can write $\Delta_{i} F ( f )$ as a convolution,
  and obtain for $i \geqslant 1$ that
  \begin{align*}
     \Delta_{i} F ( f ) ( x ) & = \int K_{i} ( x-y ) K_{<i-1} ( x-z ) F ( f ( y) ) \mathd y \mathd z \\
     & = \int K_{i} ( x-y ) K_{<i-1} ( x-z ) ( F ( f ( y ) ) -F ( f ( z ) ) ) \mathd y \mathd z,
  \end{align*}
  where $K_{i} = \CF^{-1} \rho_{i}$, $K_{<i-1} = \sum_{j<i-1} K_{i}$,
  and where we used that $\int K_{i} ( x ) \mathd x= \rho_{i} ( 0 ) =0$ for $i \geqslant 0$ and $\int K_{<i-1} ( x ) \mathd x=1$ for $i \geqslant 1$. Now a
  first order Taylor expansion applied to $F$ yields
  \begin{equation}\label{eq:paralin-step}
    \Delta_{i} F ( f ) =S_{i-1} F' ( f ) \Delta_{i} f+A_{i,F} ( f ),
  \end{equation}
  where the map $A_{i,F}$ is defined by
  \[
     A_{i,F} ( \varphi ) ( x )= \int_{0}^{1} \!\!\mathd r\!\! \int K_{i} ( x-y ) K_{<i-1} ( x-z ) ( F' ( ( 1-r  ) \varphi ( z ) +r \varphi ( y ) ) -F' ( \varphi ( z ) ) ) ( \varphi ( y ) - \varphi ( z ) ) \mathd y \mathd z. 
  \]
  From~{\eqref{eq:paralin-step}} we see that $A_{i,F} ( \varphi )$ has a Fourier
  transform supported in an annulus $2^{i} \CA$. Since $\alpha \in ( 0,1 )$, we can identify $\CC^{\alpha}$ with the classical
  H{\"o}lder space, and this implies that
  \[
     \| A_{i,F} ( f ) \|_{L^{\infty}} \lesssim \mathcal{M}_{F' ,1} ( f ) \| f  \|_{\alpha}^{2} 2^{-2i \alpha}.
  \]
  So if we let $R_{F} ( f ) =S_{1} F ( f ) + \sum_{i \geqslant 1} A_{i,F} ( f )$, then we have $F ( f ) =F' ( f ) \prec f+R_{F} ( f )$, and for $j$ large
  enough so that $S_{1} \Delta_{j} =0$ we get
  \[
     \| \Delta_{j} R_{F} ( f ) \|_{L^{\infty}} \lesssim \sum_{i:i \sim j} \| A_{i,F} ( f ) \|_{L^{\infty}} \lesssim 2^{-2 \alpha j} \mathcal{M}_{F' ,1} ( f ) \| f \|_{\alpha}^{2}.
  \]
  The estimate for small $j$ is easy, and this proves that $R_{F}$ is bounded.
  The local Lipschitz continuity of $R_{F}$ follows by similar arguments.
\end{proof}

Let $g$ be a distribution belonging to $\CC^{\beta}$ for some $\beta <0$. Then
the map $f \mapsto f \circ g$ behaves, modulo smoother correction terms, like
a derivative operator:

\begin{lemma}\label{lemma:para-taylor}
  Let $\alpha \in ( 0,1 )$ and $\beta <0$ be such
  that $\alpha + \beta <0$ and $2 \alpha + \beta >0$, and let $F \in C^{2}$.
  Then there exists a locally bounded map $\Pi_{F} \colon \CC^{\alpha} \times \CC^{\beta} \rightarrow \CC^{2 \alpha + \beta}$ such that
  \begin{equation}\label{eq:para-taylor}
      F ( f ) \circ g=F' ( f ) ( f \circ g ) + \Pi_{F} ( f,g )
  \end{equation}
  for all $f \in \CC^{\alpha}$ and all smooth $g$. More precisely, we have
  \[
     \| \Pi_{F} ( f,g ) \|_{2 \alpha + \beta} \lesssim (\mathcal{M}_{F,0} ( f ) +\mathcal{M}_{F' ,1} ( f ) \| f \|_{\alpha}^{2} ) \| g \|_{\beta}.
  \]
  If $F \in C^{3}$, then the map $\Pi_{F}$ is locally Lipschitz continuous:
  \[
     \| \Pi_{F} ( f,g ) - \Pi_{F} ( u,v ) \|_{2 \alpha + \beta} \lesssim (\mathcal{M}_{F' ,2} ( f ) +\mathcal{M}_{F' ,2} ( u ) ) ( \| f-u \|_{\alpha} + \| f-u \|_{\alpha}^{2} ) \| g-v \|_{\beta}.
  \]
\end{lemma}

\begin{proof}
  Just use the paralinearization and commutator lemmas above to deduce that
  \[
     F ( f ) \circ g= ( F' ( f ) \prec f ) \circ g+R_{F} ( f ) \circ g=F' ( f ) ( f \circ g ) +C ( F' ( f ) ,f,g ) +R_{F} ( f ) \circ g
  \]
  and set
  \[
     \Pi_{F} ( f,g ) =C ( F' ( f ) ,f,g ) +R_{F} ( f ) \circ g.
  \]
  Since $\alpha + \beta < 0$ and $2 \alpha + \beta >0$, the operator $\Pi_{F} ( f,g )$ is well defined
  and maps into $\CC^{2 \alpha + \beta}$, and the claimed bounds follow
  easily.
\end{proof}

\begin{corollary}\label{cor:para-taylor product}
  Let $\alpha \in ( 0,1 )$ and $\beta <0$ be such that $2 \alpha + \beta >0$.
  Then there exists a locally bounded, locally Lipschitz continuous map
  $\Pi_{\times} \colon \CC^{\alpha} \times \CC^{\alpha} \times \CC^{\beta} \rightarrow \CC^{2 \alpha + \beta}$, such that
  \[
     ( u v ) \circ w=u ( v \circ w ) +v ( u \circ w ) + \Pi_{\times} ( u,v,w )
  \]
  for all $u,v \in \CC^{\alpha} ( \mathbb{R} )$ and all smooth $g$. More
  precisely, we have
  \begin{align*}
     \| \Pi_{\times} ( u,v,w ) \|_{2 \alpha + \beta} &\lesssim ( \| u \|_{\alpha} + \| v \|_{\alpha} + \| u \|_{\alpha}^{2} + \| v \|_{\alpha}^{2} ) \| w \|_{\beta} \qquad \text{and} \\
     \| \Pi_{\times} ( u,v,w ) - \Pi_{\times} ( \tilde{u} , \tilde{v} ,  \tilde{w} ) \|_{2 \alpha + \beta} & \lesssim ( \| u- \tilde{u} \|_{\alpha} + \| v- \tilde{v} \|_{\alpha} + \| u- \tilde{u} \|_{\alpha}^{2} + \| v-     \tilde{v} \|_{\alpha}^{2} ) \| w- \tilde{w} \|_{\beta} .
  \end{align*}
\end{corollary}

\begin{proof}
  It suffices to consider the smooth function $G( x,y ) =x y$ and to apply
  Lemma~\ref{lemma:para-taylor} with $F=G$, $f= ( u,v )$, and $g=w$.
\end{proof}



\section{Rough differential equations}\label{sec:ODE}

Let us now resume the analysis of Section~\ref{sec:paracontrolled ode}. We want to study the \textsc{rde}
\begin{equation}\label{eq:rde} 
  \partial_{t} u=F ( u ) \xi, \qquad u(0) = u_0,
\end{equation}
where $u_0 \in \R^d$, $u\colon \mathbb{R} \rightarrow \mathbb{R}^{d}$ is a continuous vector
valued function, $\xi\colon \mathbb{R} \rightarrow \mathbb{R}^{n}$ is a vector valued distribution with values in
$\CC^{\alpha -1}$ for some $\alpha \in ( 1/3,1 )$, and
$F\colon \mathbb{R}^{d} \rightarrow \Lcal ( \mathbb{R}^{n} ,\mathbb{R}^{d})$ is a family of vector fields on $\mathbb{R}^{d}$.

%A natural approach is to
%understand this equation as limit of the classical ODEs
%\begin{equation}\label{eq:rde-reg}
%  \partial_{t} u^{\varepsilon} = F ( u^{\varepsilon} ) \xi^{\varepsilon}
%\end{equation}
%for smooth approximations $\xi^{\varepsilon}$ of $\xi$ such that
%$\xi^{\varepsilon} \rightarrow \xi$ in $\CC^{\alpha -1}$ as $\varepsilon \rightarrow 0$. Thus, in order to pass to the limit, we are looking for a
%priori estimates on $u^{\varepsilon}$ which require only a control on the
%$\CC^{\alpha -1}$ norm of $\xi$.
%
%To avoid cumbersome notation, we will work at the level of
%equation~{\eqref{eq:rde}} for smooth $\xi$, where it should be understood that
%our aim is to get a priori estimates for the solution, in order to safely pass
%to the limit and extend the solution map to a larger class of data. The
%natural regularity of $u$ is $\CC^{\alpha}$, since $u$ should gain one
%derivative with respect to $F ( u ) \xi$, which will not behave better than
%$\xi$, and will therefore be in $\CC^{\alpha -1}$. The paraproduct
%decomposition suggests to write the right hand side of {\eqref{eq:rde}} as a
%sum of the three terms
%\begin{equation}\label{eq:para-rhs}
%  \underbrace{F ( u ) \prec \xi}_{\alpha -1} + \underbrace{F ( u ) \circ  \xi}_{2 \alpha -1} + \underbrace{F ( u ) \succ \xi}_{2 \alpha -1}
%\end{equation}
%(where the quantity indicated by the underbrace corresponds to the expected regularity of each term). Of course, unless $2 \alpha -1>0$, the resonant term
%$F ( u ) \circ   \xi$ cannot be controlled using only the $\CC^{\alpha}$--norm
%of $u$ and the $\CC^{\alpha -1}$--norm of $\xi$. If $F$ is at least $C^{2} (\mathbb{R}^{d} , \Lcal ( \mathbb{R}^{n} ,\mathbb{R}^{d} ) )$, we can
%use Lemma~\ref{lemma:para-taylor}, to rewrite this term as
%\begin{equation}\label{eq:rde expand Fu}
%   F ( u ) \circ \xi =F' ( u )  ( u \circ \xi ) + \Pi_{F} ( u, \xi ),
%\end{equation}
%where the remainder $\Pi_{F} ( u, \xi )$ is well defined under the condition
%$3 \alpha -1>0$, provided that $F' ( u ) ,u \in \CC^{\alpha}$, and $\xi \in \CC^{\alpha -1}$. In this case, it belongs to $\CC^{3 \alpha -1}$. The
%difficulty is now localized in the linearized resonant product $u \circ \xi$. In order to control this term, we would like to exploit the fact that
%the function $u$ is not a generic element of $\CC^{\alpha}$ but it has a
%specific structure, since its derivative $\partial_{t} u$ has to match the
%paraproduct decomposition given in~\eqref{eq:para-rhs}. Thus, we postulate
%that the solution $u$ is given by the following \tmtextit{paracontrolled} ansatz:
%\[
%   u=u^{\vartheta} \prec \vartheta +u^{\sharp},
%\]
%where $u^{\vartheta} , \vartheta \in \CC^{\alpha}$ and the remainder
%$u^{\sharp}$ is in $\CC^{2 \alpha}$. Together with the commutator estimate
%Lemma~\ref{lemma:commutator}, this decomposition allows for a finer analysis
%of the resonant term $u \circ \xi$. Indeed, we have
%\begin{equation}\label{eq:rde expand resonant term}
%   u \circ   \xi = ( u^{\vartheta} \prec \vartheta ) \circ   \xi +u^{\sharp} \circ   \xi =u^{\vartheta} ( \vartheta \circ \xi ) +C ( u^{\vartheta} , \vartheta , \xi ) +u^{\sharp} \circ   \xi ,
%\end{equation}
%where we recall that the commutator is defined by $C ( u^{\vartheta} ,\vartheta , \xi ) = ( u^{\vartheta} \prec \vartheta ) \circ \xi -u^{\vartheta}( \vartheta \circ \xi )$. Observe now that the term $u^{\sharp} \circ   \xi$
%does not pose any further problem, as it is bounded in $\CC^{3 \alpha -1}$.
%Moreover, the commutator is a bounded linear function of its arguments as
%long as the sum of their regularities adds up to a positive number, see
%Lemma~\ref{lemma:commutator} above. By assumption, we have $3 \alpha -1>0$,
%and therefore $C ( u^{\vartheta} , \vartheta , \xi ) \in \CC^{3 \alpha -1}$.
%The only problematic term which remains to be handled is thus $( \vartheta \circ   \xi )$. Here we need to make the assumption that $( \vartheta \circ \xi ) \in \CC^{2 \alpha -1}$ in order for the product $u^{\vartheta} (\vartheta \circ   \xi )$ to be well defined. This assumption is not guaranteed
%by the analytical estimates we have at hand, and it has to be added as a
%further requirement. Granting this, we have obtained that the right hand side
%of equation~{\eqref{eq:rde}} is well defined and a continuous function of $(u,u^{\vartheta} ,u^{\sharp} , \vartheta , \xi , \vartheta \circ \xi )$.
%
%The paracontrolled ansatz and the Leibniz rule for the paraproduct imply
%that~{\eqref{eq:rde}} can be rewritten as
%\[
%   \partial_{t} u= \partial_{t} ( u^{\vartheta} \prec \vartheta +u^{\sharp} ) = \partial_{t} u^{\vartheta} \prec \vartheta +u^{\vartheta} \prec \partial_{t} \vartheta + \partial_{t} u^{\sharp} =F ( u ) \prec \xi +F ( u ) \circ   \xi +F ( u ) \succ \xi.
%\]
%If we choose $\vartheta$ such that $\partial_{t} \vartheta = \xi$ and
%$\vartheta ( 0 ) =0$ and we set $u^{\vartheta} =F ( u )$, then we can
%use~{\eqref{eq:rde expand Fu}} and~{\eqref{eq:rde expand resonant term}} to
%obtain the following equation for the remainder $u^{\sharp}$:
%\begin{align*}%\label{eq:Phi definition} 
%  \partial_{t} u^{\sharp} & = F' ( u ) F ( u ) (\vartheta \circ \xi ) +F ( u ) \succ \xi - ( \partial_{t} F ( u ) \prec \vartheta ) \\
%  &\qquad +F' ( u ) C ( F ( u ) , \vartheta , \xi ) +F' ( u )  ( u^{\sharp} \circ \xi) + \Pi_{F} ( u, \xi ).
%\end{align*}
%Together with the equation
%\[
%   u=F ( u ) \prec \vartheta +u^{\sharp},
%\]
%this completely describes the solution.
In order to obtain concrete estimates,
we have to localize the equation. Therefore, we introduce a smooth cut-off
function $\varphi$ with support on $[ -2,2 ]$, which is equal to 1 on $[ -1,1]$. For $T>0$ we set $\varphi_{T} ( x ) = \varphi ( x/T )$ and modify the equation as
\[
   \partial_{t} u= \varphi_{T} F ( u ) \xi, \qquad u(0) = u_0.
\]
In the regular setting, if $u$ is solution to this equation, it is also a
solution of the original equation on $[ -T,T ]$, and thus it is sufficient to
study the last equation for local bounds. To avoid problems with the fact
that the paraproduct is a (mildly) non-local operation, we modify the
paracontrolled ansatz as follows:
\begin{equation}\label{eq:ode paracontrolled ansatz localized}
   u= \varphi_{T} ( F ( u )\prec \vartheta ) +u^{\sharp}.
\end{equation}
If $F \in C^2$, an easy computation gives
\begin{align*}
   \partial_{t} u^{\sharp} & = \varphi_{T} F ( u ) \xi - ( \partial_{t} \varphi_{T} )  ( F ( u ) \prec \vartheta ) - \varphi_{T}  ( \partial_{t} F( u ) \prec \vartheta ) - \varphi_{T} ( F ( u ) \prec \xi )\\
   & = \varphi_{T} \left[ ( F ( u ) \succ \xi ) +F' ( u )  ( ( u-u_{0} ) \circ \xi ) + \Pi_{F_{u_{0}}} ( u-u_{0} , \xi ) - ( \partial_{t} F ( u ) \prec \vartheta ) \right] \\
   &\qquad - ( \partial_{t} \varphi_{T} )  ( F ( u ) \prec \vartheta ) ,
\end{align*}
where we set $F_{u_{0}} ( x ) =F ( u_{0} +x )$ and used that $( F_{u_{0}} )' (x-u_{0} ) =F' ( x )$ for all $x \in \R^d$. We subtract the contribution of the initial condition,
because this will eventually allow us to solve the equation on a small
interval whose length does not depend on $u_0$. If we plug in
the modified paracontrolled ansatz for $u$, then $F' ( u ) ( ( u-u_{0} ) \circ \xi )$ becomes
\[
   F' ( u )  ( ( u-u_{0} ) \circ \xi ) =F' ( u ) ( ( \varphi_{T} ( F ( u ) \prec \vartheta ) ) \circ \xi )  +F' ( u ) ( ( u^{\sharp} -u_{0} ) \circ \xi ).
\]
For the first term on the right hand side we can further use that
\[
   ( \varphi_{T} ( F ( u ) \prec \vartheta ) ) \circ \xi = \varphi_{T} ( ( F (u ) \prec \vartheta ) \circ \xi ) + ( F ( u ) \prec \vartheta ) (\varphi_{T} \circ \xi ) + \Pi_{\times} ( \varphi_{T} ,F ( u ) \prec   \vartheta , \xi ),
\]
where we recall that $\Pi_\times$ was defined in Corollary~\ref{cor:para-taylor product}. So finally, an application of our commutator lemma yields
\begin{align}\label{eq:ode phi sharp modified def} \nonumber
   \partial_{t} u^{\sharp} & = \varphi_{T} \bigg[ ( F ( u ) \succ \xi ) + \Pi_{F_{u_{0}}} ( u-u_{0} , \xi ) +F' ( u ) ( ( u^{\sharp} -u_{0} ) \circ \xi ) + ( F ( u ) \prec \vartheta ) ( \varphi_{T} \circ \xi ) \\ \nonumber
   &\hspace{40pt} + \Pi_{\times} ( \varphi_{T} ,F ( u ) \prec \vartheta , \xi ) + \varphi_{T} C ( F ( u ) , \vartheta , \xi ) +F' ( u ) \varphi_{T} F ( u ) ( \vartheta \circ \xi ) - ( \partial_{t} F ( u ) \prec \vartheta ) \bigg] \\ \nonumber
   &\qquad - ( \partial_{t} \varphi_{T} )  ( F ( u ) \prec \vartheta ) \\
   & = \varphi_{T} \Phi^{\sharp} - (\partial_{t} \varphi_{T} )  ( F ( u ) \prec \vartheta ),
\end{align}
where $\Phi^\sharp$ is defined to be the term in the square large brackets. Let us summarize our observations so far.

\begin{lemma}
  Let $\xi$ be a smooth path, let $\vartheta$ be such that $\partial_{t}  \vartheta = \xi$, and let $F \in C^{2}_{b}$. Then
  $u$ solves the ODE
  \[
     \partial_{t} u= \varphi_{T} F ( u ) \xi, \qquad u ( 0 ) =u_{0},
  \]
  if and only if $u= \varphi_{T} ( F (u ) \prec \vartheta ) +u^{\sharp}$, where $u^{\sharp}$ solves
  \[
     \partial_{t} u^{\sharp} = \varphi_{T} \Phi^{\sharp} - ( \partial_{t} \varphi_{T} )  ( F ( u ) \prec \vartheta ), \qquad u^{\sharp} ( 0 ) =u_{0} - \varphi_{T} ( F ( u )\prec \vartheta ) ( 0 ),
  \]
  and where $\Phi^{\sharp}$ is defined in~\eqref{eq:ode phi sharp modified def}. Moreover, for $\alpha \in ( 1/3,1/2 )$ we have the estimate
  \[
     \| \Phi^{\sharp} \|_{2 \alpha -1} \lesssim C_{F} C_{\xi} ( 1+ \| u-u_{0}  \|_{\alpha} + \| u-u_{0} \|^{2}_{\alpha} + \| u^{\sharp} -u_{0} \|_{2 \alpha} ),
  \]
  where
  \[
     C_{\xi} = \| \xi \|_{\alpha -1} + \| \vartheta \|_{\alpha} + \| \vartheta \circ   \xi \|_{2 \alpha -1} + \| \vartheta \|_{\alpha} \| \xi \|_{\alpha -1} \qquad \tmop{and} \qquad C_{F} = \| F \|_{C^{2}_{b}} + \| F \|_{C^{2}_{b}}^{2} .
  \]   
\end{lemma}

The estimate for $\Phi^{\sharp}$ follows from a somewhat lengthy but
elementary calculation based on the decomposition~\eqref{eq:ode phi sharp modified def}, where we estimate the $L^{\infty}$ norm rather than the
$\CC^{2 \alpha -1}$ norm for each term where this is possible.

Plugging in the correct initial condition for $u^{\sharp}$ leads to
\begin{align*}
   u^{\sharp} ( t ) &= u_{0} - ( F ( u ) \prec \vartheta ) ( 0 ) + \int_{0}^{t}\partial_{s} u^{\sharp} ( s ) \mathd s \\
   & =u_{0} - ( F ( u ) \prec \vartheta ) ( 0 ) + \int_{0}^{t} \varphi_{T} \Phi^{\sharp} ( s ) \mathd s- \int_{0}^{t} ( \partial_{s} \varphi_{T} ) ( s )  ( F ( u ) \prec \vartheta ) ( s ) \mathd s.
\end{align*}
Now Lemma~\ref{lemma:integral} gives estimates for the integrals appearing on
the right hand side in terms of distributional norms of the integrands, and we
obtain the bound
\begin{align*}
   \| u^{\sharp} -u_{0} \|_{2 \alpha} & \leqslant c T \| F ( u ) \prec \vartheta \|_{2 \alpha -1} +c T \| \Phi^{\sharp} \|_{2 \alpha -1} \\
   & \leqslant c T C_{F} C_{\xi} ( 1+ \| u-u_{0} \|_{\alpha} + \| u-u_{0} \|^{2}_{\alpha} + \| u^{\sharp} -u_{0} \|_{2 \alpha} ),
\end{align*}
where $c>0$ denotes a generic constant whose value may change from line to
line. Using that $u= \varphi_{T} (F ( u ) \prec \vartheta ) +u^{\sharp}$, we moreover have
\[
   \| u-u_{0} \|_{\alpha} \leqslant c  \| F \|_{L^{\infty}} \| \vartheta \|_{2 \alpha} + \| u^{\sharp} -u_{0} \|_{2 \alpha}.
\]
Since $\| F \|_{L^{\infty}} \leqslant C_{F}$, we deduce from these two
estimates that if $C_{F}$ is small enough (depending only on $C_{\xi}$ and $T$
but not on $| u_{0} |$), then $\| u^{\sharp} \|_{2 \alpha} \leqslant | u_{0} | +1$. 
This is the required uniform estimate on the problem.



Similarly we can show that if $F \in C^{3}_{b}$ and if $\| F \|_{C^{3}_{b}}$
is small enough, then the map
\[
   (u_0, \xi , \vartheta , \xi \circ   \vartheta) \mapsto ( u,u^{\sharp} )
\]
is locally Lipschitz continuous from $\CC^{\alpha -1} \times \CC^{\alpha}\times \CC^{2 \alpha -1} \times \mathbb{R}^{d}$ to $\CC^{\alpha} \times \CC^{2 \alpha -1}$. To summarize:

\begin{lemma}\label{lemma:ode conditional existence uniqueness}
  Let $a>0$ and let $\| F\|_{C^{3}_{b}}$ be small enough (depending on $a$). Let $\xi$, $\vartheta$,
  and $\varphi$ be smooth functions with $\xi = \partial_{t} \vartheta$ and
  such that $\varphi$ has compact support. If $\alpha >1/3$ and
  \begin{equation}\label{eq:bounds on ode data}
     \max \{ \| \xi \|_{\alpha -1} , \| \vartheta \|_{\alpha} , \| \xi \circ \vartheta \|_{2 \alpha -1} ,T, \| \varphi \|_{C^{1}_{b}} \} \leqslant a,
  \end{equation}
  then for every $u_0 \in \R^d$ there exists a unique global solution $u$ to
  \[
     \partial_{t} u= \varphi_{T} F ( u ) \xi, \qquad u(0) = u_0.
  \]
  Moreover, for fixed $T$,
  $\varphi$, and $F$, the solution $u$ depends in a Lipschitz continuous way on $(u_0, \xi, \vartheta , \xi \circ \vartheta)$ satisfying~\eqref{eq:bounds on ode data}.
\end{lemma}

In order to ensure that $\| F \|_{C^{3}_{b}}$ is small enough we can use a
dilation argument. Recall that the scaling operator $\Lambda_\lambda$ is defined for $\lambda > 0$ by $\Lambda_\lambda u = u(\lambda \cdot)$. If we let $u^{\lambda} = \Lambda_{\lambda} u$ and
$\xi^{\lambda} = \lambda^{1- \alpha} \Lambda_{\lambda} \xi$ for $\lambda >0$,
then $u^{\lambda}$ solves
\[
   \partial_{t} u^{\lambda} = \lambda^{\alpha} F ( u^{\lambda} ) \xi^{\lambda}, \qquad u^\lambda(0) = u_0.
\]
The rescaling of $\xi^{\lambda}$ is chosen so that its $\CC^{\alpha}$ norm is
uniformly bounded by that of $\xi$ as $\lambda \rightarrow 0$. Indeed,
Lemma~\ref{lemma:scaling} yields
\[
   \| \xi^{\lambda} \|_{\alpha -1} = \lambda^{1- \alpha} \| \Lambda_{\lambda} \xi \|_{\alpha -1} \lesssim ( 1+ \lambda^{1- \alpha} ) \| \xi \|_{\alpha -1} \lesssim \| \xi \|_{\alpha -1}
\]
for $\lambda \leqslant 1$. If moreover we let $\vartheta^{\lambda} = \lambda^{- \alpha} \Lambda_{\lambda} \vartheta$, then $\| \vartheta^{\lambda} \circ \xi^{\lambda} \|_{2 \alpha -1} \lesssim \| \vartheta \circ   \xi \|_{2 \alpha -1}$. Thus, we deduce from Lemma \ref{lemma:ode conditional existence uniqueness} that for fixed $T>0$ and $\varphi$
with compact support, there exists $\lambda >0$, such that for all $u_{0} \in \mathbb{R}^{d}$ we have a unique
global solution $u^{\lambda}$ to
\[
   \partial_{t} u^{\lambda} = \varphi_{T} \lambda^{\alpha} F ( u^{\lambda} ) \xi^{\lambda}, \qquad u^\lambda(0) = u_0.
\]
The rescaled problem is
equivalent to the original one upon the change $F \rightarrow \lambda^{\alpha} F$, $\xi \rightarrow \xi^{\lambda}$ and $\vartheta \circ   \xi \rightarrow \vartheta^{\lambda} \circ   \xi^{\lambda}$. So if we set $u= \Lambda_{\lambda^{-1}} u^{\lambda}$, then $u$ is the unique global solution to
\[
   \partial_{t} u= \varphi_{\lambda T} F ( u ) \xi, \qquad u(0) = u_0.
\]
In particular, if $\varphi \equiv 1$
on $[ -1,1 ]$, then $u$ is the unique solution to the original \textsc{rde} in the
interval $[ - \lambda T, \lambda T ]$.
Since $\lambda$ can be chosen independently of $u_{0}$, we can now iterate on
intervals of length $2\lambda T$, and obtain a global solution $u \in \CC^{\alpha}_{\tmop{loc}}$.



This analysis can be summarized in the following statement.

\begin{theorem}\label{theorem:rde existence uniqueness}
  Let $\alpha >1/3$. Assume that $(\xi^{\varepsilon} )_{\varepsilon >0}$ is a family of smooth functions with
  values in $\mathbb{R}^{n}$, $(u^\varepsilon_0)$ is a family of initial conditions in $\R^d$, and $F$ is a family of $C^{3}_{b}$ vector
  fields on $\mathbb{R}^{d}$. Suppose that there exist $u_0 \in \R^d$, $\xi \in \CC^{\alpha -1}$ and $\eta \in \CC^{2 \alpha -1}$ such that $(u_0^\varepsilon, \xi^{\varepsilon} , \vartheta^\varepsilon, (\vartheta^{\varepsilon} \circ   \xi^{\varepsilon} ) )$ converges to $(u_0, \xi , \vartheta, \eta )$ in $\CC^{\alpha -1} \times \CC^\alpha \times \CC^{2 \alpha -1}$, where
  $\vartheta^{\varepsilon}$ and $\vartheta$ are solutions to $\partial_{t}  \vartheta^{\varepsilon} = \xi^{\varepsilon}$ and $\partial_t \vartheta = \xi$, respectively. Let for $\varepsilon >0$ the function $u^{\varepsilon}$ be the
  unique global solution to the Cauchy problem
  \[
     \partial_{t} u^{\varepsilon} =F ( u^{\varepsilon} ) \xi^{\varepsilon} , \hspace{2em} u^{\varepsilon} ( 0 ) =u^\varepsilon_{0} .
  \]
  Then there exists $u \in \CC^{\alpha}_\mathrm{loc}$ such that $u^{\varepsilon} \rightarrow u$ in $\CC^{\alpha}_{\tmop{loc}}$ as $\varepsilon \rightarrow 0$. The limit $u$ depends only on $( u_{0} , \xi , \vartheta , \eta )$, and not on the approximating family $(u_0^\varepsilon, \xi^{\varepsilon} , \vartheta^\varepsilon,  ( \vartheta^{\varepsilon} \circ   \xi^{\varepsilon} ) )$.
\end{theorem}

\begin{proof}
  The only point which remains to be shown is the convergence of $(u^{\varepsilon} )$ to $u$ in $\CC^{\alpha}_{\tmop{loc}}$. A priori,
  we only know that for sufficiently small $\lambda >0$, the solutions
  $\tilde{u}^{\varepsilon}$ to $\partial_{t} \tilde{u}^{\varepsilon} = \varphi_{\lambda T} F ( \tilde{u}^{\varepsilon} ) \xi^{\varepsilon}$ with
  $\tilde{u}^{\varepsilon} ( 0 ) =u_{0}$ converge, as $\varepsilon \rightarrow 0$, in $  \CC^{\alpha}$ to a unique limit $\tilde{u}$.
  But since $\varphi_{\lambda T} \equiv 1$ on $[ - \lambda T, \lambda T ]$, we
  have $\tilde{u}^{\varepsilon} |_{[ - \lambda T, \lambda T ]} =u^{\varepsilon} |_{[ - \lambda T, \lambda T ]}$. So if we define $u|_{[ -  \lambda T, \lambda T ]} = \tilde{u} |_{[ - \lambda T, \lambda T ]}$, then
  $u|_{[ - \lambda T, \lambda T ]}$ does not depend on $\varphi_{\lambda T}$.
  Moreover, for every $\psi \in \CD$ with support contained in $[ - \lambda T, \lambda T ]$, we also have that $\| \psi ( u^{\varepsilon} -u ) \|_{\alpha}$ converges to zero as $\varepsilon \rightarrow 0$. Now we can
  iterate this construction of $u$ on intervals of length $2 \lambda T$. We
  end up with a distribution $u \in \CS'$, which only depends
  on $( u_{0} ,F, \xi , \vartheta , \eta )$, but not on $\varphi_{\lambda T}$
  or on the approximating sequence $(u_0^\varepsilon, \xi^{\varepsilon} , \vartheta^{\varepsilon} , \xi^{\varepsilon} \circ \vartheta^{\varepsilon} )_{\varepsilon >0}$. If $\psi \in \CD$, then it can be written as a
  finite sum of smooth functions with support contained in intervals of length
  $2 \lambda T$, and therefore $\psi u= \lim_{\varepsilon \rightarrow} \psi u^{\varepsilon}$, where convergence takes places in $\CC^{\alpha}$.
\end{proof}

\begin{remark}\label{remark:ode initial condition theta}
  Note that we did not fix the initial conditions $\vartheta^{\varepsilon}(0)$ and $\vartheta(0)$. It turns out that they do not have any influence on the convergence: If $\partial_t \vartheta^\varepsilon = \xi^\varepsilon$, $\vartheta^\varepsilon(0) = 0$, and
  $\partial_{t} \tilde{\vartheta}^{\varepsilon} = \xi^{\varepsilon}$,
  $\tilde{\vartheta}^{\varepsilon} ( 0 ) =c \in \mathbb{R}^d$, then
  $\tilde{\vartheta}^{\varepsilon} - \vartheta^{\varepsilon} \equiv c$, so
  that
  \begin{align*}
     \| \tilde{\vartheta}^{\varepsilon} \circ \xi^{\varepsilon} - ( \eta +c \circ \xi ) \|_{2 \alpha -1} & \lesssim \| c \circ ( \xi^{\varepsilon} - \xi ) \|_{2 \alpha -1} + \| \vartheta^{\varepsilon} \circ
     \xi^{\varepsilon} - \eta \|_{2 \alpha -1} \\
     & \lesssim c \| \xi^{\varepsilon} - \xi \|_{\alpha -1} + \| \vartheta^{\varepsilon} \circ \xi^{\varepsilon} - \eta \|_{2 \alpha -1} .
  \end{align*}
  In other words, $( \xi^{\varepsilon} , \vartheta^\varepsilon, \vartheta^{\varepsilon} \circ \xi^{\varepsilon} )$ converges to $(\xi , \vartheta, \eta)$ in
  $\CC^{\alpha -1} \times \CC^\alpha \times \CC^{2 \alpha -1}$ as $\varepsilon \rightarrow 0$ if and only if $  ( \xi^{\varepsilon} , \tilde{\vartheta}^\varepsilon, \tilde{\vartheta}^{\varepsilon} \circ \xi^{\varepsilon} )$ converges to $(\xi , \vartheta, \eta +c \circ \xi )$ in $\CC^{\alpha -1} \times \CC^\alpha \times \CC^{2 \alpha -1}$, and in particular the solution $u$ in Theorem~\ref{theorem:rde existence uniqueness} does not depend on the particular choice of initial conditions for $\vartheta$.
\end{remark}

\begin{remark}
   If $\xi$ has compact support, then the $\CC^\alpha$--norm of $\vartheta$ can be controlled in terms of the $\CC^{\alpha-1}$--norm of $\xi$, so that it suffices to have convergence of the pairs $(\xi^\varepsilon, \vartheta^\varepsilon \circ \xi^\varepsilon)$, rather than the triples $( \xi^{\varepsilon} , \vartheta^{\varepsilon} , \xi^{\varepsilon} \circ \vartheta^{\varepsilon} )$.
\end{remark}

\begin{remark}
  If we only suppose $F \in C^{3}$ and not that $F$ and its derivatives are
  bounded, we still obtain local existence and uniqueness of solutions. In
  that case we may consider a function $G \in C^{3}_{b}$ that coincides with
  $F$ on $\{ | x | \leqslant a \}$ for some $a> | u_{0} |$. The Cauchy problem
  \[
     \partial_{t} v=G ( v ) \xi , \hspace{2em} v ( 0 ) =u_{0},
  \]
  then has a unique global solution in the sense of Theorem~\ref{theorem:rde
  existence uniqueness}. If we stop $v$ upon leaving the set $\{ | x | \leqslant a \}$, we obtain a local solution to the \textsc{rde} with vector field
  $F$.
\end{remark}


\subsection{Interpreting the solution to an RDE}

  So far we showed that under the assumptions of Theorem~\ref{theorem:rde existence uniqueness} there exists a unique limit $u$ of the solutions to the
  regularized equations, which does not depend on the particular approximating
  sequence. In that sense, one may formally call $u$ the unique solution to
  \[
     \partial_{t} u=F ( u ) \xi , \hspace{2em} u ( 0 ) =u_{0} .
  \]
  But $u$ is actually a weak solution to the equation if we interpret the product $F(u) \xi$ appropriately. Below we will introduce a map $\pi_F$ which extends the pointwise product $F(u) \xi$ from smooth $\xi$ to $\xi \in \CC^{\alpha-1}$ by a continuity argument. But first we present an auxiliary result which shows that the considered topologies and operators do not depend on the particular dyadic
partition of unity that we use to describe them.


\begin{lemma}\label{l:paraproduct difference is bounded operator}
   Let $\alpha, \beta \in \R$. Let $( \chi, \rho )$ and $( \tilde{\chi}, \tilde{\rho} )$ be two dyadic partitions of unity. Let $(\Delta_i)_{i \ge -1}$ be the Littlewood-Paley blocks corresponding to $( \chi, \rho )$, and $(\widetilde{\Delta}_i)_{i \ge -1}$ those corresponding to $( \tilde{\chi}, \tilde{\rho} )$. Let $(u \prec v, u \succ v, u \circ v)$ denote paraproducts and resonant term defined in terms of $(\Delta_i)_{i \ge -1}$, and let $(u \ \widetilde{\prec}\ v, u \ \widetilde{\succ} \ v, u\ \widetilde{\circ}\ v)$ denote paraproducts and resonant term defined in terms of $(\widetilde{\Delta}_i)_{i \ge -1}$.
%   \begin{gather*}
%      u\prec v = \sum_{i < j - N} \Delta_i u \Delta_j v \qquad \text{and} \qquad u \ \widetilde{\prec}\ v = \sum_{i < j- \tilde{N}} \widetilde{\Delta}_i u \widetilde{\Delta}_j v,\\
%      u \circ v = \sum_{| i - j | \leqslant N} \Delta_i u \Delta_j v \qquad \text{and} \qquad u\ \widetilde{\circ}\ v = \sum_{|i-j|\leqslant \tilde{N}} \widetilde{\Delta}_i u \widetilde{\Delta}_j v,\\
%      u \succ v = \sum_{j < i - N} \Delta_i u \Delta_j v \qquad \text{and} \qquad  u \ \widetilde{\succ} \ v = \sum_{j < i- \tilde{N}} \widetilde{\Delta}_i u \widetilde{\Delta}_j v,
%   \end{gather*}
%   where $N$ is large enough so that $\Delta_i \Delta_j = 0$ for $| i -j | > N$, and similarly for $\tilde{N}$. 
Then
   \[
      (u,v) \mapsto ( u \prec v - u \ \widetilde{\prec} \ v, u \circ v - u \ \widetilde{\circ} \ v, u \succ v - u \ \widetilde{\succ} \ v)
   \]
   is a bounded bilinear operator from $\CC^\alpha \times \CC^\beta$ to $(\CC^{\alpha+\beta})^3$.%, such that for all  $u \in \CC^{\alpha}$, $v \in \CC^{\beta}$
%   \begin{gather*}
%      \lVert \pi_{<} ( u, v ) - \widetilde{\pi}_{<} (u, v) \rVert_{\alpha + \beta} \lesssim \lVert u \rVert_{\alpha} \lVert v \rVert_{\beta},\\
%      \lVert \pi_{\circ} ( u, v ) - \widetilde{\pi}_{\circ} (u, v) \rVert_{\alpha + \beta} \lesssim \lVert u \rVert_{\alpha} \lVert v \rVert_{\beta},\\
%      \lVert \pi_{>} ( u, v ) - \widetilde{\pi}_{>} (u, v) \rVert_{\alpha + \beta} \lesssim \lVert u \rVert_{\alpha} \lVert v \rVert_{\beta}.
%   \end{gather*}
\end{lemma}

\begin{proof}
   The statement for $(u,v) \mapsto (u \prec v - u \ \widetilde{\prec} \ v)$ (and thus for $(u,v) \mapsto (u \succ v - u \ \widetilde{\succ} \ v)$) is shown in Bony~\cite{Bony1981}, Theorem 2.1. But for smooth functions $u$ and $v$ we have $u \circ v = uv - u \prec v - u \succ v$, and similarly for $u \ \widetilde{\circ} \ v$. Thus, the bound on $u \circ v - u \ \widetilde{\circ} \ v$ follows from the bounds on $u \prec v - u \ \widetilde{\prec} \ v$ and on $u \succ v - u \ \widetilde{\succ} \ v$ in combination with a continuity argument.
\end{proof}

\begin{theorem}\label{theorem:product operator}
  Let $\alpha \in ( 0,1 )$ and let $\beta <0$ be such that $\alpha + \beta < 0$ and $2 \alpha + \beta >0$. Let $F \in C^{2}$. Define the operator
  \[
     \pi_{F} \colon \CC^\alpha \times \CC^{\alpha} \times \CC^{2 \alpha}  \times \CC^{\alpha} \times \CC^{\beta} \times \CC^{\alpha + \beta} \rightarrow \CC^{\beta}
  \]
  by
  \begin{align*}
     \pi_{F} ( u,u^{v} ,u^{\sharp} ,v,w, \eta ) &= F ( u ) \prec w+F ( u ) \succ w+ \Pi_{F} ( u,w ) +F' ( u ) ( u^{\sharp} \circ w ) \\
     &\qquad +F' ( u ) C ( u^{v} ,v,w ) +F' ( u ) u^{v} \eta.
  \end{align*}
  Then $\pi_{F}$ is locally Lipschitz continuous. If $w$ is smooth and if $u^{\sharp} =u- ( u^{v} \prec v )$ and $\eta =v \circ w$, then $\pi_{F} ( u,u^{v} ,u^{\sharp} ,v,w, \eta ) =F ( u ) w$, where the right hand side denotes the pointwise product.
  
  If $u^{\sharp} =u- ( u^{v} \prec v)$ and if there exist sequences of smooth functions $( v^{n} )_{n \in \mathbb{N}}$ and $( w^{n} )_{n \in \mathbb{N}}$, such that $( v^{n} ,w^{n},v^{n} \circ w^{n} )$ converges to $( v,w, \eta )$ in $\CC^{\alpha}\times \CC^{\beta} \times \CC^{\alpha + \beta}$, then $\pi_{F}$ is independent of the particular dyadic partition of unity used to define it.
\end{theorem}

\begin{proof}
  The local Lipschitz continuity of $\pi_{F}$ follows from its definition in
  combination with Lemma~\ref{lemma:commutator},
  Lemma~\ref{lemma:para-taylor}, and the paraproduct estimates
  Lemma~\ref{thm:paraproduct}. %We also use that $\lVert \cdummy \rVert_{\beta} \lesssim \lVert \cdummy \rVert_{L^{\infty}}$, which is for example helpful in
  %estimating $\| F' ( u ) ( u^{\sharp} \circ w ) \|_{\beta}$, and that $\| u \|_{\alpha} \leqslant 1+ \| u \|_{\alpha}^{2}$.
  If $w$ is smooth and if $u^{\sharp} =u- ( u^{v} \prec v )$ and $\eta =v \circ w$, then
  \[
     F' ( u ) C ( u^{v} ,v,w ) +F' ( u ) u^{v} \eta =F' ( u ) ( ( u^{v} \prec v ) \circ w ) ,
  \]
  and therefore
  \[
     \Pi_{F} ( u,w ) +F' ( u ) ( u^{\sharp} \circ w ) +F' ( u ) C ( u^{v} ,v,w ) +F' ( u ) u^{v} \eta = \Pi_{F} ( u,w ) +F' ( u ) ( u \circ w ) =F ( u ) \circ w,
  \]
  leading to $\pi_{F} =F ( u ) \prec w+F ( u ) \succ w+F ( u ) \circ w = F (u) w$.
  
  Assume now that $u^{\sharp} =u- ( u^{v} \prec v)$ and that there exist sequences of smooth functions $( v^{n} )_{n \in \mathbb{N}}$ and $( w^{n} )_{n \in \mathbb{N}}$, such that $( v^{n} ,w^{n},v^{n} \circ w^{n} )$ converges to $( v,w, \eta )$ in $\CC^{\alpha}\times \CC^{\beta} \times \CC^{\alpha + \beta}$. Let $(\tilde{\chi},\tilde{\rho})$ be another dyadic partition of unity, and let $\widetilde{\prec}$, $\widetilde{\circ}$, $\widetilde{\succ}$ be the operators defined in terms of this partition. We set
   \begin{align*}
      P_1(u,u^v,v,v^n,w^n) &= F ( u ) \prec w^n + F ( u ) \succ w^n + \Pi_{F} ( u,w^n ) +F' ( u ) ( (u- ( u^{v} \prec v)) \circ w^n ) \\
     &\qquad +F' ( u ) C ( u^{v} ,v^n ,w^n ) +F' ( u ) u^{v} (v^n \circ w^n) \\
     & = F(u) w^n + F'(u) (u^v \prec( v^n - v)) \circ w^n.
   \end{align*}
   By continuity of the operators involved, $P_1(u,u^v,v,v^n,w^n)$ converges to $\pi_F(u,u^v,u^\sharp, v, w, \eta)$ as $n$ tends to $\infty$. Similarly we define $P_2(u,u^v,v,v^n,w^n) = F(u) w^n + F'(u) (u^v \ \widetilde{\prec} \ ( v^n - v)) \ \widetilde{\circ} \ w^n$. If $\widetilde{\pi}_F$ is defined like $\pi_F$, except using the dyadic partition $(\tilde{\chi}, \tilde{\rho})$ instead of $(\chi, \rho)$, then by Lemma~\ref{l:paraproduct difference is bounded operator} $P_2(u,u^v,v,v^n,w^n)$ converges to $\widetilde{\pi}_F(u,u^v,\tilde{u}^\sharp, v, w, \tilde{\eta})$ as $n$ tends to $\infty$, where
   \begin{align*}
      \tilde{u}^\sharp = u -  u^{v} \ \widetilde{\prec} \ v \qquad \text{and} \qquad \tilde{\eta}= \eta + u \prec v + u \succ v - u \ \widetilde{\prec} \ v - u \ \widetilde{\succ} \ v.
   \end{align*}
   For the difference $P_1 - P_2$, Lemma~\ref{l:paraproduct difference is bounded operator} implies that
   \begin{align*}
      \lVert (P_1-P_2) (u,u^v,v,v^n,w^n)\rVert_{\alpha+\beta} & \lesssim \lVert F'(u)\rVert_\alpha \lVert (u^v \prec (v^n - v) - u^v \ \widetilde{\prec} \ (v^n - v) ) \circ w^n ) \rVert_{2 \alpha + \beta} \\
      &\hspace{5pt} + \lVert F'(u)\rVert_\alpha \lVert (u^v \ \widetilde{\prec} \ (v^n - v) ) \circ w^n - (u^v \ \widetilde{\prec} \ (v^n - v) ) \ \widetilde{\circ} \ w^n \rVert_{\alpha + \beta}\\
      & \lesssim \lVert F'(u) \rVert_\alpha u^v\rVert_\alpha \lVert v^n - v\rVert_\alpha \lVert w^n \rVert_\beta,
   \end{align*}
   which converges to zero by assumption, and thus $\pi_F(u,u^v,u^\sharp, v, w, \eta) = \widetilde{\pi}_F(u,u^v,\widetilde{u}^\sharp, v, w, \widetilde{\eta})$.
\end{proof}
  
   Let now $\tilde{u}$ be the limit of the solutions $( \tilde{u}^{\varepsilon} )_{\varepsilon >0}$ to $\partial_{t} \tilde{u}^{\varepsilon} = \varphi_{\lambda T} F ( \tilde{u}^{\varepsilon} ) \xi^{\varepsilon}$ with
  $\tilde{u}^{\varepsilon} ( 0 ) =u_{0}$, as in the proof of Theorem
  \ref{theorem:rde existence uniqueness}. From {\eqref{eq:ode paracontrolled ansatz localized}} we know that $\tilde{u} - \varphi_{\lambda T} ( F ( \tilde{u} ) \prec \vartheta ) \in \CC^{2 \alpha}$, provided that
  $\lambda>0$ is sufficiently small. So if $\psi \in \CD$ has its
  support in $[ - \lambda T, \lambda T ]$, then $\psi ( u- ( F ( \tilde{u} ) \prec \vartheta ) ) \in \CC^{2 \alpha}$, because $u|_{[ - \lambda T, \lambda T ]} = \tilde{u} |_{[ - \lambda T, \lambda T ]}$. But using the
  increment characterization of H{\"o}lder-Besov regularity and the fact that
  $\tilde{u} |_{[ - \lambda T, \lambda T ]} =u|_{[ - \lambda T, \lambda T ]}$,
  one can show that $\psi ( ( F ( u ) \prec \vartheta ) - ( F ( \tilde{u} ) \prec \vartheta ) ) \in \CC^{2 \alpha}$. Hence, we have $\psi ( u- ( F ( u ) \prec \vartheta ) ) \in \CC^{2 \alpha}$ for every $\psi \in \CD$ with support contained in $[ - \lambda T, \lambda T ]$, and
  just as in the proof of Theorem~\ref{theorem:rde existence uniqueness}, this
  extends to general $\psi \in \CD$, so that $u- ( F ( u ) \prec \vartheta ) \in \CC^{2 \alpha}_{\tmop{loc}}$. Let us assume for
  simplicity that $\xi$ has compact support. Then $u$ and $\vartheta$ are constant outside of
  some finite interval, and therefore $u- ( F ( u ) \prec \vartheta ) \in \CC^{2 \alpha}$. Recall now the operator $\pi_{F}$ defined in
  Theorem~\ref{theorem:product operator}, and observe that $u$ is the unique weak
  solution to
  \[
     \partial_{t} u= \pi_{F} ( u,F ( u ) ,u- ( F ( u ) \prec \vartheta ) , \vartheta , \xi , \eta ) , \hspace{2em} u ( 0 ) =u_{0}.
  \]
  This follows
  from the continuity of $\pi_{F}$ and $\partial_t$, and because
  \[
     \partial_{t} u^{\varepsilon} =F ( u^{\varepsilon} ) \xi^{\varepsilon} =\pi_{F} ( u^{\varepsilon} ,F ( u^{\varepsilon} ) ,u^{\varepsilon} -F (u^{\varepsilon} \prec \vartheta^{\varepsilon} ) , \vartheta^{\varepsilon} , \xi^{\varepsilon} , \vartheta^{\varepsilon} \circ \xi^{\varepsilon} ) .
  \]
  The same holds true if $\xi$ does not necessarily have compact support. But
  then $u$ is only in $\CC^{\alpha}_{\tmop{loc}}$, so that we would first have to extend the operator $\pi_{F}$ to local Besov spaces, which is
  possible but takes some work.


\subsection{Alternative approach}

We briefly describe an alternative approach to \textsc{rde}s which avoids the
paracontrolled ansatz. The idea is to control $u \circ   \xi$ directly by
exploiting that $u$ solves the differential equation $\partial_{t} u=F ( u ) \xi$. Indeed, let as above $\vartheta$ be a solution to $\partial_{t} \vartheta = \xi$ and observe that the Leibniz rule yields
\[
   u \circ   \xi =u \circ   \partial_{t} \vartheta = \partial_{t} ( u \circ  \vartheta ) - \partial_{t} u \circ \vartheta = \partial_{t} ( u \circ   \vartheta ) - ( F ( u ) \xi ) \circ   \vartheta.
\]
Now the second term on the right hand side can be rewritten as
\begin{align*}
   ( F ( u ) \xi ) \circ   \vartheta & = ( F ( u ) \prec \xi ) \circ   \vartheta + ( F ( u ) \circ \xi ) \circ   \vartheta + ( F ( u ) \succ \xi ) \circ  \vartheta \\
   & =F ( u ) ( \xi \circ   \vartheta ) +C ( F ( u ) , \xi , \vartheta ) + ( F'( u ) ( u \circ   \xi ) ) \circ   \vartheta + \Pi_{F} ( u, \xi ) \circ  \vartheta + ( F ( u ) \succ \xi ) \circ   \vartheta.
\end{align*}
Combining these two equations, we see that
\begin{gather*}
   u \circ   \xi = \Phi - ( F' ( u ) ( u \circ   \xi ) ) \circ   \vartheta , \qquad \text{where} \\
   \Phi = \partial_{t} ( u \circ   \vartheta ) -F ( u ) ( \xi \circ \vartheta ) -C ( F ( u ) , \xi , \vartheta ) - \Pi_{F} ( u, \xi ) \circ \vartheta - ( F ( u ) \succ \xi ) \circ   \vartheta.
\end{gather*}
This is an implict equation for $u \circ   \xi$ which can be solved by fixed point methods. For example, it is easy to obtain the estimate
\[
   \| u \circ   \xi \|_{2 \alpha -1} \lesssim \| \Phi \|_{2 \alpha -1} +C_{F} \| u \circ   \xi \|_{2 \alpha -1} \| \vartheta \|_{\alpha},
\]
and if $C_{F}$ is small enough this leads to $\| u \circ   \xi \|_{2 \alpha -1} \lesssim \| \Phi \|_{2 \alpha -1}$. Moreover, we have $\| \Phi \|_{2 \alpha -1} \lesssim C_{\xi} [ \| u \|_{\alpha} +C_{F} ( 1+ \| u \|_{\alpha} )^{2} ]$. These estimates can be
reinjected into the equation
\[
   \partial_{t} u=F ( u ) \xi =F ( u ) \prec \xi +F' ( u ) ( u \circ   \xi ) +F ( u ) \succ \xi +R ( F' ( u ) ,u, \xi )
\]
to obtain a local estimate for $u$.

\subsection{Structure of solutions to RDEs}

In this section we would like to discuss how the combination of analytic and algebraic requirements generate very interesting phenomena in the context of irregular PDEs. We will discuss only the simple case of \textsc{rde}s but similar considerations apply also to the other models. These remarks are intentionally sketchy and have only a heuristic purpose, we plan to come back to them more systematically in a further publication.

\paragraph{The stable form of an RDE.}
As we have seen, the solution to the \textsc{rde}~{\eqref{eq:rde}} for smooth $\xi$ can
be understood as a regular function of $(u_0, \xi , \vartheta \circ   \xi ) \in \R^d \times \CC^{\alpha -1} ( \mathbb{R},\mathbb{R}^{n} ) \times \CC^{2 \alpha -1} (\mathbb{R},\mathbb{R}^{n} \times \mathbb{R}^{n} )$, where $\vartheta$ is a
solution to $\partial_{t} \vartheta = \xi$. Let us denote this function by $u=\Psi ( u_{0} , \xi , \vartheta \circ   \xi )$. The structure of the solution
which we derived above shows that for $\eta \in \CC^{2 \alpha -1}$, the \textsc{rde}
\begin{equation}\label{eq:rde-ren}
  \partial_{t} v=F ( v ) \xi +F' ( v ) F ( v ) \eta, \qquad v(0) = u_0,
\end{equation}
has a solution which can be obtained from the same map by replacing $\vartheta \circ   \xi$ with $\vartheta \circ   \xi + \eta$, i.e.
\[
   v= \Psi ( u_{0} , \xi , \vartheta \circ   \xi + \eta ).
\]
Moreover, the solution $v$ depends continuously on the data $\vartheta \circ \xi + \eta \in \CC^{2 \alpha -1}$. First note that this regularity hypothesis
concerns only the combination $\vartheta \circ   \xi + \eta$, and second note
that many equations share a similar structure and one can pass from solutions
of one equation to solutions of other equations via a transformation of the
(extended) data of the problem.



For many reasons, all these differential equations should be considered to be
the same object, especially when dealing with data of low regularity. To
understand this point of view, consider two different families $( \xi^{\varepsilon})_{\varepsilon >0}$ and $( \tilde{\xi}^\varepsilon )_{\varepsilon > 0}$ of smooth
functions, such that both converge to $\xi$ in $\CC^{\alpha -1}$. The
corresponding solutions $( u^\varepsilon )$ and $( \tilde{u}^\varepsilon)$ to
equation~{\eqref{eq:rde}} are given by $u^\varepsilon = \Psi ( u_{0} , \xi^\varepsilon ,\vartheta^\varepsilon \circ   \xi^\varepsilon )$ and $\tilde{u}^\varepsilon = \Psi ( u_{0} , \tilde{\xi}^\varepsilon,
\tilde{\vartheta}^\varepsilon \circ \tilde{\xi}^\varepsilon)$. Let us take the limit as $n \rightarrow \infty$ and assume that $\vartheta^\varepsilon \circ   \xi^\varepsilon \rightarrow \vartheta$
in $\CC^{2 \alpha -1}$ and also that $\tilde{\vartheta}^\varepsilon \circ   \tilde{\xi}^\varepsilon \rightarrow \tilde\vartheta$ in $\CC^{2 \alpha -1}$, where $\vartheta \neq \tilde\vartheta$. This could happen in principle, and it is not difficult to find specific and relevant examples of this multiplicity of limits. Of course, we have
\[
   \vartheta^\varepsilon \circ \xi^\varepsilon - \tilde{\vartheta}^\varepsilon \circ \tilde{\xi}^\varepsilon = \vartheta^\varepsilon \xi^\varepsilon - \tilde{\vartheta}^\varepsilon \tilde{\xi}^\varepsilon - \vartheta^\varepsilon \prec \xi^\varepsilon -  \vartheta^\varepsilon \succ \xi^\varepsilon + \tilde{\vartheta}^\varepsilon \prec \tilde{\xi}^\varepsilon +  \tilde{\vartheta}^\varepsilon \succ \tilde{\xi}^\varepsilon.
\]
Since $( \vartheta^\varepsilon )$ and $( \tilde{\vartheta}^\varepsilon )$ both converge to
$\vartheta$, where $\partial_{t} \vartheta = \xi$, we deduce from the
continuity of the paraproduct that all the terms on the right hand side cancel, except the first two, and we remain with
\[
   \eta = \vartheta - \vartheta' = \lim_{n \rightarrow \infty} ( \vartheta^\varepsilon \xi^\varepsilon - \tilde{\vartheta}^\varepsilon \tilde{\xi}^\varepsilon ).
\]
Incidentally, this line of reasoning also shows that if the limit exists, it
does not depend on the particular Littlewood-Paley decomposition we use to
compute it. From the continuity of $\Psi$ it moreover follows that $u^\varepsilon \rightarrow u$ and $\tilde{u}^\varepsilon \rightarrow \tilde{u}$ where $u= \Psi ( u_{0} , \xi , \vartheta )$ and $\tilde{u} = \Psi ( u_{0} , \xi , \vartheta + \eta )$. That is,
different approximations of the same equation could lead to different
equations in the limit. In particular, if $\xi$ is smooth enough (but
$(\xi^\varepsilon)$ does not converge to $\xi$ in a space of sufficiently high
regularity), we can interpret $u$ as a classical solution to a differential
equation, and in this case $u'$ will solve a modified equation.



\paragraph{Rough paths as a transformation group.}

We can therefore identify $\CC^{\alpha -1} \times \CC^{2 \alpha -1}$
with a transformation group $( T_{f,g} )_{f,g \in \CC^{\alpha -1}\times \CC^{2 \alpha -1}}$, which acts on solutions to ODEs via
\[
   T_{f,g} \Psi ( u_{0} , \xi , \eta ) = \Psi ( u_{0} , \xi +f, \eta + \vartheta \circ f+\Phi \circ   \xi +\Phi \circ  f+g ).
\]
where $\Phi$ solves $\partial_t \Phi = f$.
In particular,
$
   \Psi ( u_{0} , \xi , \eta ) =T_{\xi , \eta - \vartheta \circ   \xi} \Psi ( u_{0} ,0,0 ).
$
The neutral element is $T_{0,0}$ and the group operation is
$
   T_{f,g} T_{f' ,g'} =T_{f+f' ,g+g'} ,
$ so that the group is abelian. A simple distance is given by
\[
   d ( T_{f,g} ,T_{f' ,g'} ) =d ( T_{f-f' ,g-g'} ,T_{0,0} ) = \| f-f' \|_{\alpha -1} + \| g-g' + ( \Phi - \Phi' ) \circ ( f-f' ) \|_{2 \alpha -1},
\]
where $\Phi$ and $\Phi'$ are the definite integrals with $\Phi ( 0 ) = \Phi' (0 ) =0$ of $f$ and $f'$ respectively.




\paragraph{Geometric conditions.}
Note that if $J_{S}$ and $J_{A}$ denote respectively the projections onto
symmetric and antisymmetric tensors, then for smooth $\xi , \vartheta$ we have
\[
   J_{S} ( \vartheta \circ \xi ) = J_{S} ( \vartheta \circ \partial_{t}  \vartheta ) = \frac{1}{2} \partial_{t} J_{S} ( \vartheta \circ  \vartheta ),
\]
where the right hand side is now well defined for all $\vartheta \in \CC^{\alpha}$ and defines a distribution in $\CC^{2 \alpha -1}$. For any
family $\xi^\varepsilon \rightarrow \xi$ in $\CC^{\alpha -1}$ we have that
\[
   \vartheta^\varepsilon \circ \xi^\varepsilon = \frac{1}{2} \partial_{t} J_{S} (\vartheta^\varepsilon \circ   \vartheta^\varepsilon ) \rightarrow \frac{1}{2} \partial_{t} J_{S} ( \vartheta \circ   \vartheta )
\]
as $\varepsilon \rightarrow 0$, and only the antisymmetric part $J_{A} (\vartheta^\varepsilon \circ   \xi^\varepsilon )$ could feature multiple accumulation points.
To highlight the geometric meaning of the symmetric and antisymmetric parts of $\vartheta \circ \xi$ we can submit the \textsc{rde} to a  nontrivial transformation given by the application of a smooth
diffeomorphism $\phi \colon\mathbb{R}^{d} \rightarrow \mathbb{R}^{d}$. If $u$ is
a solution to the \textsc{rde} $\partial_{t} u=F ( u ) \xi$, then $v= \phi ( u )$
solves
\[
   \partial_{t} v= \phi' ( u ) \partial_{t} u= \phi' ( u ) F ( u ) \xi,
\]
i.e. an \textsc{rde} with vector field $( \phi' F ) \circ \phi^{-1}$. We have
\[
   \partial_{t} \phi ( T_{f,g} u ) = \phi' ( T_{f,g} u ) \partial_{t} (T_{f,g} u ) = \phi' ( T_{f,g} u ) ( F ( T_{f,g} u ) ( \xi +f ) +F' (T_{f,g} u ) F ( T_{f,g} u ) g ).
\]
On the other hand, we obtain for $T_{f,g} v=T_{f,g} \phi ( u )$ that
\[
   \partial_{t} T_{f,g} v= ( ( \phi' F ) \circ \phi^{-1} ) ( T_{f,g} v ) ( \xi+f ) + \mathD ( ( \phi' F ) \circ \phi^{-1} ) ( T_{f,g} v ) ( ( \phi' F ) \circ \phi^{-1} ) ( T_{f,g} v ) g,
\]
and we have
$
   \mathD ( ( \phi' F ) \circ \phi^{-1} ) = ( ( \phi'' F+ \phi' F' ) \circ \phi^{-1} ) ( \phi' )^{-1} \circ \phi^{-1} ,
$
where we slightly abuse notation by writing $( \cdummy )^{-1}$ both for the
inverse function of $\phi$ and for the inverse matrix of $\phi' ( x )$. Thus,
we obtain
\[
   \partial_{t} T_{f,g} v= ( ( \phi' F ) \circ \phi^{-1} ) ( T_{f,g} v ) ( \xi +f ) + ( ( \phi'' FF+ \phi' F' F ) \circ \phi^{-1} ) ( T_{f,g} v ) g,
\]
so that
$
   \phi ( T_{f,g} u ) =T_{f,g} ( \phi ( u ) )
$
only if
$
   ( \phi'' FF ) ( \phi^{-1} ( T_{f,g} v ) ) g=0.
$
If this holds for all $F$ and $\phi$, then the symmetric part $J_{S} g$ of
$g$ must be 0. In this sense, $T_{f,g}$ acts geometrically only if $J_S g=0$.




\subsection{Connections to rough paths and existence of the area}

We saw in the previous section that the solution $u$ to an \textsc{rde} of the form
$\partial_{t} u=F ( u ) \xi$ depends on the driving signal in a continuous way, provided
that we not only keep track of $\xi$ but also of $\vartheta \circ \xi$. From
the theory of rough paths it is well known that the same holds true if we keep
track of $\vartheta$ and its iterated integrals $\int_{s}^{t} \int_{s}^{r_{2}} \mathd \vartheta ( r_{1} ) \mathd \vartheta ( r_{2} )$. But in fact the convergence of $( \vartheta^{\varepsilon} \circ \xi^{\varepsilon} )$
is equivalent to the convergence of the iterated integrals $\int \int \mathd \vartheta^{\varepsilon} \mathd \vartheta^{\varepsilon}$:

\begin{corollary}\label{corollary:classical rough path}
  Let $( u^{\varepsilon} ,v^{\varepsilon} )_{\varepsilon >0}$ be a family of smooth functions on
  $\mathbb{R}$. Define for every $\varepsilon >0$ the ``area''
  \[
     A^{\varepsilon}_{s,t} = \int_{s}^{t} \int_{s}^{r_{2}} \mathd u^{\varepsilon} ( r_{1} ) \mathd v^{\varepsilon} ( r_{2} ) , \hspace{2em} s<t \in \mathbb{R}.
  \]
  Let $\alpha , \beta \in ( 0,1 )$ be such that $\alpha + \beta <1$ and let $u \in \CC^\alpha ,v \in \CC^\beta , \eta \in \CC^{\alpha + \beta -1}$. Then $( u^{\varepsilon} ,v^{\varepsilon} ,u^{\varepsilon} \circ \partial_{t} v^{\varepsilon} )$ converges to $( u,v, \eta )$ in $\CC^{\alpha} \times \CC^\beta \times \CC^{\alpha + \beta -1}$ if and only if $( u^{\varepsilon},v^{\varepsilon} )$ converges to $( u,v )$ in $\CC^{\alpha} \times \CC^\beta$, and if moreover
  \begin{equation}\label{eq:classical rough path}
     \lim_{\varepsilon \rightarrow 0} \left( \sup_{s \neq t \in \mathbb{R}, | s-t | \leqslant 1} \frac{| A_{s,t} -A^{\varepsilon}_{s,t} |}{| t-s |^{\alpha + \beta}} \right) =0,
  \end{equation}
  where we set $A_{s,t} = \int_{s}^{t} ( \eta + ( u  \prec \partial_{t} v ) + ( u \succ \partial_{t} v ) ) ( r ) \mathd r-u ( s ) ( v ( t ) -v ( s ) ) $ for $s,t \in \mathbb{R}$.
%  \begin{equation}\label{eq:classical area definition}
%     A_{s,t} = \int_{s}^{t} ( \eta + ( u  \prec \partial_{t} v ) + ( u \succ \partial_{t} v ) ) ( r ) \mathd r-u ( s ) ( v ( t ) -v ( s ) ) .
%  \end{equation}
\end{corollary}

\begin{proof}
  First suppose that $( u^{\varepsilon} ,v^{\varepsilon} ,u^{\varepsilon} \circ \partial_{t} v^{\varepsilon} )$ converges to $( u,v, \eta )$ in
  $\CC^{\alpha} \times \CC^\beta \times \CC^{\alpha + \beta -1}$, and let $s, t \in \mathbb{R}$ with $| s-t | \leqslant 1$.
  We have
  \begin{align}\label{eq:classical rough path pr1} \nonumber
     A_{s,t} -A^{\varepsilon}_{s,t} %= \int_{s}^{t} ( \eta +u  \prec  \partial_{t} v+u \succ \partial_{t} v ) ( r ) \mathd r-u ( s ) ( v ( t ) -v ( s ) ) - \int_{s}^{t} ( u^{\varepsilon} \circ \partial_{t} v^{\varepsilon} +u^{\varepsilon} \prec \partial_{t} v^{\varepsilon} +u^{\varepsilon} \succ \partial_{t} v^{\varepsilon} ) ( r ) \mathd r-u^{\varepsilon} ( s ) ( v^{\varepsilon} ( t ) -v^{\varepsilon} ( s ) )
     & = \int_{s}^{t} ( \eta +u \succ \partial_{t} v-u^{\varepsilon} \circ  \partial_{t} v^{\varepsilon} -u^{\varepsilon} \succ \partial_{t}  v^{\varepsilon} ) ( r ) \mathd r \\ \nonumber
     &\qquad + \int_{s}^{t} ( ( u^{\varepsilon} -u ) \prec \partial_{t} v^{\varepsilon} ) ( r ) \mathd r- ( u^{\varepsilon} -u ) ( s ) ( v^{\varepsilon} ( t ) -v^{\varepsilon} ( s ) ) \\
     &\qquad + \int_{s}^{t} ( u \prec \partial_{t} ( v^{\varepsilon} -v ) ) ( r ) \mathd r- ( u ) ( s ) ( ( v^{\varepsilon} -v ) ( t ) - ( v^{\varepsilon} -v ) ( s ) ) .
  \end{align}
  The first term on the right hand side can be estimated with the help of
  Lemma~\ref{lemma:integral}, which roughly states that time integral and paraproduct commute with each other, at the price of introducing a smoother remainder term. We get
  \begin{align*}
     &\left| \int_{s}^{t} ( \eta +u \succ \partial_{t} v-u^{\varepsilon} \circ  \partial_{t} v^{\varepsilon} -u^{\varepsilon} \succ \partial_{t} v^{\varepsilon} ) ( r ) \mathd r \right| \\
     &\hspace{60pt} \lesssim ( \| \eta -u^{\varepsilon} \circ \partial_{t} v^{\varepsilon} \|_{\alpha + \beta -1} + \| u-u^{\varepsilon} \|_{\alpha} \| \partial_{t}  v \|_{\beta -1} + \| u^{\varepsilon} \|_{\alpha} \| \partial_{t} (v^{\varepsilon} -v ) \|_{\beta -1} ) | t-s |^{\alpha + \beta} ,
  \end{align*}
  using also the paraproduct estimates. Since $\| \partial_{t} ( v^{\varepsilon} -v ) \|_{\beta -1} \lesssim \| v^{\varepsilon} -v \|_{\beta}$, the right hand side goes to zero if we divide it by $| t-s |^{\alpha + \beta}$ and let $\varepsilon \rightarrow 0$.
  
  The second term on the right hand side of~{\eqref{eq:classical rough path pr1}} can be estimated using Lemma~\ref{lemma:controlled old and new}:
  \[
     \left| \int_{s}^{t} ( ( u^{\varepsilon} -u ) \prec \partial_{t} v^{\varepsilon} ) ( r ) \mathd r- ( u^{\varepsilon} -u ) ( s ) ( v^{\varepsilon} ( t ) -v^{\varepsilon} ( s ) ) \right| \lesssim | t-s |^{\alpha + \beta} \| u^{\varepsilon} -u \|_{\alpha} \| v^{\varepsilon} \|_{\beta},
  \]
  The third term on the right hand side of~{\eqref{eq:classical rough path pr1}} is of the same type as the second term, and therefore the convergence
  in~{\eqref{eq:classical rough path}} follows.
  
  Conversely, assume that $( u^{\varepsilon} ,v^{\varepsilon} )$ converges to
  $( u,v )$ in $\CC^{\alpha} \times \CC^\beta$, and that the
  convergence in~{\eqref{eq:classical rough path}} holds. It follows from the
  representation~{\eqref{eq:classical rough path pr1}} and the
  convergence of $( u^{\varepsilon} ,v^{\varepsilon} )$ to $( u,v )$ in
  $\CC^{\alpha} \times \CC^\beta$, that also
  \[
     \lim_{\varepsilon \rightarrow 0} \Bigg( \sup_{s \neq t \in \mathbb{R}, | s-t | \leqslant 1} \frac{\big| \int_{s}^{t} ( \eta -u^{\varepsilon} \circ \partial_{r} v^{\varepsilon} ) ( r ) \mathd r \big|}{| t-s |^{\alpha + \beta}} \Bigg) =0.
  \]
  Due to the restriction $| s-t | \leqslant 1$, it is not entirely obvious
  that this implies the convergence of $u^{\varepsilon} \circ \partial_{r} v^{\varepsilon}$ to $\eta$ in $\CC^{\alpha + \beta -1}$. However, here
  we can use an alternative characterization of Besov spaces in terms of
  local means. Let $k^{0}$ and $k$ be infinitely differentiable functions on
  $\mathbb{R}$ with support in $( -1,1 )$, such that $\mathcal{F}k^{0} (0 ) \neq 0$, and such that there exists $\delta >0$ with $\mathcal{F}k ( z ) \neq 0$ for all $0< | z | < \delta$. Then an equivalent norm on $\CC^{\alpha + \beta -1} ( \mathbb{R} )$ is given by
  \[
     \| w \|_{\alpha + \beta -1} \simeq \max \Big\{ \| k^{0} \ast w \|_{L^{\infty}} , \sup_{j \geqslant 0}  2^{j ( \alpha + \beta -1 )} \| 2^{j} k ( 2^{j}  \cdummy ) \ast w \|_{L^{\infty}} \Big\} ,
  \]
  see~\cite{Triebel2006}, Theorem 1.10. Let us write $f= \int_{0}^{\cdummy} (\eta -u^{\varepsilon} \circ \partial_{r} v^{\varepsilon} ) ( r ) \mathd r$
  and let $t \in \mathbb{R}$ and $j \ge 0$. Then
  \begin{gather*}
     | 2^{j} k ( 2^{j} \cdummy ) \ast (\partial_t f) ( t ) | = 2^{2j} \left| \int_{\mathbb{R}} (\partial_{t} k) (2^j (t-s) ) ( f ( t ) -f ( s ) ) \mathd s \right| \\
     \lesssim 2^{2j} \int_{\mathbb{R}} | ( \partial_{t} k ) ( 2^{j} ( t-s ) ) |   | t-s |^{\alpha + \beta} \mathd s \sup_{| a-b | \leqslant 1} \frac{| f( b ) -f ( a ) |}{| b-a |^{\alpha + \beta}} \lesssim 2^{-j ( \alpha + \beta -1 )} \sup_{| a-b | \leqslant 1} \frac{| f ( b ) -f ( a ) |}{| b-a |^{\alpha + \beta}},
  \end{gather*}
  where we used that $\int_{\mathbb{R}} \partial_{t} k ( t ) \mathd t=0$,
  and that $k$ is supported in $( -1,1)$. Similarly, we obtain
  \[
     | k^{0} \ast ( \partial_{t} f ) ( t ) |  \lesssim \int_{\mathbb{R}} | \partial_{t} k^{0} ( t-s ) |   | t-s |^{\alpha + \beta} \mathd s \sup_{| a-b | \leqslant 1} \frac{| f ( b ) -f ( a ) |}{| b-a |^{\alpha + \beta}} \lesssim \sup_{| a-b | \leqslant 1} \frac{| f ( b ) -f ( a ) |}{| b-a |^{\alpha + \beta}},
  \]
  from where the convergence of $u^{\varepsilon} \circ \partial_{t} v^{\varepsilon}$ to $\eta$ in $\CC^{\alpha + \beta -1}$ follows.
\end{proof}

\begin{corollary}\label{corollary:existence of rough path area}
  Let $X$ be an $n$--dimensional centered Gaussian process with independent components, whose covariance function satisfies for some $H \in ( 1/4,1 )$ the inequalities
  \begin{gather}\label{eq:classical rough path covariance} \nonumber
     \mathbb{E} [ | X_{t} -X_{s} |^{2} ] \lesssim | t-s |^{2H} \hspace{2em} \tmop{and} \\
     | \mathbb{E} [ ( X_{s+r} -X_{s} ) ( X_{t+r} -X_{t} ) ] | \lesssim | t-s |^{2H-2}  r^{2}
  \end{gather}
  for all $s,t \in \mathbb{R}$ and all $r \in [ 0, | t-s | )$. Let $\varphi \in \CD$. Then $\varphi X \in \CC^{\alpha}$ for every $\alpha <H$, and there
  exists $\eta \in \CC^{2 \alpha -1}$ such that for every $\psi \in \CS$ with
  $\int \psi \mathd t=1$ and for every $\delta >0$ we have
  \[
     \lim_{\varepsilon \rightarrow 0} \mathbb{P} \left( \| \psi^{\varepsilon}  \ast ( \varphi X ) - ( \varphi X ) \|_{\alpha} + \| ( \psi^{\varepsilon}  \ast ( \varphi X ) ) \circ \partial_{t} ( \psi^{\varepsilon} \ast (\varphi X ) ) - \eta \|_{\CC^{2 \alpha -1}} > \delta \right) =0,
  \]
  where we define $\psi^{\varepsilon} = \varepsilon^{-1} \psi (\varepsilon^{-1} \cdummy )$.
\end{corollary}

\begin{proof}
  Since $\varphi$ is smooth and of compact support, it is easy to see that
  also the Gaussian process $\varphi X$ satisfies the covariance
  condition~{\eqref{eq:classical rough path covariance}}. Moreover, $\varphi X$ has compact support. So by Theorem 15.45 of~{\cite{Friz2010}}, for every $p \geqslant 1$, the
  iterated integrals $\int_{s}^{t} \int_{s}^{r_{2}} \mathd \psi^{\varepsilon} \ast ( \varphi X ) ( r_{1} ) \mathd \psi^{\varepsilon} \ast ( \varphi X ) ( r_{2} )$ converge in $L^{p}$ in the sense of~{\eqref{eq:classical rough path}}. The statement then follows from Corollary~\ref{corollary:classical rough path}.
\end{proof}

\begin{remark}
  The proof of Corollary~\ref{corollary:classical rough path} actually shows
  more than the equivalence of the convergence of $A^{\varepsilon}$ and of
  $u^{\varepsilon} \circ \partial_{t} v^{\varepsilon}$: it shows that the norm
  of $( u^{\varepsilon} \circ \partial_{t} v^{\varepsilon} - \eta )$ can be
  controlled by a polynomial of the norms of $( A^{\varepsilon} -A )$, $(u^{\varepsilon} -u )$, and $(v^{\varepsilon} -v)$. So in fact we
  have $L^{p}$--convergence in Corollary~\ref{corollary:existence of rough path area}, and not just convergence in probability. Alternatively, the $L^p$--convergence is obtained from the convergence in probability because we are considering elements living in a fixed Gaussian chaos, see Theorem 3.50 of~\cite{Janson1997}.
\end{remark}

Combining Corollary~\ref{corollary:existence of rough path area} with Theorem~\ref{theorem:rde existence uniqueness}, we obtain the following corollary:

\begin{corollary}
  Let $X$ be a $n$--dimensional centered Gaussian process satisfying the conditions of Corollary~\ref{corollary:classical rough path} for some $H > 1/3$, and let $\varphi \in \CD$ and $F \in C^3_b$. Then there exists a unique solution $u$ to
  \[
     \partial_t u = F(u) \partial_t (\varphi X), \qquad u(0) = u_0,
  \]
  in the following sense: If $\psi \in \CS$ with $\int \psi \mathd t=1$ and if for $\varepsilon > 0$ the function $u^\varepsilon$ solves
  \[
     \partial_t u^\varepsilon = F(u^\varepsilon) \partial_t (\varphi X)^\varepsilon, \qquad u(0) = u_0,
  \]
  where $(\varphi X)^\varepsilon = \varepsilon^{-1} \psi(\varepsilon \cdot) \ast (\varphi X)$, then $\lim_{\varepsilon \rightarrow 0} \E[\| u - u^\varepsilon \|_{\alpha}^p] = 0$ for all $p \ge 1$ and all $\alpha < H$.
\end{corollary}


\section{Rough Burgers equation}\label{sec:burgers}

Fix now $\sigma > 5/6$ and consider the following PDE on $[ 0,T ] \times \mathbb{T}$ for some fixed $T>0$:
\begin{equation}\label{eq:burgers-rough}
  L u=G ( u ) \partial_{x} u+ \xi, \qquad u(0) = u_0,
\end{equation}
where $L= \partial_{t} + ( - \Delta )^{\sigma}$ and $u_0 \in \CC^\alpha$ for a suitable $\alpha$. We would like to consider solutions $u$ in the case of a
distributional $\xi$, and in particular we want to allow $\xi$ to be a typical
realization of a space-time white noise. We will see below that in this case
the solution $\vartheta$ to the linear equation $L  \vartheta = \xi$,
$\vartheta ( 0 ) =0$, belongs (locally in time) to $\CC^{\alpha} ( \mathbb{T})$ for any $\alpha < \sigma -1/2$, but it is not better than that. This is
also the regularity to be expected from the solution $u$ of the non-linear
problem~{\eqref{eq:burgers-rough}}, and therefore the term $G ( u ( t ) ) \partial_{x} u ( t )$ is not well defined since $G ( u ( t ) ) \in \CC^{\alpha} ( \mathbb{T} )$ and $\partial_{x} u ( t ) \in \CC^{\alpha -1} (\mathbb{T} )$, and the sum of their regularities fails to be positive.

For $\sigma =1$, this equation has been solved by
Hairer~{\cite{Hairer2011}}, who used rough path integrals to define the
product $G ( u ) \partial_{x} u$. In the following, we will show how to solve
the equation using our paracontrolled calculus.



While in general it is possible to set up this equation in a space-time Besov
space, the fact that the distribution $\xi$ (which is a genuine space-time
distribution) enters the problem linearly allows for a small simplification.
Indeed, if we let $w=u- \vartheta$, then $w$ solves the PDE
\begin{equation}\label{eq:burgers-alt}
  L w=G ( \vartheta +w ) \partial_{x} ( \vartheta +w ),
\end{equation}
which can be studied as an evolution equation for a continuous function of
time with values in a suitable H{\"o}lder-Besov space:

Recall that for $T>0$
and $\beta \in \mathbb{R}$ we defined the spaces $C_{T} \CC^{\beta} =C( [ 0,T ], \CC^{\beta} ( \mathbb{T}^{d},\mathbb{R}^{n} ))$
with norm $\| u \|_{C_{T} \CC^{\beta}} = \sup_{0 \leqslant s \leqslant T} \| u ( s ) \|_{\beta}$. By the regularity theory for $L$ we expect $w \in C_{T} \CC^{\alpha -1+2 \sigma}$ whenever $G ( \vartheta +w ) \partial_{x} ( \vartheta +w ) \in C_{T}\CC^{\alpha -1}$ (at least in the sense of uniform estimates as the
regularization goes to zero). The paraproduct allows us to decompose the right
hand side of~{\eqref{eq:burgers-alt}} as
\[
   G ( \vartheta +w ) \partial_{x} ( \vartheta +w ) =G ( \vartheta +w ) \prec \partial_{x} \vartheta +G ( \vartheta +w ) \circ \partial_{x} \vartheta +G( \vartheta +w ) \succ \partial_{x} \vartheta +G ( \vartheta +w ) \partial_{x} w,
\]
where we have expanded only the term containing $\partial_{x} \vartheta$ since
the one linear in $\partial_{x} w$ is well defined under the hypothesis that
$w \in C_{T} \CC^{\alpha -1+2 \sigma}$. Note that here we only let the
paraproduct act on the spatial variables, i.e. $G ( \vartheta +w ) \prec \partial_{x} \vartheta$ should really be understood as $t \mapsto G (\vartheta ( t ) +w ( t ) ) \prec \partial_{x} \vartheta ( t )$, an element of
$C_{T} \CC^{\alpha -1}$. A simple modification of the proof of
Lemma~\ref{lemma:paralinearization} shows that, for $\alpha \in (0,1/2)$, we have
\[
   \| G ( \vartheta +w ) -G' ( \vartheta +w ) \prec \vartheta \|_{2 \alpha}  \lesssim \| G \|_{C^{2}_{b}} ( 1+ \| \vartheta \|^{2}_{\alpha} + \| w \|_{2 \alpha} ) \lesssim \| G \|_{C^{2}_{b}} ( 1+ \| \vartheta \|^{2}_{\alpha} +  \| w \|_{\alpha -1+2 \sigma} ) ,
\]
where we used that $\alpha -1+2 \sigma >2 \alpha$, which holds because $\alpha < \sigma -1/2<2 \sigma -1$. The linear dependence on the norm of $w$ will be
crucial for obtaining global solutions. We can now rewrite
\[
   G ( \vartheta +w ) \circ \partial_{x} \vartheta = ( G ( \vartheta +w ) -G'( \vartheta +w ) \prec \vartheta ) \circ \partial_{x} \vartheta +C ( G' ( \vartheta +w ) , \vartheta , \partial_{x} \vartheta ) +G' ( \vartheta +w )( \vartheta \circ \partial_{x} \vartheta ).
\]
So if we assume that $( \vartheta \circ \partial_{x} \vartheta ) \in C_{T} \CC^{2 \alpha -1}$, then we have a well behaved representation of the resonant
term $G ( \vartheta +w ) \circ \partial_{x} \vartheta$, and
\begin{align}\label{eq:burgers resonant estimate} \nonumber
   \| G ( \vartheta +w ) \circ \partial_{x} \vartheta \|_{2 \alpha -1} & \lesssim \| G \|_{C^{2}_{b}} ( 1+ \| \vartheta \|^{2}_{\alpha} + \| w \|_{\alpha -1+2 \sigma} ) \| \partial_{x} \vartheta \|_{\alpha -1} + \| G'( \vartheta +w ) \|_{\alpha} \| \vartheta \|_{\alpha}^{2} \\
   &\qquad + \| G' ( \vartheta +w ) \|_{\alpha} \| \vartheta \circ \partial_{x} \vartheta \|_{2 \alpha -1} \lesssim C_{G} C_{\vartheta} ( 1+ \| w \|_{\alpha -1+2 \sigma} ) ,
\end{align}
where we set
\[
   C_{G} = \| G \|_{C^{2}_{b}} \hspace{2em} \tmop{and} \hspace{2em}  C_{\vartheta} = \sup_{t \in [ 0,T ]} ( 1+ \| \vartheta ( t ) \|_{\alpha} )^{3} ( 1+ \| \vartheta ( t ) \circ \partial_{x} \vartheta ( t ) \|_{2 \alpha -1} ) .
\]
Let us now define
\[
   \Phi =G ( \vartheta +w ) \prec \partial_{x} \vartheta +G ( \vartheta +w ) \succ \partial_{x} \vartheta +G ( \vartheta +w ) \circ \partial_{x} \vartheta ,
\]
so that~{\eqref{eq:burgers resonant estimate}} and the paraproduct estimates yield
\begin{equation}\label{eq:burgers phi-estimate}
   \| \Phi \|_{\alpha -1} \lesssim C_{G}C_{\vartheta} ( 1+ \| w \|_{\alpha -1+2 \sigma} ) ,
\end{equation}
and $w$ satisfies $L w= \Phi +G ( \vartheta +w ) \partial_{x} w$. So if we denote by $( P_{t} )_{t \geqslant 0}$ the semigroup generated by $- (- \Delta )^{\sigma}$, then
\begin{equation}\label{eq:mild-burgers}
  w ( t ) =P_{t} u_{0} + \int_{0}^{t} P_{t-s} \Phi ( s ) \mathd s+ \int_{0}^{t} P_{t-s} ( G ( \vartheta ( s ) +w ( s ) ) \partial_{x} w ( s ) ) \mathd s,
\end{equation}
where we assumed that $\vartheta(0) = 0$. Applying the Schauder estimates for the fractional Laplacian (Lemma~\ref{lemma:schauder} and Lemma~\ref{lemma:heat flow smoothing}) to~{\eqref{eq:mild-burgers}}, we obtain for all $t>0$ that
\begin{align*}
   \| w ( t ) \|_{\alpha -1+2 \sigma} & = \left\| P_{t} u_{0} + \int_{0}^{t} P_{t-s} \Phi ( s ) \mathd s+ \int_{0}^{t} P_{t-s} ( G ( \vartheta ( s ) +w ( s ) ) \partial_{x} w ( s ) ) \mathd s \right\|_{\alpha -1+2 \sigma} \\
   & \leqslant c t^{- ( 2 \sigma -1 ) /2 \sigma} \Big( \| u_{0} \|_{\alpha} + \sup_{s \in [ 0,t ]} ( s^{( 2 \sigma -1 ) / ( 2 \sigma )} \| \Phi ( s )\|_{\alpha -1} ) \Big) \\
   &\qquad +c \int_{0}^{t} \frac{\| G ( \vartheta ( s ) +w ( s ) ) \partial_{x} w ( s ) \|_{L^{\infty}}}{( t-s )^{( \alpha -1+2 \sigma ) / ( 2 \sigma )}} \mathd s,
\end{align*}
where $c>0$ is a generic constant whose value may change in every step. But
now recall from {\eqref{eq:burgers phi-estimate}} that $\| \Phi ( s )\|_{\alpha -1} \lesssim C_{G} C_{\vartheta} ( 1+ \| w ( s ) \|_{\alpha -1+2 \sigma} )$. Moreover, if we choose $\alpha \in ( 1/3, \sigma -1/2 )$ close enough to $\sigma -1/2$, then $\alpha +2 \sigma -2>0$ (recall that $\sigma >5/6$), and therefore
\[
   \| G ( \vartheta ( s ) +w ( s ) ) \partial_{x} w ( s ) \|_{L^{\infty}} \lesssim \| G \|_{L^{\infty}} \| \partial_{x} w ( s ) \|_{\alpha -2+2 \sigma} \lesssim \| G \|_{L^{\infty}} \| w ( s ) \|_{\alpha -1+2 \sigma} .
\]
Thus, we get for all $t \in [ 0,T ]$ that
\begin{align*}
   ( t^{1-1/ ( 2 \sigma )} \| w ( t ) \|_{\alpha -1+2 \sigma} ) & \leqslant c \| u_{0} \|_{\alpha} +c C_{\vartheta} C_{G} ( 1+ \sup_{s \in [ 0,t ]} ( s^{( 2 \sigma -1 ) / ( 2 \sigma )} \| w ( s ) \|_{\alpha -1+2 \sigma} ) ) \\
   &\quad +c C_{G} t^{1-1/ ( 2 \sigma )} \int_{0}^{t} ( s^{1-1/ ( 2 \sigma )} \| w ( s ) \|_{\alpha +1} ) \frac{\mathd s}{( t-s )^{( \alpha -1+2 \sigma ) / ( 2 \sigma )} s^{1-1/ ( 2 \sigma )}}
\end{align*}
Since $( \alpha -1+2 \sigma ) / ( 2 \sigma ) <1$, we have
\[
   t^{1-1/ ( 2 \sigma )} \int_{0}^{t} \frac{\mathd s}{( t-s )^{( \alpha -1+2 \sigma ) / ( 2 \sigma )} s^{1-1/ ( 2 \sigma )}} \lesssim t^{1- ( \alpha -1+2 \sigma ) / ( 2 \sigma )} \lesssim 1
\]
for $t \in [ 0,T ]$. Putting everything together, we conclude that
\[
   ( t^{1-1/ ( 2 \sigma )} \| w ( t ) \|_{\alpha -1+2 \sigma} ) \leqslant c \| u_{0} \|_{\alpha} +c C_{\vartheta} C_{G} ( 1+ \sup_{s \in [ 0,t ]} ( s^{( 2 \sigma -1 ) / ( 2 \sigma )} \| w ( s ) \|_{\alpha -1+2 \sigma} ) ) .
\]

In order to turn this into a bound on $\| w \|_{C_{T} \CC^{\alpha -1+2 \sigma}}$, we use again a scaling argument. We extend the scaling
transformation to the time variable in such a way that it leaves the operator
$L$ invariant. More precisely, for $\lambda >0$ we set $\Lambda_{\lambda} u (t,x ) =u ( \lambda^{2 \sigma} t, \lambda x )$, so that $L  \Lambda_{\lambda} =\lambda^{2 \sigma} \Lambda_{\lambda} L$. Now let $u^{\lambda} =\Lambda_{\lambda} u$, $w^{\lambda} = \Lambda_{\lambda} w$, and
$\vartheta^{\lambda} = \Lambda_{\lambda} \vartheta$. Note that $u^{\lambda} \colon [ 0,T/ \lambda^{2 \sigma} ] \times \mathbb{T}_{\lambda} \rightarrow \mathbb{R}$, where $\mathbb{T}_{\lambda} = \R / (2\pi\lambda^{-1}\Z)$ is a
rescaled torus, and that $w^{\lambda}$ solves the equation
\[
   L w^{\lambda} = \lambda^{2 \sigma} \Lambda_{\lambda} L w= \lambda^{2 \sigma} \Lambda_{\lambda} ( \Phi +G ( w+ \vartheta ) \partial_{x} w ) = \lambda^{2 \sigma} \Lambda_{\lambda} \Phi + \lambda^{2 \sigma -1} G ( w^{\lambda} + \vartheta^{\lambda} ) \partial_{x} w^{\lambda}.
\]
The same derivation as above shows that
\[
   \| \Lambda_{\lambda} \Phi ( t ) \|_{\alpha -1} \lesssim C_{G} C_{\vartheta^{\lambda}} ( 1+ \| w^{\lambda} ( t ) \|_{\alpha -1+2 \sigma} ),
\]
where
\[
   C_{\vartheta^{\lambda}} = \sup_{t \in [ 0,T ]} ( 1+ \| \vartheta^{\lambda} ( t ) \|_{\alpha} )^{3} ( 1+ \| \Lambda_{\lambda} ( \vartheta ( t ) \circ \partial_{x} \vartheta ( t ) ) \|_{2 \alpha -1} ) \lesssim \lambda^{2 \alpha -1} C_{\vartheta} \leqslant \lambda^{-1} C_{\vartheta}
\]
as long as $\lambda \in ( 0,1 ]$. Thus, we finally conclude that
\begin{align*}
   ( t^{1-1/ ( 2 \sigma )} \| w^{\lambda} ( t ) \|_{\alpha -1+2 \sigma} ) & \leqslant c \| \Lambda_{\lambda} u_{0} \|_{\alpha} + \lambda^{2 \sigma -1} c C_{\vartheta} C_{G} ( 1+ \sup_{s \in [ 0,t ]} ( s^{( 2 \sigma -1 ) / ( 2\sigma )} \| w^{\lambda} ( s ) \|_{\alpha -1+2 \sigma} ) ) \\
   & \leqslant c \| u_{0} \|_{\alpha} + \lambda^{2 \sigma -1} c C_{\vartheta} C_{G} ( 1+ \sup_{s \in [ 0,t ]} ( s^{( 2 \sigma -1 ) / ( 2 \sigma )} \| w^{\lambda} ( s ) \|_{\alpha -1+2 \sigma} ) )
\end{align*}
for all $\lambda \in ( 0,1 ]$. Since $2 \sigma -1>0$, we get for small enough
$\lambda >0$, depending only on $C_{\vartheta}$ and $C_{G}$ but not on
$u_{0}$, that
\[
   \sup_{t \in [ 0,T ]} ( t^{1-1/ ( 2 \sigma )} \| w^{\lambda} ( t ) \|_{\alpha -1+2 \sigma} ) \leqslant 2 ( c \| u_{0} \|_{\alpha} +1 ).
\]
But $u= \Lambda_{\lambda^{-1}} ( w^{\lambda} + \vartheta^{\lambda} )$, and therefore
\[
   \sup_{t \in [ 0, \lambda^{2 \sigma} T ]} \| u ( t ) \|_{\alpha} \lesssim_{\lambda} \| u_{0} \| +C_{\vartheta}.
\]
This provides the key ingredient for obtaining a uniform estimate on the full time interval $[ 0,T ]$, and then the existence of global solutions to the
Burgers equation.

Uniqueness in the space of solutions $u$ with decomposition $u= \vartheta +w$
with $w \in C_{T} \CC^{\alpha -1+2 \sigma}$ can be handled easily along the
lines above, and we obtain the following result:

\begin{theorem}\label{theorem:burgers existence uniqueness}
  Let $\sigma > 5/6$, $\alpha \in ( 1/3, \sigma -1/2 )$, let $T>0$,
  and assume that $( \xi^{\varepsilon} )_{\varepsilon >0}$ is a family of
  smooth functions on $[ 0,T ] \times \mathbb{T}$ with values in
  $\mathbb{R}^{n}$, and $G \in C^{3}_{b} ( \mathbb{R}^{n} ,\Lcal (\mathbb{R}^{n} ,\mathbb{R}^{n} ) )$. Suppose that there exist $\vartheta \in C_{T} \CC^{\alpha}$ and $\eta \in C_{T} \CC^{2 \alpha -1}$ such that $(\vartheta^{\varepsilon} , ( \vartheta^{\varepsilon} \circ   \partial_{x} \vartheta^{\varepsilon} ) )$ converges to $( \vartheta , \eta )$ in $C_{T}
  \CC^{\alpha -1} \times C_{T} \CC^{2 \alpha -1}$, where
  $\vartheta^{\varepsilon}$ are solutions to $L \vartheta^{\varepsilon} = \xi^{\varepsilon}$ and $\vartheta^{\varepsilon} ( 0 ) =0$, and where $L= \partial_{t} + ( - \Delta )^{\sigma}$. Let for $\varepsilon >0$ the function $u^{\varepsilon}$ be the unique global solution to the Cauchy problem
  \[
     L u^{\varepsilon} =G ( u^{\varepsilon} ) \partial_{x} u^{\varepsilon} + \xi^{\varepsilon} , \hspace{2em} u^{\varepsilon} ( 0 ) =u_{0} ,
  \]
  where $u_{0} \in \CC^{\alpha}$. Then there exists $u \in C_{T} \CC^{\alpha}$ such that $u^{\varepsilon} \rightarrow u$ in $C_{T} \CC^\alpha$. The limit $u$ depends only on $( u_{0} , \vartheta ,\eta )$, and not on the approximating family $( \vartheta^{\varepsilon} , (\vartheta^{\varepsilon} \circ   \partial_{x} \vartheta^{\varepsilon} ) )$.
\end{theorem}

\begin{remark}
  As in the case of \textsc{rde}s, it is not important that $\vartheta^{\varepsilon} ( 0 ) =0$, and we can take any other initial condition as long as it is
  sufficiently regular: If $L \tilde{\vartheta}^{\varepsilon} = \xi^{\varepsilon}$ and $\tilde{\vartheta}^{\varepsilon} ( 0 ) =f \in \CC^{2 \alpha}$, then $L ( \tilde{\vartheta}^{\varepsilon} -  \vartheta^{\varepsilon} ) =0$ and $( \tilde{\vartheta}^{\varepsilon} - \vartheta^{\varepsilon} ) ( 0 ) =f$, so that $(\tilde{\vartheta}^{\varepsilon} - \vartheta^{\varepsilon} ) ( t ) =P_{t} f$.  Since $\sup_{t \in [ 0,T ]} \| P_{t} f \|_{2 \alpha} \lesssim \| f \|_{2 \alpha}$, we conclude, as in Remark~\ref{remark:ode initial condition theta}, that $\vartheta^{\varepsilon} \circ \partial_{x}  \vartheta^{\varepsilon}$ converges to $\eta$ in $C_{T} \CC^{2 \alpha -1}$ if and only $\tilde{\vartheta}^{\varepsilon} \circ \partial_{x} \tilde{\vartheta}^{\varepsilon}$ converges to $\eta + ( P_{\cdummy} f ) \circ \partial_{x} \tilde{\vartheta}$ in $C_{T} \CC^{2 \alpha -1}$.
\end{remark}

\begin{remark}
  As for \textsc{rde}s, the limit $u$ of the regularized solutions $u^{\varepsilon}$ actually solves the equation
  \[
     L u=G ( u ) \partial_{x} u+ \xi , \hspace{2em} u ( 0 ) =u_{0}
  \]
  in the weak sense, as long as we interpret the product $G ( u ) \partial_{x} u$ correctly, using our product operator $\pi_{G}$ defined in
  Theorem~\ref{theorem:product operator}.
\end{remark}

\begin{remark}
   Of course the solution $u$ to the fractional Burgers type equation also depends continuously on the initial condition $u_0$.
\end{remark}






\subsection{Construction of the area}

%It remains to show that if $\xi$ is a space time white noise, then the solution $\vartheta$ to
%\begin{align}\label{eq:burgers linear}
%   L \vartheta = \xi, \qquad \vartheta(0) = 0,
%\end{align}
%is in $C_T\CC^\alpha$ for every $\alpha < \sigma - 1/2$, and that the area $\vartheta \circ \partial_x \vartheta$ is in $C_T \CC^{2 \alpha - 1}$. Some general results on the existence of the area for Gaussian processes indexed by a one-dimensional spatial variable are shown in~\cite{Friz2012}. However, in the present setting it is relatively straightforward to construct the area ``by hand'', using Fourier analytic methods.
%
%Recall that $\xi$ is a centered Gaussian space-time distribution, such that
%\[
%   \E[\langle \xi^i, \varphi \rangle \langle \xi^{i'}, \psi \rangle] = \mathbf{1}_{i=j} \int_{\R \times\mathbb{T}} \varphi(t,x) \psi (t,x) \mathd t \mathd x
%\]
%for all test function $\varphi, \psi \in L^2([0,T] \times \mathbb{T})$ and for all $i, i'  \in \{ 1, \dots, n\}$. Formally we write $\E[\xi^i(t,x) \xi^{i'}(t',x')] = \mathbf{1}_{i = i'} \delta(t-t') \delta(x-x')$, where $\delta$ denotes the Dirac delta. Let $\CF$ denote the spatial Fourier transform, i.e. $\CF u(t,\cdot) (k) = \int_\mathbb{T} e^{-\imath k x} u(t,x) \mathd x$. A simple calculation shows that $\CF \xi$ is a complex valued, centered Gaussian space-time distribution, whose covariance is formally given by
%\[
%   \E[\CF \xi^i (t,\cdot)(k) \overline{\CF \xi^{i'} (t',\cdot)(k')}] = 2\pi \mathbf{1}_{i=i'} \mathbf{1}_{k=k'} \delta(t - t')
%\]
%for $i,i'\in \{1, \dots, n\}$, $t, t' \in [0,T]$, and $k, k' \in \Z$. If $(P_t)_{t\ge 0} = (e^{-t |\cdot|^{2\sigma}}(\mathD))_{t \ge 0}$ denotes the semigroup generated by $-(-\Delta)^\sigma$, then the solution $\vartheta$ to~\eqref{eq:burgers linear} is given by $\vartheta(t,x) = \int_{0}^t (P_{t-s} \xi)(x)\mathd s$, $t \in [0,T]$, from where another straightforward calculation yields the following result:

It remains to show that if $\xi$ is a space-time white noise, then the solution $\vartheta$ to $L \vartheta = \xi$, $\qquad \vartheta(0) = 0$, is in $C_T\CC^\alpha$ for every $\alpha < \sigma - 1/2$, and that the area $\vartheta \circ \partial_x \vartheta$ is in $C_T \CC^{2 \alpha - 1}$. Some general results on the existence of the area for Gaussian processes indexed by a one-dimensional spatial variable are shown in~\cite{Friz2012}. However, in the present setting it is relatively straightforward to construct the area ``by hand'', using Fourier analytic methods.
%
%Recall that $\xi$ is a centered Gaussian space-time distribution, whose covariance is formally given as $\E[\xi^i(t,x) \xi^{i'}(t',x')] = \mathbf{1}_{i = i'} \delta(t-t') \delta(x-x')$, where $\delta$ denotes the Dirac delta. 

In this section, we use $\CF$ to denote the spatial Fourier transform, i.e. $\CF u(t,\cdot) (k) = \int_\mathbb{T} e^{-\imath k x} u(t,x) \mathd x$. Recall that $\CF \xi$ is a complex valued, centered Gaussian space-time distribution, whose covariance is formally given by
\[
   \E[\CF \xi^i (t,\cdot)(k) \overline{\CF \xi^{i'} (t',\cdot)(k')}] = 2\pi \mathbf{1}_{i=i'} \mathbf{1}_{k=k'} \delta(t - t')
\]
for $i,i'\in \{1, \dots, n\}$, $t, t' \in [0,T]$, $k, k' \in \Z$, where $\delta$ denotes the Dirac delta. If $(P_t)_{t\ge 0} = (e^{-t |\cdot|^{2\sigma}}(\mathD))_{t \ge 0}$ denotes the semigroup generated by $-(-\Delta)^\sigma$, then $\vartheta(t,x) = \int_{0}^t (P_{t-s} \xi)(x)\mathd s$, $t \in [0,T]$, from where a straightforward calculation yields the following result:


\begin{lemma}\label{l:heat covariance}
   The spatial Fourier transform $\CF \vartheta$ of $\vartheta$ is a complex-valued Gaussian process with zero mean and covariance
   \begin{equation*}%\label{e:stationary covariance}
      \E[\CF{\vartheta}^i(t,\cdot)(k) \overline{\CF{\vartheta}^{i'}(t',\cdot)(k')}] = \begin{cases}
                                                                                                       2 \pi \mathbf{1}_{i=i'} \mathbf{1}_{k=k'} (e^{-|t'-t| | k |^{2 \sigma}} - e^{-(t+t')|k|^{2\sigma} })/(2|k|^{2 \sigma}), & k \neq 0,\\
                                                                                                       2\pi \mathbf{1}_{i=i'} \mathbf{1}_{k=k'} t \wedge t', & k = 0,
                                                                                                   \end{cases}
   \end{equation*}
   for $i,i' \in \{1, \dots, n\}$, $k, k' \in \mathbb{Z}$, and $t,t' \in [0,T]$. As a consequence, $\E[\CF{\vartheta}^i_{s,t}(0) \overline{\CF{\vartheta}^{i}_{s,t}(k')}] = 2\pi \mathbf{1}_{i=i'} \mathbf{1}_{k'=0} |t-s|$, and for $k\neq 0$
   \begin{equation*}
      \E[\CF{\vartheta}^i_{s,t}(k) \overline{\CF{\vartheta}^{i'}_{s,t}(k')}] = \pi \mathbf{1}_{i=i'} \mathbf{1}_{k=k'} (2 - e^{-2 s | k |^{2 \sigma}}  - e^{-2 t | k |^{2 \sigma}} - 2 e^{- 2 |t-s| |k|^{2\sigma}} + 2 e^{-(s+t)|k|^{2\sigma}})/|k|^{2 \sigma},
   \end{equation*}
   where we write $\CF{\vartheta}^i_{s,t}(k) = \CF{\vartheta}^i(t,\cdot)(k) - \CF{\vartheta}^i(s,\cdot)(k)$ for all $0 \leqslant s < t \leqslant T$. In particular,
   \begin{equation}
      |\E[\CF{\vartheta}^i_{s,t}(k) \overline{\CF{\vartheta}^{i}_{s,t}(k)}]| \lesssim |t-s|^{\delta} |k|^{-2\sigma(1 - \varepsilon)}
   \end{equation}
   for all $\delta \in [0,1]$ and all $k \neq 0$.
\end{lemma}

%\begin{proof}
%   We give a formal derivation, which can be rendered rigorous by considering the action on test functions. For the spatial Fourier transform $\widehat{\dot{W}}$ of $\dot{W}$ we obtain
%   \begin{align*}
%      E\Bigl(\widehat{\dot{W}^j}(t,k)\widehat{\dot{W}^\ell}(t',k')\Bigr) & = \int_{\mathbb{T}^d}\int_{\mathbb{T}^d} e^{-\imath \langle k,x\rangle} e^{-\imath \langle k',x'\rangle} E(\dot{W}^j(t,x)\dot{W}^\ell(t',x')) \mathd x' \mathd x \\
%      & = \delta_{j\ell} \delta(t-t') \int_{\mathbb{T}^d} e^{-\imath \langle k,x\rangle} e^{-\imath \langle k',x\rangle} \mathd x \\
%      & = (2\pi)^d \delta_{j\ell} \delta(t-t') \delta_{k(-k')}.
%   \end{align*}
%   Hence, the covariance of $\widehat{\vartheta}$ is given by
%   \begin{align*}
%      &E(\widehat{\vartheta}^j(t,k) \widehat{\vartheta}^\ell(t',k' ))\\
%      &\hspace{40pt} = \int_{-\infty}^t \int_{-\infty}^{t'} e^{-(t-s)(\lambda + |k|^{2\sigma})} e^{-(t'-s')(\lambda + |k'|^{2\sigma})}E\Bigl(\widehat{\dot{W}^j}(s,k)\widehat{\dot{W}^\ell}(s',k')\Bigr) \mathd s' \mathd s \\
%      &\hspace{40pt} = (2\pi)^d \delta_{j\ell} \delta_{k(-k')} \int_{-\infty}^{t \wedge t'} e^{-(t-s)(\lambda + |k|^{2\sigma})} e^{-(t'-s)(\lambda + |k|^{2\sigma})} \mathd s \\
%      &\hspace{40pt} = (2\pi)^d \delta_{j\ell} \delta_{k(-k')} \frac{e^{2(t\wedge t')(\lambda + |k|^{2\sigma})}}{2(\lambda + |k|^{2\sigma})} e^{-(t+t')(\lambda + |k|^{2\sigma})}\\
%      &\hspace{40pt} = (2 \pi)^{d} \delta_{j\ell} \delta_{k(-k')}  \frac{e^{-|t'-t| (\lambda+| k |^{2 \sigma})}}{2(\lambda +|k|^{2 \sigma})}.
%   \end{align*}
%\end{proof}

Our first concern is to study the H\"older-Besov regularity of the process $\vartheta$.

\begin{lemma}\label{l:stationary besov}
   For any $\alpha < \sigma - 1 / 2$ and any $p\ge 1$, the process $\vartheta$ satisfies
   \[
      \E[ \| \vartheta \|_{C_T \CC^{\alpha}(\mathbb{T})}^p] < \infty.
   \]
\end{lemma}

\begin{proof}
   Let $s, t \in [ 0, T ]$ and $\ell \ge 0$. The case $\ell = 0$ can be treated using essentially the same arguments, except that then we need to distinguish the cases $k= 0$ and $k\neq 0$, where $k$ is the argument in the Fourier transform. %Writing $\Delta_{\ell}$ as a convolution operator, we see that $\Delta_{\ell}\vartheta_{s,t}$ is a Gaussian process indexed by $\mathbb{T}$, with values in $\R^n$.
   Using Gaussian hypercontractivity (\cite{Janson1997}, Theorem 3.50), we obtain for $p \ge 1$ that
   \begin{equation}\label{e:burgers hypercontractivity}
      \E[\Vert  \Delta_{\ell} \vartheta_{s,t} \Vert_{L^{2 p} (\mathbb{T})}^{2 p}] \lesssim_p \Vert  \E [| \Delta_{\ell} \vartheta_{s,t}( x ) |^2] \Vert_{L^p_x ( \mathbb{T})}^p.
   \end{equation}
   %where we write $\Vert  f(x) \Vert ^p_{L^p_x (\mathbb{T} )} := \int_{\mathbb{T}} | f ( x )|^p \mathd x$ for $f \in L^p(\mathbb{T})$.
   Using Fourier inversion and Lemma~\ref{l:heat covariance}, we have
%   By definition we have
%   \begin{align*}
%      \Delta_{\ell} \vartheta_{s,t} ( x ) = ( 2 \pi )^{- 1}\sum_{k \in \mathbb{Z}^d} \rho_{\ell} ( k ) \CF{\vartheta}_{s,t}( k ) e^{\imath \langle k, x \rangle},
%   \end{align*}
%   and therefore
   \begin{align*}
      \E[| \Delta_{\ell} \vartheta_{s,t} (x)|^2] & = (2 \pi)^{-2} \sum_{k, k' \in \Z} \rho_{\ell} (k) \rho_{\ell}(k') e^{\imath (k - k') x } \E[\CF{\vartheta}_{s,t} ( k ) \overline{\CF{\vartheta}_{s,t} ( k')}] \\
      & \lesssim \sum_{k \in \Z} \rho^2_\ell(k) |t-s|^{\varepsilon} |k|^{2\sigma(\varepsilon - 1)} \lesssim |t-s|^{\varepsilon} \sum_{k \in \tmop{supp} (\rho_{\ell})} |k|^{2\sigma(\varepsilon - 1)} \lesssim |t-s|^{\varepsilon} 2^{\ell(1 - 2 \sigma (1-\varepsilon))}.
   \end{align*}
%   
%   Since $\rho_{\ell}$ has compact support we are considering finite sums, and therefore exchanging expectation and summation is justified. Now \eqref{e:stationary covariance} implies for $1 \leqslant j \leqslant n$ and independently of $x \in \mathbb{T}$ that
%   \begin{align*}
%      E(|\Delta_{\ell} \vartheta^j_{s,t} ( x )|^2) = (2\pi)^{-d} \sum_{k \in \mathbb{Z}^d} \rho_{\ell}^2 (k)\frac{1 - e^{- |t-s|(\lambda + |k|^{2 \sigma})}}{\lambda + |k|^{2 \sigma}}.
%   \end{align*}
%   For any $\varepsilon \in (0,1]$ we have $1-e^{-x} \lesssim x^\varepsilon$, and therefore
%   \begin{align*}
%      (2\pi)^{-d} \sum_{k \in \mathbb{Z}^d} \rho_{\ell}^2 (k)\frac{1 - e^{- |t-s|(\lambda + |k|^{2 \sigma})}}{\lambda + |k|^{2 \sigma}} & \lesssim |t-s|^{\varepsilon} \sum_{k \in \tmop{supp} (\rho_{\ell})} \frac{1}{(\lambda + | k |^{2 \sigma})^{1-\varepsilon}} \\
%      & \lesssim |t-s|^{\varepsilon} 2^{\ell(d - 2 \sigma (1-\varepsilon))}.
%   \end{align*}
   Hence, we obtain from \eqref{e:burgers hypercontractivity} that
   \[
      \E[\lVert \vartheta(t, \cdot) - \vartheta(s,\cdot)\rVert_{B^\alpha_{2p,2p}(\mathbb{T})}^{2p} ] \lesssim \sum_{\ell \ge -1} 2^{\ell \alpha 2p} \E [\lVert \Delta_\ell \vartheta_{s,t}\rVert_{L^{2p}(\mathbb{T})}^{2p} ] \lesssim \sum_{\ell \ge -1} 2^{\ell \alpha 2p} \left(|t-s|^{\varepsilon} 2^{2\ell(1/2 - \sigma (1-\varepsilon))}\right)^p
   \]
   for any $\alpha \in \R$ and any $p \ge 1$. For $\alpha < \sigma - 1/2$ there exists $\varepsilon \in (0,1]$ small enough so that the series converges. Since we can choose $p$ arbitrarily large, Kolmogorov's continuity criterion %for Banach space valued processes, Theorem A.10 of \cite{Friz2010},
   implies that $\vartheta$ has a continuous version such that $\E[ \|\vartheta \|_{C_T B^{\alpha}_{2 p, 2 p}(\mathbb{T})}^{2p}]< \infty$ for all $\alpha < \sigma - 1/2$. Now we use again that $p$ can be chosen arbitrarily large, so that the Besov embedding theorem, Lemma~\ref{l:besov embedding torus}, shows that this continuous version takes its values in $C_T \CC^{\alpha}(\mathbb{T})$ for all $\alpha < \sigma - 1/2$.
\end{proof}

%Contrary to the one-dimensional case, in higher dimensions the Laplacian does not sufficiently smoothen the white noise. In fact Lemma~\ref{l:stationary besov} shows that if $\sigma = 1$, then already for $d = 2$ we only have a distribution valued solution $\vartheta \in C([0,T], C^{- \varepsilon})$ for any $\varepsilon > 0$. Hence, in that case it is not even clear how to define $G(\vartheta(t,\cdot))$, let alone $G(\vartheta(t,\cdot))\mathD_x \vartheta(t,\cdot)$. So in higher dimensions we need to consider the fractional Laplacian of a sufficiently high order.
%
%If $\sigma - d / 2 \in (0, 1 / 2)$, then $G(\vartheta)$ makes sense, but the product $G(\vartheta)\mathD_x \vartheta$ cannot be defined using classical analytical arguments. This is why prior to {\cite{Hairer2011}} it was not known how to describe solutions to \eqref{e:burgers}. Hairer solved the case $d = 1$, and in the following we show how to use our paracontrolled calculus in order to solve \eqref{e:burgers} for $d \ge 1$. For this purpose we need to construct the area process of $\vartheta(t,\cdot)$.

Next, we construct the area $\vartheta \circ \partial_x \vartheta$.

\begin{lemma}\label{l:burgers area}
   Define
   \[
      \vartheta \circ \partial_x \vartheta = (\vartheta^k \circ \partial_x \vartheta^\ell)_{1\leqslant k, \ell \leqslant n} = \bigg(\sum_{|i-j| \le 1} \Delta \vartheta^i \Delta_j \partial_x \vartheta^j \bigg)_{1\leqslant k, \ell \leqslant n}.
   \]
   Then almost surely $\vartheta \circ \partial_x \vartheta \in C_T\CC^{2\alpha-1}(\mathbb{T}; \R^{n \times n})$ for all $\alpha < \sigma - 1/ 2$. Moreover, if $\psi \in \CS$ is such that $\int \psi(x) \mathd x = 1$ and $\vartheta^\varepsilon = \psi^\varepsilon \ast \vartheta$, where $\psi^\varepsilon = \varepsilon^{-1} \psi(\varepsilon^{-1} \cdot)$, then we have for all $p \ge 1$ that
   \begin{equation}\label{eq:burgers area lp convergence}
      \lim_{\varepsilon \rightarrow 0} \E[\| \vartheta^\varepsilon \circ \partial_x \vartheta^\varepsilon - \vartheta \circ \partial_x \vartheta \|_{C_T \CC^{2\alpha - 1}}^p] = 0.
   \end{equation}
\end{lemma}

\begin{proof}
   Without loss of generality we can argue for $\vartheta^1 \circ \partial_x \vartheta^2$. The case $\vartheta^1 \circ \partial_x \vartheta^1$ is easy, because Leibniz' rule yields $\vartheta^1 \circ \partial_x \vartheta^1 = \frac{1}{2} \partial_x(\vartheta^1 \circ \vartheta^1)$.
  
   Let $\ell \in \N$. Note that if $i$ is smaller than $\ell - N$ for a suitable $N$, and if $| i - j | \leqslant 1$, then $\Delta_{\ell}(\Delta_i f \Delta_j g) = 0$ for all $f, g \in \mathcal{S}'$. Hence, the projection of $\vartheta^1\circ \partial_x\vartheta^2$ onto the $\ell$--th dyadic Fourier block is given by
   \begin{align*}
      \Delta_{\ell} \pi_{\circ}(\vartheta^1, \partial_x \vartheta^2 ) = \sum_{| i- j | \leqslant 1} \Delta_{\ell}(\Delta_i \vartheta^1 \Delta_j\partial_x \vartheta^2) = \sum_{|i-j|\leqslant 1} 1_{\ell \lesssim i} \Delta_{\ell} (\Delta_i \vartheta^1 \Delta_j \partial_x \vartheta^2).
   \end{align*}
   To avoid case distinctions, we only argue for $\ell \ge N$, so that we can always assume $i, j  \ge 0$. The case $\ell < N$ can be handled using essentially the same arguments.
      
   We can apply Gaussian hypercontractivity to obtain
   \begin{gather}\label{e:burgers area hypercontractivity} \nonumber
      \E [\Vert  (\Delta_{\ell} ( \vartheta^1\circ \partial_x \vartheta^2 - \vartheta^{1,\varepsilon} \circ \partial_x \vartheta^{2,\varepsilon}) )_{s,t} \Vert_{L^{2 p}(\mathbb{T})}^{2 p}] \\
      \lesssim \biggl\Vert  \E\biggl[\biggl| \sum_{|i-j| \leqslant 1} 1_{\ell \lesssim i} ( \Delta_{\ell} ( \Delta_i \vartheta^1 \Delta_j \partial_x \vartheta^2 -  \Delta_i \vartheta^{1,\varepsilon} \Delta_j \partial_x \vartheta^{2,\varepsilon} )(x))_{s,t}\biggr|^2\biggr] \biggr\Vert_{L^p_x(\mathbb{T})}^p,
   \end{gather}
   where we write $\vartheta^{1, \varepsilon} = \psi^\varepsilon \ast \vartheta$ and similarly for $\vartheta^{2, \varepsilon}$. 
   
   Let us start by estimating
   \begin{align}\label{e:burgers area pr1} \nonumber
      &\E\biggl[\biggl| \sum_{|i-j| \leqslant 1} 1_{\ell \lesssim i} \Delta_{\ell} ( \Delta_i \vartheta^1(t,\cdot) \Delta_j \partial_x \vartheta^2_{s,t} - \Delta_i \vartheta^{1,\varepsilon} (t,\cdot) \Delta_j \partial_x \vartheta^{2,\varepsilon}_{s,t})(x)\biggr|^2\biggr] \\ \nonumber
      &\hspace{20pt} = \sum_{| i - j | \leqslant 1} \sum_{| i' - j' | \leqslant 1} 1_{\ell \lesssim i} 1_{\ell \lesssim i'} \E[\Delta_{\ell} (\Delta_i \vartheta^1(t,\cdot) \Delta_j \partial_x \vartheta^2_{s,t} - \Delta_i \vartheta^{1,\varepsilon} (t,\cdot) \Delta_j \partial_x \vartheta^{2,\varepsilon}_{s,t})(x)\\
      &\hspace{160pt} \times \overline{\Delta_{\ell} ( \Delta_{i'} \vartheta^1(t,\cdot) \Delta_{j'} \partial_x \vartheta^2_{s,t} - \Delta_{i'} \vartheta^{1,\varepsilon} (t,\cdot) \Delta_{j'} \partial_x \vartheta^{2,\varepsilon}_{s,t})(x)}].
   \end{align}
   Taking the infinite sums outside of the expectation can be justified a posteriori, because for every finite partial sum we will obtain a bound on the $L^2$--norm below, which does not depend on the number of terms that we sum up. The Gaussian hypercontractivity \eqref{e:burgers area hypercontractivity} then provides a uniform $L^p$--bound for any $p \ge 2$, which implies that the squares of the partial sums are uniformly integrable, and thus allows us to exchange summation and expectation.
   
   Recall that $\CF(\CF^{-1}u \CF^{-1}v)(k) = (2\pi)^{-1} \sum_{k'} u(k') v(k -k')$, and $\CF(\partial_x u)(k) = \imath k \CF(u)(k)$, and therefore 
   \begin{gather*}
      \Delta_{\ell} (\Delta_i \vartheta^1(t,\cdot) \Delta_j \partial_x \vartheta^2_{s,t} - \Delta_i \vartheta^{1,\varepsilon} (t,\cdot) \Delta_j \partial_x \vartheta^{2,\varepsilon}_{s,t})(x) \\
      = (2\pi)^{-1} \sum_{k \in \Z} \rho_\ell(k) e^{\imath \langle k, x \rangle} \CF(\Delta_i \vartheta^1(t,\cdot) \Delta_j \partial_x \vartheta^2_{s,t} - \Delta_i \vartheta^{1,\varepsilon} (t,\cdot) \Delta_j \partial_x \vartheta^{2,\varepsilon}_{s,t})(k) \\
      = (2\pi)^{-2} \sum_{k,k' \in \Z} \rho_\ell(k) e^{\imath \langle k, x \rangle} \rho_i(k') \rho_j(k-k') \imath (k-k') \CF \vartheta^1(t,\cdot)(k') \CF\vartheta^2_{s,t}(k-k') \\
      \times (1 - \CF \psi(\varepsilon k') \CF \psi(\varepsilon(k-k'))).
   \end{gather*}
   From this expression it is clear that if we can show $\E[\| \vartheta^\varepsilon \circ \partial_x \vartheta^\varepsilon \|_{C_T \CC^{2\alpha - 1}}^p] < \infty$, then the convergence result in~\eqref{eq:burgers area lp convergence} will follow by dominated convergence, because $\CF \psi$ is bounded and $\CF \psi(0) = 1$ by assumption.
   
   Using the covariance of $\CF \vartheta$ that we calculated in Lemma~\ref{l:heat covariance}, we obtain
   \begin{align*}
      &\E\biggl[\biggl|\sum_{|i-j| \leqslant 1} 1_{\ell \lesssim i} \Delta_{\ell} ( \Delta_i \vartheta^1(t,\cdot) \Delta_j \partial_x \vartheta^2_{s,t})(x)\biggr|^2\biggr] \\
      &\hspace{40pt} \lesssim \sum_{| i - j | \leqslant 1} \sum_{| i' - j' | \leqslant 1} 1_{\ell \lesssim i} 1_{\ell \lesssim i'} \sum_{k, k' \in \mathbb{Z}^d} \rho_{\ell}^2 (k+k') \rho_i ( k ) \rho_{i'} ( k ) \rho_j (k') \rho_{j'}(k')\\
      &\hspace{200pt} \times \frac{1 - e^{-2t|k|^{2\sigma}}}{2 |k|^{2 \sigma}} |k'|^2 |t-s|^{\delta} |k'|^{-2\sigma(1 - \delta)}\\
      &\hspace{40pt} \lesssim \sum_{| i - j | \leqslant 1} 1_{\ell \lesssim i} \sum_{k \in \tmop{supp} ( \rho_i ), k' \in \tmop{supp} ( \rho_j )} \rho_{\ell}^2 ( k + k' ) 2^{2i(1 - 2\sigma + \sigma \delta)} |t-s|^{\delta} \\
      &\hspace{40pt} \lesssim \sum_{i \gtrsim \ell} 2^\ell 2^{2i(1 + 1/2 - 2\sigma + \sigma\delta)} |t-s|^{\delta} 
   \end{align*}
   for all $\delta \in [0,1]$. Since $\sigma > 5/6$, there exists $\delta >0$ small enough so that the sum is finite, and we obtain
   \[
      \E\biggl[\biggl|\sum_{|i-j| \leqslant 1} 1_{\ell \lesssim i} \Delta_{\ell} ( \Delta_i \vartheta^1(t,\cdot) \Delta_j \partial_x \vartheta^2_{s,t})(x)\biggr|^2\biggr]  \lesssim 2^{2i(2 - 2\sigma + \sigma \delta)}  |t-s|^{\delta},
   \]
   and by the same arguments
   \[
      \E\biggl[\biggl|\sum_{|i-j| \leqslant 1} 1_{\ell \lesssim i} \Delta_{\ell} ( \Delta_i \vartheta^1_{s,t} \Delta_j \partial_x \vartheta^2(s,\cdot))(x)\biggr|^2\biggr]  \lesssim 2^{2i(2 - 2\sigma + \sigma \delta)}  |t-s|^{\delta}.
   \]   
%   
%   
%   For any $\varepsilon\in (0,1]$ we can estimate the right hand side by
%   \begin{align*}
%      &\lesssim \sum_{| i - j | \leqslant 1} 1_{\ell\lesssim i} \sum_{k \in \tmop{supp} ( \rho_i ), k' \in\tmop{supp}(\rho_j)} \rho_{\ell}^2(k + k')|k'|^2 \frac{|t-s|^{\varepsilon}(\lambda+|k'|^{2 \sigma})^{\varepsilon}}{2(\lambda + |k'|^{2 \sigma})^2}\\
%      &\lesssim \sum_{i: i \gtrsim \ell} 2^{d i} 2^{d\ell} (2^{i})^2 \frac{|t-s|^{\varepsilon}(2^{i2\sigma})^\varepsilon}{(2^{i2\sigma})^2} = 2^{d \ell} \sum_{i: i \gtrsim \ell} 2^{2 i ( 1+ d / 2 - 2 \sigma ( 1 - \varepsilon ) )} |t - s|^{\varepsilon}.
%   \end{align*}
%    So if $1 + d / 2 - 2 \sigma < 0$, then we obtain for sufficiently small $\varepsilon > 0$ that
%   \begin{align*}
%      E\biggl(\biggl(\sum_{|i-j| \leqslant 1} 1_{\ell \lesssim i} \Delta_{\ell} ( \Delta_i \vartheta^1(t,\cdot) \Delta_j \partial_x \vartheta^2_{s,t})(x)\biggr)^2\biggr) \lesssim 2^{2 \ell ( 1 + d - 2 \sigma (1-\varepsilon))} |t - s|^{\varepsilon},
%   \end{align*}
%   and by the same arguments
%   \begin{align*}
%      E\biggl(\biggl(\sum_{|i-j| \leqslant 1} 1_{\ell \lesssim i} \Delta_{\ell} ( \Delta_i \vartheta^1_{s,t} \Delta_j \partial_x \vartheta^2(s,\cdot))(x)\biggr)^2\biggr) \lesssim 2^{2 \ell (1 + d - 2\sigma(1-\varepsilon))} |t-s|^{\varepsilon}.
%   \end{align*}
   Since
   \[
      \Delta_i \vartheta^1(t,\cdot) \Delta_j \partial_x \vartheta^2(t,\cdot) - \Delta_i \vartheta^1(s,\cdot) \Delta_j \partial_x \vartheta^2(s,\cdot) = \Delta_i \vartheta^1(t,\cdot) \Delta_j \partial_x \vartheta^2_{s,t} + \Delta_i \vartheta^1_{s,t} \Delta_j \partial_x \vartheta^2(s,\cdot),
   \]
   we get for sufficiently small $\delta > 0$ and for arbitrarily large $p\ge 1$ that
   \begin{align*}
      \E [\Vert  \Delta_{\ell} (\vartheta^1 \circ \partial_x \vartheta^2)_{s,t} \Vert_{L^{2 p}(\mathbb{T} )}^{2 p}] \lesssim 2^{-2\ell(2\sigma - 2 - \sigma \delta ) p} |t - s|^{\delta p}.
   \end{align*}
   Now we use the same arguments as in the proof of Lemma \ref{l:stationary besov} to obtain the required $L^p$-bound for $\| \vartheta^1\circ \partial_x \vartheta^2\|_{C_T \CC^{2\alpha-1}}$ with $\alpha < \sigma - 1 / 2$.
\end{proof}

%\begin{remark}\label{r:burgers area continuous}
%   We constructed the Besov area $\vartheta \partial_x \vartheta$ using one fixed partition of unity. But if $(\widetilde{\chi},\widetilde{\rho})$ is another dyadic partition of unity, with dyadic blocks $(\widetilde{\Delta}_i)_{i \ge -1}$, and if $\vartheta \widetilde{\circ} \partial_x \vartheta$ is constructed using that partition, then $\vartheta \partial_x \vartheta$ and $\vartheta \widetilde{\circ} \partial_x \vartheta$ almost surely coincide. This can be seen by considering the difference
%   \begin{align*}
%      \sum_{i,j \leqslant N} \Delta_i \vartheta \Delta_j \partial_x \vartheta - \sum_{i',j' \leqslant N} \widetilde{\Delta}_{i'} \vartheta \widetilde{\Delta}_{j'}\partial_x \vartheta,
%   \end{align*}
%   and by showing that it converges to 0 as $N$ tends to $\infty$.
%   
%   More generally, we expect that for a wide range of smooth approximations $(\vartheta_j)_{j\in\N}$ to $\vartheta$, the product $(\pi_\circ(\vartheta_j, \partial_x \vartheta_j))$ almost surely converges to $\pi_\circ(\vartheta, \partial_x \vartheta)$. In the setting of Gaussian rough paths, such results have been obtained by Friz and Victoir~\cite{Friz2010c}. See also \cite{Friz2012} for the case of SPDEs with a one-dimensional spatial index variable, and Chapter 10 of \cite{Hairer2013a} for some general results in the multidimensional case.
%\end{remark}

Now Lemma~\ref{l:burgers area} and Theorem~\ref{theorem:burgers existence uniqueness} yield the existence and uniqueness of solutions to the fractional Burgers type equation driven by space-time white noise:

\begin{corollary}
  Let $\sigma > 5/6$, $\alpha \in ( 1/3, \sigma -1/2 )$, $T>0$, $G \in C^3_b$, $u_0 \in \CC^\alpha(\mathbb{T})$, $L = \partial_t + (-\Delta)^\sigma$, and let $\xi$ be a space-time white noise on $[0,T] \times \mathbb{T}$ with values in $\R^n$. Then there exists a unique solution $u$ to
  \[
     L u = G(u) \partial_x u + \xi, \qquad u(0) = u_0,
  \]
  in the following sense: If $\psi \in \CS$ with $\int \psi \mathd t=1$ and if for $\varepsilon > 0$ the function $u^\varepsilon$ solves
  \[
     L u^\varepsilon = G(u^\varepsilon) \partial_x u^\varepsilon + \xi^\varepsilon, \qquad u(0) = u_0,
  \]
  where $\xi^\varepsilon = \varepsilon^{-1} \psi(\varepsilon \cdot) \ast \xi$, then $\lim_{\varepsilon \rightarrow 0} \E[\| u - u^\varepsilon \|_{C_T \CC^\alpha(\mathbb{T})}^p] = 0$ for all $p \ge 1$.
%  
%  
%  $( \xi^{\varepsilon} )_{\varepsilon >0}$ is a family of
%  smooth functions on $[ 0,T ] \times \mathbb{T}$ with values in
%  $\mathbb{R}^{n}$, and $G \in C^{3}_{b} ( \mathbb{R}^{n} ,\Lcal (\mathbb{R}^{n} ,\mathbb{R}^{n} ) )$. Suppose that there exist $\vartheta \in C_{T} \CC^{\alpha}$ and $\eta \in C_{T} \CC^{2 \alpha -1}$ such that $(\vartheta^{\varepsilon} , ( \vartheta^{\varepsilon} \circ   \partial_{x} \vartheta^{\varepsilon} ) )$ converges to $( \vartheta , \eta )$ in $C_{T}
%  \CC^{\alpha -1} \times C_{T} \CC^{2 \alpha -1}$, where
%  $\vartheta^{\varepsilon}$ are solutions to $L \vartheta^{\varepsilon} = \xi^{\varepsilon}$ and $\vartheta^{\varepsilon} ( 0 ) =0$, and where $L= \partial_{t} + ( - \Delta )^{\sigma}$. Let for $\varepsilon >0$ the function $u^{\varepsilon}$ be the unique global solution to the Cauchy problem
%  \[
%     L u^{\varepsilon} =G ( u^{\varepsilon} ) \partial_{x} u^{\varepsilon} + \xi^{\varepsilon} , \hspace{2em} u^{\varepsilon} ( 0 ) =u_{0} ,
%  \]
%  where $u_{0} \in \CC^{\alpha}$. Then there exists $u \in C_{T} \CC^{\alpha}$ such that $u^{\varepsilon} \rightarrow u$ in $C_{T} \CC^\alpha$. The limit $u$ depends only on $( u_{0} , \vartheta ,\eta )$, and not on the approximating family $( \vartheta^{\varepsilon} , (\vartheta^{\varepsilon} \circ   \partial_{x} \vartheta^{\varepsilon} ) )$.
\end{corollary}

\begin{remark}
   There is no problem in extending the problem from $\mathbb{T}$ to $\mathbb{T}^d$, %In higher dimensions the space-time white noise is more irregular, and therefore we need to consider the fractional Laplacian $-(-\Delta)^\sigma$ for larger values of $\sigma$ to sufficiently smoothen the noise. But apart from that
   and the analysis works exactly as in the one-dimensional case. The proof of Lemma~\ref{l:stationary besov} shows that if $\xi$ is a space-time white noise on $[0,T] \times \mathbb{T}^d$, then the solution $\vartheta$ to $L \vartheta = \xi$, $\vartheta(0) = 0$, will be in $C_T\CC^\alpha(\mathbb{T}^d)$ for every $\alpha < \sigma - d/2$. So as long as $\sigma - d/2 > 1/3$, we can solve the Burgers equation on $\mathbb{T}^d$. For the existence of the area $\vartheta \circ \partial_x \vartheta$ we need the additional condition $2 \sigma - d/2 - 1 > 0$. But if $\sigma - d/2 > 1/3$, then this is always satisfied. 
\end{remark}



\section{A generalized parabolic Anderson model}\label{sec:PAM}

Consider now the following PDE on $[ 0,T ] \times \mathbb{T}^{2}$ for some fixed $T>0$:
\begin{equation}\label{eq:pam}
  L u=F ( u ) \xi, \qquad u(0) = u_0,
\end{equation}
where $L = \partial_{t} - \Delta$, the function $F$ is continuous from
$\mathbb{R}$ to $\mathbb{R}$, $\xi$ is a spatial white noise, and $u_0 \in \CC^\alpha$ for suitable $\alpha \in \R$.
Formally, this equation is very similar to the rough differential
equation~{\eqref{eq:rde}}.

The regularity of a spatial white noise $\eta$ on $\mathbb{T}^{d}$ is $\eta \in \CC^{-d/2- \varepsilon}$ for every $\varepsilon >0$. Since we are
in dimension $d=2$, we have $\xi \in \CC^{-1- \varepsilon}$. The
Laplacian increases the regularity by 2, so we expect that for fixed $t>0$ we
have $u ( t ) \in \CC^{1- \varepsilon}$, and therefore the product
$F ( u ) \xi$ is ill-defined.

However, let us assume that $\xi \in \CC^{\alpha -2} ( \mathbb{T}^{2} )$ for
some $2/3< \alpha <1$, and let $\vartheta \in \CC^{\alpha}$ be such that $L \vartheta = \xi$. Consider the paracontrolled ansatz
\[
   u=F ( u ) \prec \vartheta +u^{\sharp}
\]
with $u^{\sharp} \in C_{T} \CC^{2 \alpha}$, and where as in Section~\ref{sec:burgers} the paraproduct $\prec$ is only acting on the spatial
variables. If $u$ is of this form, then Lemma~\ref{lemma:para-taylor} and
Lemma~\ref{lemma:commutator} yield
\[
   F ( u ) \xi =F ( u ) \prec \xi +F ( u ) \succ \xi +F' ( u ) F ( u )  (
   \vartheta \circ \xi ) +F' ( u ) C ( F ( u ) , \vartheta , \xi ) +F' ( u )  ( u^{\sharp} \circ \xi ) + \Pi_{F} ( u, \xi )
\]
is well defined provided that $( \vartheta \circ   \xi ) \in C_{T} \CC^{2 \alpha -2}$. Moreover, the algebraic rules for $\partial_{t}$ and $\Delta$
acting on products imply that
\[
   L u= ( L F ( u ) ) \prec \vartheta +F ( u ) \prec L  \vartheta +2 \mathD_{x} F ( u ) \prec \mathD_{x} \vartheta +L u^{\sharp},
\]
and thus we find the following equation for $u^{\sharp}$:
\begin{gather*}
   L u^{\sharp} =-2 \mathD_{x} F ( u ) \prec \mathD_{x} \vartheta - ( L F ( u ) ) \prec \vartheta +F ( u ) \succ \xi +F' ( u ) F ( u )  ( \vartheta \circ \xi ) \\
    +F' ( u ) C ( F ( u ) , \vartheta , \xi ) +F' ( u )  ( u^{\sharp} \circ \xi) + \Pi_{F} ( u, \xi ).
 \end{gather*}
We would like all the terms on the right hand side to be in $C_{T} \CC^{2 \alpha -2}$. However, it is not easy to estimate the term $( LF ( u ) ) \prec \vartheta$ in $C_{T} \CC^{\beta}$ for any $\beta \in \mathbb{R}$: the term $\Delta F ( u )$ can be controlled in $\CC^{\alpha-2}$, but there are no straightforward estimates available for the time derivative $\partial_{t} F ( u )$ appearing in $L F ( u )$. Indeed, it would
be more convenient to treat the generalized parabolic Anderson model in a
space-time parabolic Besov space adapted to the operator $L$ and to use the
natural paraproduct associated to this space. An alternative strategy would be
to stick with the simpler space $C_{T} \CC^{\alpha -2}$ and to observe that
\[
   L F ( u ) =F' ( u ) L u + F'' ( u ) ( \mathD_{x} u )^{2} =F' ( u ) F ( u ) \xi + F'' ( u ) ( \mathD_{x} u )^{2},
\]
and that the terms on the right hand side can be analyzed using the
paracontrolled ansatz. Since this strategy seems to require a lot of
regularity from $F$, we do not pursue it further.



Instead, we keep working on $C_{T} \CC^{\alpha -2}$, but we modify the paraproduct
appearing in the paracontrolled ansatz. Let $\varphi \colon \mathbb{R} \rightarrow \mathbb{R}_{+}$ be a positive smooth function with compact support and total
mass $1$, and for all $i \geqslant -1$ define the operator $Q_{i} \colon C_{T} \CC^{\beta} \rightarrow C_{T} \CC^{\beta}$ by
\[
   Q_{i} f ( t ) = \int_{\mathbb{R}} 2^{-2i} \varphi ( 2^{2i} ( t-s ) ) f ( (s \wedge T ) \vee 0 ) \mathd s.
\]
For $Q_{i}$ we have the following standard estimates, which we leave to the reader to prove:
\begin{gather}\label{eq:Q-bounds}
   \| Q_{i} f ( t ) \|_{L^{\infty}} \leqslant \| f \|_{C_{T} L^{\infty}} , \hspace{1em} \| \partial_{t} Q_{i} f ( t ) \|_{L^{\infty}} \leqslant 2^{( 2-2 \gamma ) i} \| f \|_{C_{T}^{\gamma} L^{\infty}}, \\ \nonumber
   \| ( Q_{i} f-f ) ( t ) \|_{L^{\infty}} \leqslant 2^{-2 \gamma i} \| f \|_{C_{T}^{\gamma} L^{\infty}}
\end{gather}
for all $t \in [ 0,T ]$ and all $\gamma \in ( 0,1 )$; for the second estimate
we use that $\int \varphi' ( t ) \mathd t=0$, and for the third estimate we
use that $\varphi$ has total mass 1.
With the help of $Q_{i}$, let us define a modified paraproduct by setting
\[
   f \precprec g= \sum_{i} ( S_{i-1} Q_{i} f ) \Delta_{i} g
\]
for $f,g \in C_{T} \CS'$. It is easy to show that for this paraproduct we have
essentially the same estimates as for the pointwise paraproduct $f \prec g$,
only that we have to bound $f$ uniformly in time; for example
\[
   \| ( f \precprec g ) ( t ) \|_{\alpha} \lesssim \| f \|_{C_{T} L^{\infty}} \| g ( t ) \|_{\alpha} .
\]
for all $t \ [0,T]$. The modified commutator
\[
   C' ( f,g,h ) = ( f \precprec g ) \circ h-f ( g \circ h )
\]
satisfies a similar estimate as the commutator $C$:
\begin{equation}\label{eq:mod commutator}
   \| C' ( f,g,h ) ( t ) \|_{\alpha + \beta + \gamma} \lesssim ( \| f \|_{C_{T} \CC^\alpha} + \| f \|_{C^{\alpha /2}_{T} L^{\infty}} ) \| g ( t ) \|_{\beta} \| h ( t ) \|_{\gamma}
\end{equation}
for all $t \in [ 0,T ]$. This is shown in the same way as
Lemma~\ref{lemma:commutator}, noting that the proof of
Lemma~\ref{lemma:paraprojection} can be slightly adapted to yield
\[
   \| ( \Delta_{i} ( f \precprec g ) -f \Delta_{j} g ) ( t ) \|_{L^{\infty}} \lesssim 2^{-i ( \alpha + \beta )} ( \| f ( t ) \|_{\alpha} + \| f \|_{C^{\alpha /2}_{T} L^{\infty}} ) \| g ( t ) \|_{\beta}
\]
for $\alpha \in ( 0,1 )$ and $t \in [0,T]$; see also the proof of the following Lemma~\ref{lemma:mod-para-est} for similar arguments.
Letting
\begin{equation}\label{eq:mod-paracontrolled}
  u=F ( u ) \precprec \vartheta +u^{\sharp}
\end{equation}
and redoing the same computation as above gives
\begin{align}\label{eq:phi sharp def pam}\nonumber
   L u^{\sharp} = \Phi^{\sharp} & = -2 \mathD_{x} F ( u ) \precprec \mathD_{x} \vartheta - ( L F ( u ) ) \precprec \vartheta +F ( u ) \succ \xi +F' ( u ) F ( u )  ( \vartheta \circ \xi ) \\
   &\qquad +F' ( u ) C' ( F ( u ) , \vartheta , \xi ) +F' ( u )  ( u^{\sharp} \circ  \xi ) + \Pi_{F} ( u, \xi ) + [ F ( u ) \prec \xi -F ( u ) \precprec \xi ] .
\end{align}
The modified paraproduct takes care of the two most singular terms. As shown
by the following lemma, the remainder $[ F ( u ) \prec \xi -F ( u ) \precprec \xi ]$ is indeed sufficiently smooth, and also the term $( L F ( u ) ) \precprec \vartheta$ can now be easily estimated.

\begin{lemma}\label{lemma:mod-para-est}
  Let $\alpha \in (0,2)$, let $u, \vartheta \in C_{T} \CC^{\alpha}$ and $F \in C^{1}$, and assume moreover that $u \in C_{T}^{\alpha /2} L^{\infty}$.
  Then
  \begin{equation}\label{eq:mod-para-1}
     \| ( ( L F ( u ) ) \precprec \vartheta ) ( t ) \|_{2 \alpha -2} \lesssim (\| F ( u ( t ) ) \|_{C_{T} L^{\infty}} + \| F' ( u ) \|_{C_{T} L^{\infty}} ) ( \| u \|_{C_{T}^{\alpha /2} L^{\infty}} + \| u \|_{C_{T} \CC^\alpha} ) \| \vartheta ( t ) \|_{\alpha}    
  \end{equation}
  for all $t \in [ 0,T ]$, as well as
  \[
     \| ( F ( u ) \prec \xi -F ( u ) \precprec \xi ) ( t ) \|_{2 \alpha -2} \lesssim \| F' ( u ) \|_{C_{T} L^{\infty}} \| u \|_{C_{T}^{\alpha /2} L^{\infty}} \| \xi ( t ) \|_{\alpha -2}.
  \]
\end{lemma}

\begin{proof}
  Consider first
  \[
     ( L F ( u ) ) \precprec \vartheta = \sum_{i} ( S_{i-1} Q_{i} L F ( u )  )  \Delta_{i} \vartheta = \sum_{i} ( L S_{i-1} Q_{i}  F ( u )  ) \Delta_{i} \vartheta.
  \]
  Observe that, as for the standard paraproduct, $( L S_{i-1} Q_{i}  F ( u )) \Delta_{i} \vartheta$ has a space Fourier transform localized in an
  annulus of the form $2^{i} \CA$, so that according to Lemma~\ref{lem: Besov regularity of series} it will be enough to show the following estimate to
  prove the bound~{\eqref{eq:mod-para-1}}:
  \begin{align*}
     &\| ( ( L S_{i-1} Q_{i}  F ( u ) ) \Delta_{i} \vartheta ) ( t ) \|_{L^{\infty}} \\
     &\hspace{60pt} \lesssim ( \| F ( u ( t ) ) \|_{C_{T} L^{\infty}} + \| F' ( u ) \|_{C_{T}  L^{\infty}} ) ( \| u \|_{C_{T}^{\alpha /2} L^{\infty}} + \| u \|_{C_{T} \CC^\alpha} ) \| \vartheta ( t ) \|_{\alpha} 2^{- ( 2 \alpha -2 ) i}
  \end{align*}
  for all $i \geqslant -1$. But
  \begin{align*}
     \| ( L S_{i-1} Q_{i}  F ( u ) ) \|_{L^{\infty}} & \leqslant \| \partial_{t} Q_{i} S_{i-1}  F ( u )  \|_{L^{\infty}} + \| Q_{i} \Delta  S_{i-1}  F ( u )  \|_{L^{\infty}} \\
     & \lesssim 2^{- ( \alpha -2 ) i} \| S_{i-1}  F ( u )  \|_{C_{T}^{\alpha /2} L^{\infty}} + \sum_{j<i-1} \| \Delta \Delta_{j}  F ( u )  \|_{C_{T} L^{\infty}},
  \end{align*}
  where we used the bounds in equation~{\eqref{eq:Q-bounds}}. Now
  \[
     \| S_{i-1}  F ( u )  \|_{C_{T}^{\alpha /2} L^{\infty}} \lesssim \| F ( u )  \|_{C_{T}^{\alpha /2} L^{\infty}} \lesssim \| F' ( u ) \|_{C_{T} L^{\infty}} \| u \|_{C_{T}^{\alpha /2} L^{\infty}},
  \]
  %where in the last step we used that for fixed $x$, the map $t \mapsto u(t,x)$ is continuous from $[0,T]$ to $\R$, and therefore its image is convex. Moreover, we have
  and moreover
  \begin{align*}
     \sum_{j \leqslant i-1} \| \Delta \Delta_{j}  F ( u )  \|_{C_{T} L^{\infty}} & \lesssim \sum_{j \leqslant i-1} 2^{( 2- \alpha ) j} \| F ( u ) \|_{C_{T} \CC^\alpha} \lesssim 2^{( 2- \alpha ) i} \| F ( u ) \|_{C_{T} \CC^\alpha} \\
     & \lesssim 2^{( 2- \alpha ) i} ( \| F ( u ) \|_{C_{T} L^{\infty}} + \| F' ( u ) \|_{C_{T} L^{\infty}} ) \| u \|_{C_{T} \CC^\alpha}.
  \end{align*}
  Thus, we obtain
  \begin{align*}
     \| ( L S_{i-1} Q_{i}  F ( u ) ) \Delta_{i} \vartheta \|_{L^{\infty}} & \lesssim \| ( L S_{i-1} Q_{i}  F ( u )  ) \|_{L^{\infty}} \| \Delta_{i}\vartheta \|_{L^{\infty}} \\
     &\lesssim 2^{( 2-2 \alpha ) i} ( \| F ( u ) \|_{C_{T} L^{\infty}} + \| F' ( u ) \|_{C_{T} L^{\infty}} ) ( \| u \|_{C_{T}^{\alpha /2} L^{\infty}} + \| u \|_{C_{T} \CC^\alpha} ) \| \vartheta \|_{\alpha}
  \end{align*}
  as needed.
  For the second inequality we have
  \[
     F ( u ) \prec \xi -F ( u ) \precprec \xi = \sum_{i} ( Q_{i} S_{i-1}  F ( u ) -S_{i-1}  F ( u ) ) \Delta_{i} \xi,
  \]
  so it will be sufficient to establish that
  \[
     \| ( ( Q_{i} S_{i-1}  F ( u ) -S_{i-1}  F ( u ) ) \Delta_{i} \xi ) ( t )  \|_{L^{\infty}} \lesssim \| F' ( u ) \|_{C_{T} L^{\infty}} \| u \|_{C_{T}^{\alpha /2} L^{\infty}} \| \xi ( t ) \|_{\alpha -2} 2^{( 2-2     \alpha ) i} .
  \]
  But using again~{\eqref{eq:Q-bounds}}, we obtain
  \begin{align*}
     \| ( ( Q_{i} S_{i-1}  F ( u ) -S_{i-1}  F ( u ) ) \Delta_{i} \xi ) ( t ) \|_{L^{\infty}} & \lesssim 2^{-i \alpha} \| S_{i-1} F ( u ) \|_{C_{T}^{\alpha /2} L^{\infty}} \| \Delta_{i} \xi ( t ) \|_{L^{\infty}} \\
     & \lesssim 2^{( 2-2 \alpha ) i} \| F' ( u ) \|_{C_{T} L^{\infty}} \| u \|_{C_{T}^{\alpha /2} L^{\infty}} \| \xi ( t ) \|_{\alpha -2},
  \end{align*}
  and the result is proved.
\end{proof}

Let us now assume $F \in C^{2}_{b}$ and set
\[
   C_{F} = \| F \|_{C^{2}_{b}} , \hspace{1em} C_{\xi} = \| \xi \|_{\CC^{\alpha -2}} + \| \vartheta \|_{C_{T} \CC^\alpha} + \| \vartheta \circ   \xi \|_{C_{T} \CC^{2 \alpha -2}}.
\]
Recall from~{\eqref{eq:phi sharp def pam}} that
\begin{align*}
   \Phi^{\sharp} & = -2 \mathD_{x} F ( u ) \precprec \mathD_{x} \vartheta - ( L F( u ) ) \precprec \vartheta +F ( u ) \succ \xi +F' ( u ) F ( u )  (\vartheta \circ \xi ) \\
   &\qquad  +F' ( u ) C' ( F ( u ) , \vartheta , \xi ) +F' ( u )  ( u^{\sharp} \circ \xi ) + \Pi_{F} ( u, \xi ) + [ F ( u ) \prec \xi -F ( u ) \precprec \xi ],
\end{align*}
so that Lemma~\ref{lemma:mod-para-est} and~{\eqref{eq:mod commutator}} allow
us to derive the estimate
\[
   \| \Phi^{\sharp} ( t ) \|_{2 \alpha -2} \lesssim C_{F} C_{\xi} ( 1+ \| u \|_{C_{T} \CC^\alpha} + \| u \|_{C_{T} \CC^\alpha}^{2} + \| u \|_{C^{\alpha /2}_{T} L^{\infty}} + \| u^{\sharp} ( t ) \|_{2   \alpha} ) +C_{F}^{2} C_{\xi}^{2} ( 1+ \| u \|_{C_{T} \CC^\alpha}),
\]
where we pay attention to indicate that, for fixed $t \in [ 0,T ]$, the
estimate on $\Phi^{\sharp} ( t )$ depends only on the $\CC^{2 \alpha}$ norm of
$u^{\sharp} ( t )$. This will come useful below when introducing the right
norm to control the contribution of the initial condition. Let us summarize our observations so far.

\begin{lemma}
  Let $\alpha \in (2/3, 1)$, $T>0$, $u_0 \in \CC^\alpha$, let $\xi$ be a smooth function, let $\vartheta$ be such that $L \vartheta = \xi$, and let $F \in C^{2}_{b}$. Then
  $u$ solves the PDE
  \[
     L u=F ( u ) \xi , \hspace{2em} u ( 0 ) =u_{0} \in \CC^\alpha
  \]
  on $[ 0,T ]$ if and only if $u=F ( u ) \precprec \vartheta +u^{\sharp}$,
  where $u^{\sharp}$ solves
  \[
     L u^{\sharp} = \Phi^{\sharp} , \hspace{2em} u^{\sharp} ( 0 ) =u_{0} - ( F ( u ) \precprec \vartheta ) ( 0 )
  \]
  on $[ 0,T ]$, and where $\Phi^{\sharp}$ is as defined in {\eqref{eq:phi
  sharp def pam}}. Moreover, for all $t \in [ 0,T ]$ we have the estimate
  \begin{equation}\label{eq:phi sharp estimate pam}
     \| \Phi^{\sharp} ( t ) \|_{2 \alpha -2} \lesssim C_{F} C_{\xi} ( 1+ \| u \|_{C_{T} \CC^\alpha} + \| u \|_{C_{T} \CC^\alpha}^{2} + \| u \|_{C^{\alpha /2}_{T} L^{\infty}} + \| u^{\sharp} ( t ) \|_{2 \alpha} ) +C_{F}^{2} C_{\xi}^{2} ( 1+ \| u \|_{C_{T} \CC^\alpha} ) ,
  \end{equation}
  where
  \[
     C_{\xi} = \| \xi \|_{\CC^{\alpha -2}} + \| \vartheta \|_{C_{T} \CC^\alpha} + \| \vartheta \circ   \xi \|_{C_{T} \CC^{2 \alpha -2}} \hspace{2em} \tmop{and} \hspace{2em} C_{F} = \| F \|_{C^{2}_{b}}.
  \]
\end{lemma}

Next, we would like to close the estimate~\eqref{eq:phi sharp estimate pam},
so that the right hand side only depends on $\Phi^{\sharp}$. In order to
estimate $\| u \|_{C_{T} \CC^\alpha}$ in terms of $\Phi^{\sharp}$,
we observe that
\begin{equation}\label{eq:u-estimate}
  u ( t ) =P_{t} u_{0} + \int_{0}^{t} P_{t-s} \Phi ( s ) \mathd s,
\end{equation}
where $\Phi$ is defined by
\begin{align*}
   \Phi =L u & =F ( u ) \prec \xi +F ( u ) \succ \xi +F' ( u ) F ( u )  (\vartheta \circ \xi ) \\
   &\qquad +F' ( u ) C' ( F ( u ) , \vartheta , \xi ) +F' ( u )  ( u^{\sharp} \circ \xi ) + \Pi_{F} ( u, \xi )
\end{align*}
and can be bounded by
\[
   \| \Phi ( t ) \|_{\alpha -2} \lesssim C_{F} C_{\xi} ( 1+ \| u ( t ) \|_{\alpha} + \| u ( t ) \|_{\alpha}^{2} + \| u^{\sharp} ( t ) \|_{2 \alpha} ) +C_{F}^{2} C_{\xi}^{2} \| u ( t ) \|_{\alpha}.
\]
Using Lemma~\ref{lemma:schauder} (more precisely equation~{\eqref{eq:holder-est-heat}}) and standard estimates for the heat
semigroup, we get from~{\eqref{eq:u-estimate}} that
\[
   \| u \|_{C_{T}^{\alpha /2} L^{\infty}} \lesssim \| u_{0} \|_{\alpha} + \| \Phi \|_{C_{T} \CC^{\alpha -2}},
\]
while from the paracontrolled ansatz~{\eqref{eq:mod-paracontrolled}} we easily get
\begin{equation}\label{eq:pam u bounded by phi sharp}
   \| u \|_{C_{T} \CC^\alpha}
  \lesssim C_{F} C_{\xi} + \| u^{\sharp} \|_{C_{T} \CC^\alpha}.
\end{equation}
Plugging this into~{\eqref{eq:phi sharp estimate pam}}, the estimates for
$\Phi^{\sharp}$ now read as
\[
   \| \Phi^{\sharp} ( t ) \|_{2 \alpha -2} \lesssim C_{F} C_{\xi} ( 1+C_{F} C_{\xi} )^{2} ( 1+ \| u^{\sharp} \|_{C_{T} \CC^\alpha} )^{2} ( 1+ \| u^{\sharp} ( t ) \|_{2 \alpha} ).
\]
Moreover, since $u^{\sharp} ( 0 ) =u_{0} - ( F ( u ) \precprec \vartheta ) ( 0)$ and $L u^{\sharp} = \Phi^{\sharp}$, we obtain
\begin{equation}\label{eq:u sharp differential equation}
   u^{\sharp} ( t ) =P_{t} [ u_{0} - (F ( u ) \precprec \vartheta ) ( 0 ) ] + \int_{0}^{t} P_{t-s} \Phi^{\sharp} (s ) \mathd s.
\end{equation}
Using Lemma~\ref{lemma:heat flow smoothing} with $\delta = \alpha$ and
Lemma~\ref{lemma:schauder} with $\gamma = \alpha /2$ and $\beta = 2\alpha-2$, this leads to
\[
   t^{\alpha /2} \| u^{\sharp} ( t ) \|_{2 \alpha} \lesssim \| u_{0} - ( F ( u) \precprec \vartheta ) ( 0 ) \|_{\alpha} + \sup_{s \in [ 0,t ]} (s^{\alpha /2} \| \Phi^{\sharp} ( s ) \|_{2 \alpha -2} ).
\]
We deduce that
\[
   t^{\alpha /2} \| u^{\sharp} ( t ) \|_{2 \alpha} \lesssim \| u_{0} \|_{\alpha} +C_{F} C_{\xi} + \sup_{s \in [ 0,t ]} ( s^{\alpha /2} \| \Phi^{\sharp} ( s ) \|_{2 \alpha -2} ),
\]
so that our new estimate for $\Phi^{\sharp}$ is
\[
   t^{\alpha /2} \| \Phi^{\sharp} ( t ) \|_{2 \alpha -2} \lesssim C_{F} C_{\xi} ( 1+C_{F} C_{\xi} )^{3} ( 1+ \| u^{\sharp} \|_{C_{T} \CC^\alpha} )^{2} \Big( 1+ \| u_{0} \|_{\alpha} + \sup_{s \in [0,t]} ( s^{\alpha /2} \| \Phi^{\sharp} ( s ) \|_{2 \alpha -2} ) \Big),
\]
which holds uniformly over all $t \in [ 0,T ]$. Using again {\eqref{eq:u sharp differential equation}} in conjuction with Lemma~\ref{lemma:heat flow smoothing} (where we set $\delta =2- \alpha$), we have
\[
   \| u^{\sharp} ( t ) \|_{\alpha} \lesssim \| u_{0} \|_{\alpha} +C_{F} C_{\xi} + \int_{0}^{t} s^{\alpha /2} \| \Phi^{\sharp} ( s ) \|_{2 \alpha -2} \frac{\mathd s}{( t-s )^{1- \alpha /2} s^{\alpha /2}}.
\]
Since the time integral can be bounded by
$
   \int_{0}^{t} \frac{\mathd s}{( t-s )^{1- \alpha /2} s^{\alpha /2}} = \int_{0}^{1} \frac{\mathd s}{( 1-s )^{1- \alpha /2} s^{\alpha /2}} \lesssim 1,
$
we also get that
\begin{equation}\label{eq:pam u sharp uniform alpha bound}
   \| u^{\sharp} \|_{C_{T} \CC^\alpha} \lesssim \| u_{0} \|_{\alpha} +C_{F} C_{\xi} + \sup_{s \in [ 0,T ]} ( s^{\alpha /2} \| \Phi^{\sharp} ( s ) \|_{2 \alpha -2} ) .
\end{equation}
This gives us our final estimate for $\Phi^{\sharp}$:
\begin{equation}\label{eq:pam phi sharp final estimate}
   t^{\alpha /2} \| \Phi^{\sharp} ( t ) \|_{2 \alpha -2} \lesssim C_{F} C_{\xi} ( 1+C_{F} C_{\xi} )^{4} \Big( 1+ \| u_{0} \|_{\alpha} + \sup_{s \in [ 0,t ]} ( s^{\alpha /2} \| \Phi^{\sharp} ( s ) \|_{2 \alpha -2} ) \Big)^{3} .
\end{equation}
In order to use this estimate to bound $\Phi^{\sharp}$, we will apply the
usual scaling argument. Again, we extend the scaling transformation to the
time variable in such a way that it leaves the operator $L$ invariant. More
precisely, we set $\Lambda_{\lambda} u ( t,x ) =u ( \lambda^{2} t, \lambda x)$, so that $L  \Lambda_{\lambda} = \lambda^{2} \Lambda_{\lambda} L$. Now let
$u^{\lambda} = \Lambda_{\lambda} u$, $\xi^{\lambda} = \lambda^{2- \alpha} \Lambda_{\lambda} \xi$, and $\vartheta^{\lambda} = \lambda^{\alpha} \Lambda_{\lambda} \vartheta$. Note that $u^{\lambda} \colon [ 0,T/ \lambda^{2} ] \times \mathbb{T}_{\lambda}^{2} \rightarrow \mathbb{R}$, where
$\mathbb{T}_{\lambda}^{2} = (\R / (2\pi \lambda^{-1} \Z))^{2}$ is a rescaled torus,
and that $u^{\lambda}$ solves the equation
\[
   L u^{\lambda} = \lambda^{2} F ( u^{\lambda} ) \Lambda_{\lambda} \xi = \lambda^{\alpha} F ( u^{\lambda} ) \xi^{\lambda}, \qquad u^\lambda_0 = \Lambda_\lambda u_0.
\]
The scaling is chosen in such a way that $\| \xi^{\lambda} \|_{\CC^{\alpha -2}} \lesssim \| \xi \|_{\CC^{\alpha -2}}$, $\| \vartheta^{\lambda} \|_{C_{T} \CC^\alpha} \lesssim \| \vartheta \|_{C_{T} \CC^\alpha}$, and $\| \vartheta^{\lambda} \circ \xi^{\lambda} \|_{C_{T} \CC^{2 \alpha -2}} \lesssim \| \vartheta \circ \xi \|_{C_{T} \CC^{2 \alpha -2}}$, $\| \Lambda_\lambda u_0 \|_\alpha \lesssim \|u_0\|_\alpha$, uniformly over $\lambda \in (0,1]$. In particular we have $C_{\xi^{\lambda}} \lesssim C_{\xi}$ and
$C_{\lambda^{\alpha} F} = \lambda^{\alpha} C_{F}$ for all $\lambda >0$.
%
We define $u^{\sharp , \lambda}$ in the same way as $u^{\sharp}$, with
$u^{\lambda}$ and $\vartheta^{\lambda}$ replacing $u$ and $\vartheta$, and
similarly for $\Phi^{\sharp , \lambda}$. The same derivation as for
{\eqref{eq:pam phi sharp final estimate}} shows that
\[
   \sup_{t \in [ 0,T ]}  ( t^{\alpha /2} \| \Phi^{\sharp , \lambda} ( t ) \|_{2 \alpha -2} ) \lesssim \lambda^{\alpha} C_{F} C_{\xi} ( 1+C_{F} C_{\xi} )^{4} \Big( 1+ \| \Lambda_{\lambda} u_{0} \|_{\alpha} + \sup_{t \in [ 0,T ]} ( t^{\alpha /2} \| \Phi^{\sharp , \lambda} ( t ) \|_{2 \alpha -2} ) \Big)^{3},
\]
%Since we also know from Lemma~\ref{lemma:scaling} that
%\[%\label{eq:pam phi sharp final estimate} 
%   \sup_{t \in [ 0,T ]} \label{eq:pam phi sharp final estimate} ( t^{\alpha /2} \| \Phi^{\sharp} ( t ) \|_{2 \alpha -2} ) \lesssim \sup_{t \in [ 0,T ]} ( t^{\alpha /2} \| \Phi^{\sharp , \lambda} ( t ) \|_{2 \alpha -2} ),
%\]
%this proves that
and therefore
\[
   \sup_{t \in [ 0,T ]}  ( t^{\alpha /2} \| \Phi^{\sharp , \lambda} ( t ) \|_{2 \alpha -2} ) \lesssim 1
\]
as long as $\lambda >0$ is small enough. From~\eqref{eq:pam u bounded by phi sharp} in combination with~\eqref{eq:pam u sharp uniform alpha bound} we
also obtain
\[
   \| u^{\lambda} \|_{C_{T} \CC^{\alpha}} \lesssim 1+ \| u_{0} \|_{\alpha}.
\]
Similar arguments show that if $F \in C^{3}_{b}$, then for every $a>0$ there
exists a sufficiently small $\lambda >0$, such that the map $(\Lambda_{\lambda} u_{0} , \xi^{\lambda} , \vartheta^{\lambda} , \xi^{\lambda} \circ \vartheta^{\lambda} ) \mapsto u^{\lambda} \in C ( [ 0,T ] , \CC^{\alpha} ( \mathbb{T}_{\lambda}^{2} ,\mathbb{R} ) )$ is
Lipschitz continuous on the set
\[
   \max \{ \| u_{0} \|_{\alpha} , \| \xi \|_{\alpha -2} , \| \vartheta \|_{\alpha} , \| \xi \circ \vartheta \|_{2 \alpha -2} \} \leqslant a.
\]
Since $u= \Lambda_{\lambda^{-1}} u^{\lambda}$, we also obtain that $u$
restricted to $[ 0, \lambda^{2} T ]$ depends in a locally Lipschitz continuous
way on the data.

\subsection{Renormalization}

Our previous results show that for given $\xi \in \CC^{\alpha -2}$ with associated area $\xi \circ \vartheta \in C_{T} \CC^{2 \alpha -2}$
there exists a unique local solution $u$ to {\eqref{eq:pam}} with initial
condition $u_{0}$, which is the limit of the solutions $u^{\varepsilon}$,
$\varepsilon >0$, to
\[
   L u^{\varepsilon} =F ( u^{\varepsilon} ) \xi^{\varepsilon} , \hspace{2em} u^{\varepsilon} ( 0 ) =u_{0},
\]
provided that $\xi^{\varepsilon} \rightarrow \xi$ in $\CC^{\alpha -2}$
and $\xi^{\varepsilon} \circ \vartheta^{\varepsilon} \rightarrow \xi \circ \vartheta$ in $C_{T} \CC^{2 \alpha -2}$ as $\varepsilon \rightarrow 0$. Note that here the superscript $\varepsilon$ refers to a smooth regularization of the noise, whereas in the previous section the
superscript $\lambda$ refered to a rescaling. In the following we will not
consider any more rescalings, so that there should be no confusion.

One further difficulty is the fact that the regularized term $(\vartheta^{\varepsilon} \circ \xi^{\varepsilon} )$ does not converge in some
relevant cases; in particular, if $\xi$ is a spatial white noise. However, what we will show below is that for the white noise  there
exist constants $c_{\varepsilon} \in \mathbb{R}$ such that $((\vartheta^{\varepsilon} \circ \xi^{\varepsilon} ) -c_{\varepsilon})$ converges (almost surely with respect to the white noise law)
in $C_{T} \CC^{2 \alpha -2}$. In order to make the term $c_{\varepsilon}$
appear in the equation, we can introduce a suitable correction term in the
regularized problems and consider the renormalized PDE
\begin{equation}\label{eq:pam renormalized}
   L u^{\varepsilon} =F ( u^{\varepsilon} ) \xi^{\varepsilon} -c_{\varepsilon} F' ( u^{\varepsilon} ) F ( u^{\varepsilon} ) .
\end{equation}
For this equation we use again the paracontrolled ansatz~\eqref{eq:mod-paracontrolled}. The same derivation as for~\eqref{eq:phi sharp def pam} yields
\begin{gather*}
   L u^{\sharp , \varepsilon} = \Phi^{\sharp , \varepsilon} =-2 \mathD_{x} F (u^{\varepsilon} ) \prec \mathD_{x} \vartheta^{\varepsilon} -L F ( u^{\varepsilon} ) \prec \vartheta^{\varepsilon} +F ( u^{\varepsilon} ) \succ \xi^{\varepsilon} +F' ( u^{\varepsilon} ) F ( u^{\varepsilon} )  [ (\vartheta^{\varepsilon} \circ \xi^{\varepsilon} ) -c_{\varepsilon} ] \\
   +F' ( u^{\varepsilon} ) C' ( F ( u^{\varepsilon} ) ,  \vartheta^{\varepsilon} , \xi^{\varepsilon} ) +F' ( u^{\varepsilon} )  ( u^{\sharp , \varepsilon} \circ \xi^{\varepsilon} ) +R ( F' ( u^{\varepsilon} ) ,u^{\varepsilon} , \xi^{\varepsilon} ) + [ F ( u^{\varepsilon} ) \prec \xi^{\varepsilon} -F ( u^{\varepsilon} ) \precprec \xi^{\varepsilon} ].
\end{gather*}
for $\varepsilon >0$. We see that $\Phi^{\sharp , \varepsilon}$ only depends on $\xi^{\varepsilon}$, $\vartheta^{\varepsilon}$, and $(\vartheta^{\varepsilon} \circ \xi^{\varepsilon} ) -c_{\varepsilon}$. Thus,
convergence of $(\xi^{\varepsilon})$ to $\xi$ in
$\CC^{\alpha -2}$ and convergence of $( \vartheta^{\varepsilon} \circ \xi^{\varepsilon} -c_{\varepsilon} )$ to $\eta$ in $C_{T} \CC^{2 \alpha -2}$ imply that the solutions $( u^{\varepsilon} )$ to~\eqref{eq:pam renormalized} converge to a limit which only depends on $\xi$ and $\eta$,
but not on the approximating family.

\begin{theorem}\label{theorem:pam existence uniqueness}
  Let $\alpha \in ( 2/3,1 )$ and assume that $( \xi^{\varepsilon})_{\varepsilon >0}$ is a family of smooth functions on $\mathbb{T}^{2}$
  with values in $\mathbb{R}$, and $F \in C^{3}_{b} ( \mathbb{R},\mathbb{R})$. Suppose that there exist $\xi \in \CC^{\alpha -2}$ and $\eta \in C_{T} \CC^{2 \alpha -2}$ such that $( \xi^{\varepsilon} , (\vartheta^{\varepsilon} \circ   \xi^{\varepsilon} ) -c_{\varepsilon} )$
  converges to $( \xi , \eta )$ in $C_{T} \CC^{\alpha -2} \times C_{T} \CC^{2 \alpha -2}$, where $\vartheta^{\varepsilon}$ are solutions to $L \vartheta^{\varepsilon} = \xi^{\varepsilon}$, and where $c_{\varepsilon} \in \mathbb{R}$ for all $\varepsilon >0$.
  Let for $\varepsilon >0$ the function $u^{\varepsilon}$ be the unique global
  solution to the Cauchy problem
  \[
     L u^{\varepsilon} =F ( u^{\varepsilon} ) \xi^{\varepsilon} -c_{\varepsilon} F' ( u^{\varepsilon} ) F ( u^{\varepsilon} ) , \hspace{2em} u^{\varepsilon} ( 0 ) =u_{0},
  \]
  where $u_{0} \in \CC^{\alpha}$. Then there exists $\tau = \tau (u_{0} , \xi , \eta ) >0$ and $u \in C_{\tau} \CC^{\alpha}$ such that
  $u^{\varepsilon} \rightarrow u$ in $C_{\tau} \CC^{\alpha'}$ for
  every $\alpha' < \alpha$. The limit $u$ depends only on $( u_{0} , \xi , \eta )$, and not on the approximating family $( \xi^{\varepsilon} , (\vartheta^{\varepsilon} \circ   \xi^{\varepsilon} ) -c_{\varepsilon} )$.
\end{theorem}

As for the previous equations, $u$ is a weak solution to $L u=F ( u )\diamond \xi$ with $u ( 0 ) =u_{0}$, if we use the correct interpretation
for the renormalized product $F ( u ) \diamond\xi$, and also $u$ depends continuously on the initial condition $u_0$.








\subsection{Regularity of the area and renormalized products}

It remains to study the regularity of the area $\vartheta \circ \xi$. As already indicated, we will have to renormalize the product by ``subtracting an infinite constant'' in order to obtain a well-defined object.

To simplify the arguments below, we assume that $\xi$ is given by $\tilde{\xi} - (2\pi)^{-2} \CF \tilde{\xi}(0)$, where $\tilde{\xi}$ is a spatial white noise on $\mathbb{T}^2$. Then $(\CF \xi (k))_{k\in \Z^2}$ is a complex valued, centered Gaussian process with covariance
\[
   \E[\CF \xi (k) \overline{\CF \xi(k')}] = (2\pi)^2 \mathbf{1}_{k = k'} \mathbf{1}_{k\neq 0}.
\]
Since $\xi$ is a smooth additive perturbation of $\tilde{\xi}$, this simplification will pose no problems, and we indicate below how to handle $\tilde{\xi}$ once we are able to handle $\xi$.

%and therefore if the solutions to $(u^\varepsilon)$ to
%\[
%   L u^\varepsilon = F(u^\varepsilon) \xi^\varepsilon -c_{\varepsilon} F' ( u^{\varepsilon} ) F ( u^{\varepsilon} ), \qquad u^\varepsilon(0) = u_0,
%\]
%converge to a limit $u$, then also the solutions $(\tilde{u}^\varepsilon)$ to
%\[
%   L \tilde{u}^\varepsilon = F(\tilde{u}^\varepsilon) \tilde{\xi}^\varepsilon - c_{\varepsilon} F' ( \tilde{u}^{\varepsilon} ) F ( \tilde{u}^{\varepsilon} ), \qquad \tilde{u}^\varepsilon(0) = u_0,
%\]
%converge to a limit $\tilde{u}$ - provided that $\xi^\varepsilon$ and $\tilde{\xi}^\varepsilon$ are 
%
%t is easy to later on include the mean of $\tilde{\xi}$ again, and to obtain the convergence result for the equation

Since $\xi$ is a mean zero distribution, there exists a stationary solution $\vartheta$ to $L \vartheta = \xi$. Indeed, let
\[
   \vartheta(x) = \int_0^\infty P_t \xi (x) \mathd t.
\]
Then $L \vartheta = \xi$ by definition, and it is easily verified that $(\CF \vartheta(k))$ is a centered, complex valued Gaussian process with covariance
\[
   \E[ \CF \vartheta(k) \overline{\CF \vartheta(k')}] = (2\pi)^2 \frac{1}{|k|^4} \mathbf{1}_{k = k'} \mathbf{1}_{k \neq 0}.
\]
This yields, using Gaussian hypercontractivity and Besov embedding, that $\E[\| \vartheta\|_{\CC^\alpha(\mathbb{T}^2)}^p] < \infty$ for all $\alpha < 1$ and $p \ge 1$. Alternatively, we could also show that $\E[\|\xi\|_{\CC^{\alpha-2}(\mathbb{T}^2)}^p] < \infty$ and then use Lemma~\ref{lemma:heat flow smoothing} and Lemma~\ref{lemma:schauder}, as well as the fact that $\CF \xi$ is supported away from 0, to estimate $\| \vartheta\|_{\CC^\alpha(\mathbb{T}^2)} \lesssim \|\xi \|_{\CC^{\alpha-2}(\mathbb{T}^2)}$.

%Using Lemma~\ref{lemma:schauder}, it is easily verified that $\vartheta \in \CC^{\alpha}$ for any $\alpha < 1$.
%
%
%Recall that the spatial white noise $\xi$ is a centered Gaussian process with values in $\CS'(\mathbb{T}^2)$, such that
%\[
%   \E[\langle \xi, \varphi \rangle \langle \xi, \vartheta\rangle] = \int_{\mathbb{T}^2} \varphi(x) \vartheta(x) \mathd x
%\]
%for all test function $\varphi, \vartheta \in L^2(\mathbb{T}^2)$. %Formally we write $E(\xi(x) \xi(x')) = \delta(x-x')$, $x, x' \in \R^2$.
%We can then write $\xi = (\xi - (2\pi)^{-1} \CF \xi(0)) + (2\pi)^{-1} \CF \xi(0)$, where $c = \int_{\mathbb{T}^2} \xi(x) \mathd x = \CF \xi (0)$, and where the Fourier transform of
%
%
%
%The solution $\vartheta$ to $\vartheta(0) = 0$ and $L \vartheta = \xi$ is given by
%\[
%   \vartheta(t,x) = \int_0^t P_{s} \xi (x) \mathd s
%\]
%where $(P_t)_{t \ge 0}$ denotes the semigroup generated by the Laplacian. Note that for $t > 0$ the function $P_t \xi$ is smooth, and therefore $P_{t-s} \xi \circ \xi$ is a well-defined smooth function, so that we could formally set
%\[
%   (\vartheta \circ \xi) (t,x) = \int_0^t (P_{s} \xi \circ \xi)(x) \mathd s.
%\]
%However, this expression is not well defined:

Since $P_t \xi$ is a smooth function for every $t > 0$, the resonant term $P_t \xi \circ \xi$ is a smooth function, and therefore we could formally set $\vartheta \circ \xi = \int_0^\infty (P_t \xi \circ \xi) \mathd t$. However, this expression is not well defined:

\begin{lemma}\label{l:anderson area expectation}
   For any $x \in \mathbb{T}^2$ and $t > 0$ we have
   \[
      g_t = \E[(P_t \xi \circ \xi)(x)] = \E[ \Delta_{- 1} (P_t \xi \circ \xi ) ( x ) ] = \sum_{k \in\mathbb{Z}^2\setminus \{0\}} e^{- t | k |^2}.
   \]
   In particular, $g_t$ does not depend on the specific partition of unity used to define the $\circ$ operator, and $\int_0^t g_s \mathd s = \infty$ for all $t > 0$.
\end{lemma}

\begin{proof}
   Let $x \in \mathbb{T}^2$, $t > 0$, and $\ell \ge - 1$. Then
   \begin{align*}
      \E[\Delta_\ell ( P_t \xi \circ \xi) ( x ) ] = \sum_{| i - j |\leqslant 1} \E [ \Delta_\ell ( \Delta_i ( P_t \xi) \Delta_j \xi ) ( x ) ], 
   \end{align*}
   where exchanging summation and expectation is justified because it can be easily verified that the partial sums of $\Delta_\ell ( P_t \xi \circ \xi ) ( x )$ are uniformly $L^p$--bounded for any $p \ge 1$. %Now $( \CF{\xi} ( k ))_{k \in \mathbb{Z}^2}$ is a complex valued centered Gaussian process with covariance $\E [ \CF{\xi} ( k ) \CF{\xi}(k')] = (2 \pi)^2 \mathbf{1}_{k = -k'}$.
   Now recall that $P_t = e^{-t\lvert \cdot\rvert^2}(\mathD)$, and therefore
   \begin{align*}
      \E[\Delta_\ell ( \Delta_i ( P_t \xi )\Delta_j \xi ) ( x ) ] & = ( 2 \pi )^{-2} \sum_{k, k' \in \mathbb{Z}^2 \setminus \{0\}} e^{\imath \langle k+k', x \rangle}\rho_m ( k + k' ) \rho_i ( k ) e^{- t | k|^2} \rho_j (k') \E [\CF{\xi} ( k ) \CF{\xi}(k')]\\
      & = \sum_{k \in \mathbb{Z}^2\setminus \{0\}} \rho_m ( 0 ) \rho_i ( k) e^{- t | k |^2} \rho_j ( k ) = \mathbf{1}_{m=0} \sum_{k \in \mathbb{Z}^2 \setminus \{0\}} \rho_i(k) \rho_j ( k ) e^{- t | k |^2}.
   \end{align*}
   For $| i - j | > 1$ we have $\rho_i(k) \rho_j(k) = 0$. This implies, independently of $x \in \mathbb{T}^2$, that
   \[
      g_t = \E[ ( P_t \xi \circ \xi)(x)] = \sum_{k \in \mathbb{Z}^2\setminus \{0\}} \sum_{i, j} \rho_i ( k) \rho_j ( k ) e^{- t | k |^2} = \sum_{k \in \mathbb{Z}^2 \setminus{\{0\}}} e^{- t | k |^2} ,
   \]
   while $\E[(P_t \xi \circ \xi)(x) - \Delta_{- 1} ( P_t \xi \circ \xi))( x ) ] = 0$ for every $x \in \mathbb{T}^2$.
\end{proof}

\begin{remark}
   The same calculation shows that if $\psi \in \CS$, and if $\xi^\varepsilon = \varepsilon^{-2} \psi (\varepsilon^{-1} \cdot) \ast \xi$, then
   \[
      \E[ (P_t \xi^\varepsilon \circ \xi^\varepsilon) ( x ) ] = \E[ \Delta_{- 1} (P_t \xi^\varepsilon \circ \xi^\varepsilon ) ( x ) ] = \sum_{k \in\mathbb{Z}^2\setminus \{0\}} e^{- t | k |^2} |\CF \psi (\varepsilon k)|^2.
   \]
\end{remark}

The diverging time integral motivates us to study the renormalized product $\vartheta \circ \xi  - \int_0^\infty g_t \mathd t$, where $\int_0^\infty g_t \mathd t$ is an infinite constant:

\begin{lemma}\label{l:pam area}
   Set
   \[
      (\vartheta \diamond \xi) = \int_0^\infty (P_t \xi \circ \xi - g_t) \mathd t.
   \]
   Then $\E [\| \vartheta \diamond \xi\|_{\CC^{2\alpha - 2}(\mathbb{T}^2)}^p]<\infty$ for every $\alpha < 1$ and every $p \ge 1$. Moreover, if $\psi \in \CS$ satisfies $\int \psi(x) \mathd x = 1$, and if $\xi^\varepsilon = \psi^\varepsilon \ast \xi$ with $\psi^\varepsilon = \varepsilon^{-2} \psi(\varepsilon \cdot)$ for $\varepsilon > 0$, and $\vartheta^\varepsilon = \int_0^\infty P_t \xi^\varepsilon \mathd t$, then
   \[
      \lim_{\varepsilon \rightarrow 0} \E[ \| \vartheta \diamond \xi - ( \vartheta^\varepsilon \circ \xi^\varepsilon - c_\varepsilon) \|_{\CC^{2\alpha - 2}(\mathbb{T}^2)}^p] = 0
   \]
   for all $p \ge 1$, where for $x \in \mathbb{T}^2$
   \[
      c_\varepsilon = \E[\vartheta^\varepsilon(x) \xi^\varepsilon(x)] = \E[\vartheta^\varepsilon \circ \xi^\varepsilon(x)] = \sum_{k \in \Z^2 \setminus\{0\}} \frac{|\psi(\varepsilon k)|^2}{|k|^2}.
   \]
%   , which in fact does not depend on $x \in \mathbb{T}^2$.
%   
%   For $t > 0$ we define
%   \begin{align*}
%      \Xi_t ( x ) := \pi_{\circ} ( P_t \xi, \xi ) (x) - g_t .
%   \end{align*}
%   Then for all $\varepsilon \in (0,1)$ and for all $t > 0$ we have $E ( \Vert  \Xi_t \Vert_{- \varepsilon}) \lesssim t^{- 1 + \frac{\varepsilon}{4}}$. In particular
%   \begin{align*}
%      E\biggl( \int_0^T \Vert  \Xi_t \Vert_{- \varepsilon} \mathd t \biggr) < \infty
%   \end{align*}
%   for all $T, \varepsilon > 0$.
\end{lemma}

\begin{proof}
   We split the time integral into two components, $\int_0^1 (P_t \xi \circ \xi - g_t) \mathd t$ and $\int_1^\infty (P_t \xi \circ \xi - g_t) \mathd t$. The second integral can be treated without relying on probabilistic estimates: For $x \in \mathbb{T}^2$ we have
   \begin{gather*}
      \bigg\| \int_1^\infty (P_t \xi \circ \xi - g_t) \mathd t - \int_1^\infty (P_t \xi^\varepsilon \circ \xi^\varepsilon - \E[P_t \xi^\varepsilon \circ \xi^\varepsilon(x)]) \mathd t \bigg\|_{\CC^{2\alpha-2}(\mathbb{T}^2)} \\
      \lesssim \int_1^\infty \| P_t \xi \circ \xi - P_t \xi^\varepsilon \circ \xi^\varepsilon \|_{2\alpha + 1} \mathd t + \int_1^\infty \sum_{k \in \Z^2 \setminus\{0\}} e^{-t |k|^2} |1 - |\psi(\varepsilon k)|^2| \mathd t \\
      \lesssim  \int_1^\infty (\| P_t (\xi - \xi^\varepsilon) \|_{2\alpha + 3} \| \xi \|_{\alpha-2} + \| P_t \xi^\varepsilon \|_{2 \alpha + 3} \| \xi - \xi^\varepsilon \|_{\alpha - 2}) \mathd t + \sum_{k \in \Z^2 \setminus\{0\}} \frac{e^{- |k|^2}}{|k|^2} |1 - |\CF \psi(\varepsilon k)|^2|,
   \end{gather*}
   Now since $\CF \xi (0) = 0$, the estimate $\| P_t \xi^\varepsilon \|_{\alpha + 3} \lesssim t^{-3/2} \| \xi^\varepsilon \|_\alpha$ of Lemma~\ref{lemma:heat flow smoothing} holds uniformly over $t > 0$, and thus the time integral is finite. The claimed convergence in $L^p(\mathbb{P})$ now easily follows from the dominated convergence theorem.
   
   We will treat the integral from 0 to 1 using similar arguments as in the proof of Lemma~\ref{l:burgers area}. To lighten the notation, we will only show that $\E[\| \int_0^1 (P_t \xi \circ \xi - g_t) \mathd t\|_{\CC^\alpha}^p]<\infty$. The difference
   \[
      \E[ \| \vartheta \circ \xi - ( \vartheta^\varepsilon \circ \xi^\varepsilon - c_\varepsilon) \|_{\CC^{2\alpha - 2}(\mathbb{T}^2)}^p]
   \]
   can be treated with the same arguments, we only have to include some additional factors of the form $|1 - \CF \psi (\varepsilon k))|^2$ in the sums below. The convergence of the expectation can then be shown using dominated convergence.
   
   Let $t \in (0,1]$ and define $\Xi_t = P_t \xi \circ \xi - g_t$. By Gaussian hypercontractivity we obtain for $p\ge 1$ and $m \ge - 1$ that
   \begin{equation}\label{e:anderson area hyper}
      \E [\Vert  \Delta_m \Xi_t \Vert_{L^{2 p} ( \mathbb{T}^2 )}^{2 p}] \lesssim_p \lVert \E[|\Delta_m \Xi_t ( x ) |^2] \rVert ^p_{L^p_x(\mathbb{T}^2)}.
   \end{equation}
   Lemma~\ref{l:anderson area expectation} yields $\Delta_m g_t = 0 = \E [\Delta_m ( P_t \xi\circ \xi )(x)]$ for $m \ge 0$ and $x \in \mathbb{T}^2$, and therefore
   \begin{equation}\label{e:anderson area m-block}
      \E[| \Delta_m \Xi_t(x)|^2] = \E[| \Delta_m ( P_t \xi\circ \xi ) ( x )|^2] = \tmop{Var} ( \Delta_m ( P_t \xi\circ \xi ) ( x ) ),
   \end{equation}
   where $\tmop{Var}(\cdot)$ denotes the variance. Lemma \ref{l:anderson area expectation} also gives $\Delta_{- 1} g_t = g_t= \E[ \Delta_{- 1} ( P_t \xi\circ \xi )(x) ]$, so
   \begin{equation}\label{e:anderson area -1-block}
      \E [| \Delta_{- 1} \Xi_t ( x ) |^2 ] = \E[ |\Delta_{- 1} ( P_t \xi\circ \xi ) (x ) - g_t |^2] = \tmop{Var} ( \Delta_{- 1} (P_t \xi\circ \xi ) ( x ) ) .
   \end{equation}
   Therefore, it suffices to control $\tmop{Var} ( \Delta_m ( P_t \xi\circ \xi ) ( x ) )$. We have
%   \begin{align*}
%      \CF  ( \pi_{\circ} ( P_t \xi, \xi ) ) (k) = ( 2 \pi )^{- 2} \sum_{k' \in \mathbb{Z}^2} \sum_{| i - j | \leqslant 1} \rho_i ( k-k' ) e^{- t | k-k' |^2} \CF{\xi} (k-k') \rho_j ( k' ) \CF{\xi} ( k' ),
%   \end{align*}
%   implying for all $m \ge - 1$ and $x \in \mathbb{T}^2$ that
   \begin{gather*}
      \Delta_m ( P_t \xi\circ \xi ) (x) = (2\pi)^{-2} \sum_{k \in \Z^2} e^{\imath\langle k,x\rangle} \rho_m(k) \CF(P_t \xi\circ \xi))(k)\\
       = ( 2 \pi )^{- 4} \sum_{k_1, k_2 \in \mathbb{Z}^2}\sum_{| i - j | \leqslant 1} e^{\imath \langle k_1 + k_2, x \rangle} \rho_m ( k_1 + k_2 ) \rho_i ( k_1 )     e^{- t | k_1 |^2} \CF{\xi} ( k_1 ) \rho_j (k_2) \CF{\xi} ( k_2 ).
   \end{gather*}
   Hence
   \begin{align*}
      &\tmop{Var} ( \Delta_m ( \pi_{\circ} ( P_t \xi, \xi )) ( x ) )\\
      &\hspace{20pt} = ( 2 \pi )^{- 8} \sum_{k_1, k_2 \in \mathbb{Z}^2} \sum_{k'_1, k'_2 \in \mathbb{Z}^2} \sum_{| i - j | \leqslant 1} \sum_{| i' - j' | \leqslant 1} e^{\imath \langle k_1 + k_2, x \rangle} \rho_m ( k_1 + k_2 ) \rho_i ( k_1) e^{- t | k_1 |^2} \rho_j ( k_2 )\\
      &\hspace{65pt} \times e^{\imath \langle k'_1 + k'_2, x \rangle} \rho_m (k'_1 + k'_2 ) \rho_i ( k'_1 ) e^{- t | k'_1 |^2}\rho_j ( k'_2 ) \tmop{cov}(\CF{\xi} ( k_1 )\CF{\xi} ( k_2 ), \CF{\xi} ( k'_1 ) \CF{\xi}( k'_2 ) ),
   \end{align*}
   where exchanging summation and expectation can again be justified a posteriori by the uniform $L^p$-boundedness of the partial sums, and where $\tmop{cov}(X,Y)$ denotes the covariance of $X$ and $Y$.
    
   Now $( \widehat{\xi} (k) )_{k \in \mathbb{Z}^2}$ is a centered Gaussian process, and therefore Wick's theorem (\cite{Janson1997}, Theorem 1.28) yields
   \begin{align*}
      &\tmop{cov} ( \widehat{\xi} ( k_1 ) \widehat{\xi} (k_2), \widehat{\xi} ( k'_1 ) \widehat{\xi} (k'_2)) = ( 2 \pi )^4 ( \mathbf{1}_{k_1=-k_1'} \mathbf{1}_{k_2 = - k'_2} + \mathbf{1}_{k_1 = - k'_2} \mathbf{1}_{k_2 = - k_1'} ) \mathbf{1}_{k_1 \neq 0} \mathbf{1}_{k_2 \neq 0}.
   \end{align*}
   Therefore,
   \begin{align*}
      &\tmop{Var} ( \Delta_m ( \pi_{\circ} ( P_t \xi, \xi )) ( x ) ) \\
      &\hspace{20pt} = \sum_{k_1, k_2 \in \mathbb{Z}^2 \setminus\{0\}} \sum_{| i - j|\leqslant 1} \sum_{| i' - j'| \leqslant 1} 1_{m \lesssim i} 1_{m \lesssim i'} \rho^2_m(k_1 + k_2) \rho_i ( k_1 ) \rho_j ( k_2 )\rho_{i'} ( k_1 ) \rho_{j'} ( k_2 ) e^{- 2 t |k_1 |^2}\\
      &\hspace{30pt} + \sum_{k_1, k_2 \in \mathbb{Z}^2 \setminus\{0\}} \sum_{| i - j | \leqslant 1} \sum_{| i' - j' | \leqslant 1} 1_{m \lesssim i} 1_{m \lesssim i'} \rho^2_m (k_1 + k_2 ) \rho_i ( k_1 ) \rho_j ( k_2 ) \rho_{i'} ( k_2 ) \rho_{j'} ( k_1 ) e^{- t |k_1 |^2 - t | k_2 |^2}.
   \end{align*}
   There exists $c > 0$ such that $e^{- 2 t | k|^2} \lesssim e^{- t c 2^{2 i}}$ for all $k \in \tmop{supp} ( \rho_i )$ and for all $i \ge -1$. In the  remainder of the proof the value of this strictly positive $c$ may change from line to line. If $| i - j | \leqslant 1$, then we also have $e^{- t | k |^2} \lesssim e^{- t c 2^{2 i}}$ for all $k \in \tmop{supp} ( \rho_j )$. Thus
   \begin{align}\label{e:anderson area var}\nonumber
      & \tmop{Var} ( \Delta_m ( \pi_{\circ} ( P_t \xi, \xi )) ( x ) )\\ \nonumber
      &\hspace{40pt} \lesssim \sum_{i, j, i', j'} 1_{m \lesssim i} 1_{i \sim j \sim i' \sim j'} \sum_{k_1, k_2 \in \mathbb{Z}^2 \setminus\{0\}} 1_{\tmop{supp}(\rho_m)}(k_1 + k_2) 1_{\tmop{supp}( \rho_i )}(k_1) 1_{\tmop{supp} ( \rho_j)}(k_2) e^{- 2 t c 2^{2 i}} \\
      &\hspace{40pt} \lesssim \sum_{i:i \gtrsim m} 2^{2 i} 2^{2 m} e^{- t c 2^{2 i}} \lesssim \frac{2^{2 m}}{t} \sum_{i: i \gtrsim m} e^{- t c 2^{2 i}} \lesssim \frac{2^{2 m}}{t} e^{- t c 2^{2 m}},
   \end{align}
   where we used that $t 2^{2 i} \lesssim e^{t ( c - c' ) 2^{2 i}}$ for any $c' < c$.
   
   Now let $\alpha < 1$. We apply Jensen's inequality and combine \eqref{e:anderson area hyper}, \eqref{e:anderson area m-block}, \eqref{e:anderson area -1-block}, and \eqref{e:anderson area var} to obtain
   \begin{gather*}
      \E[\Vert  \Xi_t \Vert_{B^{2 \alpha - 2}_{2 p, 2 p}}] \lesssim \biggl( \sum_{m \ge - 1} 2^{(2 \alpha - 2) m 2 p} \E[\Vert  \Delta_m \Xi_t \Vert_{L^{2 p} (\mathbb{T}^2 )}^{2 p} ] \biggr)^{\frac{1}{2 p}} \\
      \lesssim t^{- 1 / 2} \biggl( \sum_{m \ge - 1} 2^{(2 \alpha - 2) m 2 p} 2^{2 m p} e^{- t c p 2^{2 m}} \biggr)^{\frac{1}{2 p}}  \lesssim t^{- 1 / 2} \biggl( \int_{- 1}^{\infty} ( 2^x )^{2 p ( 2 \alpha - 1 )} e^{- c t p ( 2^x )^2} \mathd x \biggr)^{\frac{1}{2 p}}.
   \end{gather*}
   The change of variables $y = \sqrt{t} 2^x$ then yields
   \begin{align*}
      \E[\Vert  \Xi_t \Vert_{B^{2 \alpha - 2}_{2 p, 2 p}}] \lesssim t^{- 1 / 2} \biggl( t^{- p ( 2 \alpha -1)} \int_0^{\infty} y^{2 p ( 2 \alpha - 1 ) - 1} e^{- c p y^2} \mathd y \biggr)^{\frac{1}{2 p}}.
   \end{align*}
   If $\alpha > 1/2$, then the integral is finite for all sufficiently large $p$, and therefore $\E[\Vert  \Xi_t \Vert_{B^{2 \alpha - 2}_{2 p, 2 p}}] \lesssim_p t^{- \alpha}$, so that $\int_0^1 \E[\Vert  \Xi_t \Vert_{B^{2 \alpha - 2}_{2 p, 2 p}}] \mathd t < \infty$ for all $\alpha < 1$. Gaussian hypercontractivity now allows us to conclude that also
   \[
      \E \Big[\Big\| \int_0^1 \Xi_t \mathd t\Big\|^p_{B^{2\alpha - 2}_{2p, 2p}} \Big] < \infty
   \]
   for all $p \ge 1$. The result now follows from Besov embedding, Lemma~\ref{l:besov embedding torus}.
\end{proof}

\begin{corollary}
   Let $\tilde{\xi}$ be a spatial white noise, and let $\xi$, $\vartheta$, and $\vartheta \diamond \xi$ be as defined above. Set $\tilde{\vartheta}(t,x)= \vartheta(x) + t (2\pi)^{-1} \CF \tilde{\xi}(0)$ and
   \[
      (\tilde{\vartheta}\diamond \tilde{\xi}) (t,x) = (\vartheta \diamond \xi) (t,x) + (2\pi)^{-1} (\vartheta \circ  \CF \tilde{\xi}(0))(x) + t (2\pi)^{-1} (\CF \tilde{\xi}(0) \circ \xi)(x) + t (2\pi)^{-2} (\CF\tilde{\xi}(0))^2,
   \]
   so that $L \tilde{\vartheta} = \tilde{\xi}$. If $\psi \in \CS$ satisfies $\int \psi(x) \mathd x = 1$, and if $\tilde{\xi}^\varepsilon = \psi^\varepsilon \ast \tilde{\xi}$ and $\tilde{\vartheta}^\varepsilon = \psi^\varepsilon \ast \tilde{\vartheta}$, where $\psi^\varepsilon = \varepsilon^{-2} \psi(\varepsilon \cdot)$ for $\varepsilon > 0$, then $\lim_{\varepsilon \rightarrow 0} \E[ \| \tilde{\vartheta} \circ \tilde{\xi} - ( \tilde{\vartheta}^\varepsilon \circ \tilde{\xi}^\varepsilon - c_\varepsilon) \|_{\CC^{2\alpha - 2}(\mathbb{T}^2)}^p] = 0$ for all $p \ge 1$, where $c_\varepsilon$ is as defined in Lemma~\ref{l:pam area}.
%
%   
%   
%    Then Since the smoothing operator $v \mapsto v^\varepsilon = \varepsilon^{-2} \psi(\varepsilon^{-1} \cdot) \ast v$ is linear, all the convergence results below stay true if we replace $\xi$ by $\tilde{\xi}$ and $(\xi^\varepsilon)$ by $(\tilde{\xi^\varepsilon})$. {\color{red} Not optimal formulation. Come back later to rephrase.}
\end{corollary}

Combining the existence of the renormalized product $\vartheta \diamond \xi$ with Theorem~\ref{theorem:pam existence uniqueness}, we obtain the existence and uniqueness of local solutions to the generalized parabolic Anderson model:

\begin{corollary}
  Let $\alpha \in ( 2/3, 1)$, $G \in C^3_b$, $u_0 \in \CC^\alpha(\mathbb{T})$, $L = \partial_t + (-\Delta)^\sigma$, and let $\xi$ be a spatial white noise on $\mathbb{T}^2$. Then there exists a unique local solution $u$ to
  \[
     L u = F(u)\diamond \xi, \qquad u(0) = u_0,
  \]
  in the following sense: For $\psi \in \CS$ with $\int \psi \mathd t=1$ and for $\varepsilon > 0$ consider the solution $u^\varepsilon$ to
  \[
     L u^{\varepsilon} =F ( u^{\varepsilon} ) \xi^{\varepsilon} -c_{\varepsilon} F' ( u^{\varepsilon} ) F ( u^{\varepsilon} ) , \hspace{2em} u^{\varepsilon} ( 0 ) =u_{0},
  \]
  on $[0, \infty)\times\mathbb{T}^2$, where $\xi^\varepsilon = \varepsilon^{-1} \psi(\varepsilon \cdot) \ast \xi$, and where $c_\varepsilon$ is as defined in Lemma~\ref{l:pam area}. Then there exists a strictly positive, $\sigma(u_0, \xi)$--measurable random time $\tau$, such that $\lim_{\varepsilon \rightarrow 0} \E[\| u - u^\varepsilon \|_{C_\tau \CC^\alpha(\mathbb{T}^2)}^p] = 0$ for all $p \ge 1$.
\end{corollary}




\begin{remark}


Consider the problem of convergence of $( \vartheta^{\varepsilon} \circ
\xi^{\varepsilon} )$. Since $L  \vartheta^{\varepsilon} = \xi^{\varepsilon}$
we have
\[ ( \vartheta^{\varepsilon} \circ \xi^{\varepsilon} ) = (
   \vartheta^{\varepsilon} \circ L  \vartheta^{\varepsilon} ) = \frac{1}{2} L 
   ( \vartheta^{\varepsilon} \circ   \vartheta^{\varepsilon} ) - (
   \mathD_{x} \vartheta^{\varepsilon} \circ \mathD_{x}  
   \vartheta^{\varepsilon} ),
\]
from which we see that the only problem in passing to the limit is given by
the second term on the right hand side. This integration by parts formula is the crucial difference with what happens in the \textsc{rde} case, which otherwise shares many structural properties with the \textsc{pam} model. The fact that $L$ is a second order operator generates the term $(\mathD_{x} \vartheta^{\varepsilon} \circ \mathD_{x}  \vartheta^{\varepsilon} )$ in the above computation, which is absent in case of the operator $\partial_t$. This term, whose convergence is equivalent to the convergence of the positive term $(\mathD_x \vartheta^\varepsilon)^2$, cannot have simple cancellation properties and it is the origin for the need of introducing an additive renomalization when considering the \textsc{pam} model.

Our previous analysis easily implies  that the solutions to the modified
problem
\[
   L u^{\varepsilon} =F ( u^{\varepsilon} ) \xi^{\varepsilon} + F'( u^{\varepsilon} ) F (u^{\varepsilon} ) ( \mathD_{x} \vartheta^{\varepsilon}  )^2
\]
will converge as soon as $\xi^{\varepsilon} \rightarrow \xi$ in $\CC^{\alpha
-2} ( \mathbb{T}^{2} )$, without any requirements on the bilinear quantity $\vartheta^{\varepsilon} \circ
\xi^{\varepsilon}$ . 

\end{remark}





\appendix


\section{Besov spaces and paraproducts}\label{sec:para}

\subsection{Littlewood-Paley theory and Besov spaces}\label{sec:littlewood}

In the following, we describe the concepts from Littlewood--Paley theory which are necessary for our analysis, and we recall the definition and some properties of Besov spaces. For a general introduction to Littlewood--Paley theory, Besov spaces, and paraproducts, we refer to the nice book of Bahouri, Chemin, and Danchin {\cite{Bahouri2011}}.

Littlewood--Paley theory allows for an efficient way of characterizing the regularity of functions and distributions. It relies on the decomposition of an arbitrary distribution into a series of smooth functions whose Fourier transforms have localized support.

Let $\chi \nocomma , \rho \in \CD$ be nonnegative radial functions on $\mathbb{R}^{d}$, such that
\begin{enumerateroman}
  \item the support of $\chi$ is contained in a ball and the support of $\rho$ is contained in an annulus;
  
  \item $\chi ( z ) + \sum_{j \geqslant 0} \rho ( 2^{-j} z ) =1$ for all $z \in \mathbb{R}^{d}$;
  
  \item \label{def:dyadic partition3}$\tmop{supp} ( \chi ) \cap \tmop{supp} (\rho ( 2^{-j} \cdummy ) ) = \emptyset$ for $j \geqslant 1$ and $\tmop{supp}( \rho ( 2^{-i} \cdummy ) ) \cap \tmop{supp} ( \rho ( 2^{-j} \cdummy ) ) =  \emptyset$ for $| i-j | >1$.
\end{enumerateroman}
We call such $( \chi , \rho )$ \tmtextit{dyadic partition of unity}, and we frequently employ the notation
\[
   \rho_{-1} = \chi \qquad \text{and} \qquad \rho_{j} = \rho (2^{-j} \cdummy ) \text{ for } j \geqslant 0.
\]
For the existence of dyadic partitions of unity see {\cite{Bahouri2011}}, Proposition 2.10. The Littlewood--Paley blocks are now defined as
\[
   \Delta_{-1} u= \CF^{-1} \left( \chi \CF u \right) = \CF^{-1} \left( \rho_{-1} \CF u \right) \qquad \text{and} \qquad \Delta_{j} u= \CF^{-1} \left( \rho_{j} \CF u \right) \text{ for }  j \geqslant 0. 
\]
Then $\Delta_{j} u=K_{j} \ast u$, where $K_{j} = \CF^{-1} \rho_{j}$ for all $j \geqslant -1$. In particular, $\Delta_{j} u$ is a smooth function for every $j \geqslant -1$. We also use the notation
\[
   S_{j} f= \sum_{i \leqslant j-1} \Delta_{i} f.
\]
It is easy to see that $u =  \sum_{j \geqslant -1} \Delta_{j} u= \lim_{j \rightarrow \infty} S_{j} u$ for every $u \in \CS'$.

For $\alpha \in \mathbb{R}$, the H{\"o}lder-Besov space $\CC^{\alpha}$ is given by $\CC^{\alpha} =B^{\alpha}_{\infty , \infty} (\mathbb{R}^{d},\mathbb{R}^{n} )$, where for $p,q \in [1, \infty ]$ we define
\[
   B^{\alpha}_{p,q} (\mathbb{R}^{d} ,\mathbb{R}^{n} ) = \bigg\{ u \in \CS'(\mathbb{R}^{d} ,\mathbb{R}^{n} ):\|u\|_{B^{\alpha}_{p,q}} = \bigg(\sum_{j \geqslant -1} (2^{j \alpha} \| \Delta_{j} u\|_{L^{p}} )^{q} \bigg)^{\frac{1}{q}} < \infty \bigg\},
\]
with the usual interpretation as $\ell^{\infty}$ norm in case $q= \infty$. The $\lVert \cdummy \rVert_{L^{p}}$ norm is taken with respect to Lebesgue measure on $\mathbb{R}^{d}$. While the norm $\lVert \cdummy \rVert_{B^{\alpha}_{p,q}}$ depends on the dyadic partition of unity $( \chi , \rho )$, the space $B^{\alpha}_{p,q}$ does not, and any other dyadic partition of unity corresponds to an equivalent norm. We write $\lVert \cdummy \rVert_{\alpha}$ instead of $\lVert \cdummy \rVert_{B^{\alpha}_{\infty , \infty}}$.

We denote by $\CC^{\alpha}_{\text{\tmop{loc}}}$ the set of all distributions $u$ such that $\varphi u \in \CC^{\alpha}$ for every $\varphi \in \CD$. If for every $\varphi \in \CD$, the difference $\varphi(u_n - u)$ converges to $0$ in $\CC^\alpha$, then we say that $(u_n)$ converges to $u$ in $\CC^\alpha_{\tmop{loc}}$.

If $\alpha \in (0, \infty ) \backslash \mathbb{N}$, then $\CC^{\alpha}$ is the space of $\lfloor \alpha \rfloor$ times differentiable functions, whose partial derivatives up to order $\lfloor \alpha \rfloor$ are bounded, and whose partial derivatives of order $\lfloor \alpha \rfloor$ are ($\alpha - \lfloor \alpha \rfloor$)-H{\"o}lder continuous (see p. 99 of {\cite{Bahouri2011}}). Note however, that for $k \in \mathbb{N}$ the H{\"o}lder-Besov space $\CC^{k}$ is strictly larger than $C^{k}_{b}$.

We will use without comment that $\lVert \cdummy \rVert_{\alpha} \leqslant \lVert \cdummy \rVert_{\beta}$ for $\alpha \leqslant \beta$, that $\lVert \cdummy \rVert_{L^{\infty}} \lesssim \lVert \cdummy \rVert_{\alpha}$ for $\alpha >0$, and that $\lVert \cdummy \rVert_{\alpha} \lesssim \lVert \cdummy \rVert_{L^{\infty}}$ for $\alpha \leqslant 0$. We will also use that $\|S_{j} u\|_{L^{\infty}} \lesssim 2^{j \alpha} \|u\|_{\alpha}$ for $\alpha <0$ and $u \in C^{\alpha}$.


The following Bernstein inequalities are often useful when dealing with functions with compactly supported Fourier transform.

\begin{lemma}[Lemma 2.1 of {\cite{Bahouri2011}}]\label{lemma:Bernstein}
   Let $\CA$ be an annulus and let $\CB$ be a ball. For any $k \in \mathbb{N}$, $\lambda >0$, and $1 \leqslant p \leqslant q \leqslant \infty$ we have that
   \begin{enumeratenumeric}
      \item if $u \in L^{p} ( \mathbb{R}^{d} )$ is such that $\tmop{supp} ( \CF u ) \subseteq \lambda \CB$, then
         \[
            \max_{\mu \in \mathbb{N}^{d} : | \mu | \leqslant k} \| \partial^{\mu} u \|_{L^{q}} \lesssim_{k} \lambda^{k+d \left( \frac{1}{p} - \frac{1}{q} \right)} \| u \|_{L^{p}};
         \]
      \item if $u \in L^{p} ( \mathbb{R}^{d} )$ is such that $\tmop{supp} (\CF u ) \subseteq \lambda \CA$, then
         \[
            \lambda^{k} \| u \|_{L^{p}} \lesssim_{k} \max_{\mu \in \mathbb{N}^{d}: | \mu | \leqslant k} \| \partial^{\mu} u \|_{L^{p}} .
         \]
  \end{enumeratenumeric}
\end{lemma}

For example, it is a simple consequence of the Bernstein inequalities that $\| \mathD^k u \|_{\alpha-k} \lesssim \| u \|_\alpha$ for all $\alpha \in \R$ and $k \in \N$.

We point out that everything above and everything that follows can (and will) be applied to distributions on the torus. More precisely, let $\CD'(\mathbb{T}^{d} )$ be the space of distributions on $\mathbb{T}^{d}$. Any $u \in \CD' (\mathbb{T}^{d} )$ can be interpreted as a periodic tempered distribution on $\mathbb{R}^{d}$, with frequency spectrum contained in $\mathbb{Z}^{d}$ -- and vice versa. For details see~\cite{Schmeisser1987}, Chapter 3.2. In particular, $\Delta_{j} u$ is a periodic smooth function, and therefore $\| \Delta_{j} u\|_{L^{\infty}} = \| \Delta_{j} u\|_{L^{\infty} (\mathbb{T}^{d} )}$. In other words, we can define
\[
    \CC^\alpha(\mathbb{T}^d) = \{ u \in \CC^\alpha: u \text{ is } (2\pi)-\text{periodic}\}
\]
for $\alpha \in \R$. For $p \neq \infty$ however, this definition is not very useful, because no nontrivial periodic function is in $L^p$ for $p < \infty$. Therefore, Besov spaces on the torus with general indices $p,q \in [1,\infty]$ are defined as
\[
   B^\alpha_{p,q}(\mathbb{T}^d) := \biggl\{ u \in \mathcal{D}' (\mathbb{T}^d) : \lVert u\rVert_{B^{\alpha}_{p, q}(\mathbb{T}^d)} := \biggl( \sum_{j \ge - 1} (2^{j \alpha} \lVert \Delta_j u \rVert_{L^p(\mathbb{T}^d)})^q\biggr)^{\frac{1}{q}} < \infty \biggr\}.
\]
Note that the Fourier transform of $u \in \mathcal{D}'(\mathbb{T}^d)$ is supported in $\Z^d$, and that
\[
   u(x) = (2\pi)^{-d} \sum_{k \in \Z^d} \CF{u}(k) e^{\imath \langle k, x\rangle} = \CF^{-1} (\CF{u})(x).
\]
Apart from that, $\Delta_j u = \CF^{-1} (\rho_j \CF u)$ is defined exactly as in the non-periodic case. Note that we have $\CC^\alpha(\mathbb{T}^d) = B^\alpha_{\infty,\infty} (\mathbb{T}^d)$. Strictly speaking we will not work with $B^\alpha_{p,q}(\mathbb{T}^d)$ for $(p,q) \neq (\infty, \infty)$. But we will need the Besov embedding theorem on the torus:
\begin{lemma}\label{l:besov embedding torus}
   Let $1\leqslant p_1 \leqslant p_2 \leqslant \infty$ and $1\leqslant q_1 \leqslant q_2 \leqslant \infty$, and let $\alpha \in \R$. Then $B^{\alpha}_{p_1,q_1}(\mathbb{T}^d)$ is continuously embedded in $B^{\alpha - d(1/p_1 - 1/p_2)}_{p_2,q_2}(\mathbb{T}^d)$, and $B^{\alpha}_{p_1,q_1}(\mathbb{R}^d)$ is continuously embedded in $B^{\alpha - d(1/p_1 - 1/p_2)}_{p_2,q_2}(\mathbb{R}^d)$.
\end{lemma}
For the embedding theorem on $\R^d$ see~\cite{Bahouri2011}, Proposition~2.71. The result on the torus can be shown using the same arguments, see for example~\cite{Chemin2006}. In both cases, the proof is based on the Bernstein inequalities, Lemma~\ref{lemma:Bernstein}.


The following characterization of Besov regularity for functions which can be decomposed into pieces that are well localized in Fourier space will be useful below.

\begin{lemma}\label{lem: Besov regularity of series}(Lemma 2.69 and 2.84 of {\cite{Bahouri2011}})
   \begin{enumeratenumeric}
      \item Let $\CA$ be an annulus, let $\alpha \in \mathbb{R}$, and let $(u_{j})$ be a sequence of smooth functions such that $\CF u_{j}$ has its support in $2^{j} \CA$, and such that $\| u_{j} \|_{L^{\infty}} \lesssim 2^{-j \alpha}$ for all $j$. Then
         \[
            u= \sum_{j \geqslant -1} u_{j} \in \CC^{\alpha} \qquad  \tmop{and} \qquad \| u \|_{\alpha} \lesssim \sup_{j \geqslant -1} \{ 2^{j \alpha} \| u_{j} \|_{L^{\infty}} \} .
         \]
      \item Let $\CB$ be a ball, let $\alpha >0$, and let $( u_{j} )$ be a sequence of smooth functions such that $\CF u_{j}$ has its support in $2^{j} \CB$, and such that $\| u_{j} \|_{L^{\infty}} \lesssim 2^{-j \alpha}$ for all $j$. Then
         \[
            u= \sum_{j \geqslant -1} u_{j} \in \CC^{\alpha} \qquad \tmop{and} \qquad \| u \|_{\alpha} \lesssim \sup_{j \geqslant -1} \{ 2^{j \alpha} \| u_{j} \|_{L^{\infty}} \} .
         \]
   \end{enumeratenumeric}
\end{lemma}

\begin{proof}
  Let us first consider the case when $\CF u_{j}$ is supported in $2^{j} \CA$. This implies that $\Delta_{i} u_{j} \neq 0$ only for $i \sim j$. Hence, we obtain
  \begin{gather*}
     \| \Delta_{i} u \|_{L^{\infty}} \leqslant \sum_{j \geqslant -1} \| \Delta_{i} u_{j} \|_{L^{\infty}} = \sum_{j:j \sim i} \| \Delta_{i} u_{j} \|_{L^{\infty}} \leqslant \sup_{k \geqslant -1} \{ 2^{k \alpha} \| u_{k} \|_{L^{\infty}} \} \sum_{j:j \sim i} 2^{-j  \alpha} \\
     \simeq \sup_{k \geqslant -1} \{ 2^{k \alpha} \| u_{k} \|_{L^{\infty}} \} 2^{-i \alpha}.
  \end{gather*}
  
   If $\CF u_{j}$ is supported in $2^{j} \CB$, then $\Delta_{i} u_{j} \neq 0$ only for $i \lesssim j$. Therefore,
   \[
      \| \Delta_{i} u \|_{L^{\infty}} \leqslant \sum_{j:j \gtrsim i} \| \Delta_{i} u_{j} \|_{L^{\infty}} \leqslant \sup_{k \geqslant -1} \{ 2^{k \alpha} \| u_{k} \|_{L^{\infty}} \} \sum_{j:j \gtrsim i} 2^{-j  \alpha}  \lesssim \sup_{k \geqslant -1} \{ 2^{k \alpha} \| u_{k} \|_{L^{\infty}} \} 2^{-i \alpha},
   \]
   using $\alpha >0$ in the last step.
\end{proof}

Sometimes it will be also useful to have some control of the Besov norm of a function which is obtained as a sum of regular functions without any Fourier localization. In this case the following lemma applies.

\begin{lemma}\label{lemma:collage}
   Let $( u_{j} )_{j \geqslant -1}$ a sequence of functions on $\mathbb{R}^{d}$ for which there exist constants $c, \alpha >0$, and an integer $n> \alpha$, such that
   \[
      \| \mathD^{n} u_{j} \|_{L^{\infty}} \leqslant 2^{( n- \alpha ) j} c,  \hspace{2em} \| u_{j} \|_{L^{\infty}} \leqslant 2^{- \alpha j} c.
   \]
   Then the series $u= \sum_{j \geqslant -1} u_{j}$ converges in $L^{\infty}$, and we have $\| u \|_{\alpha} \lesssim c$.
\end{lemma}

\begin{proof}
   Start by noting that
   \[
      \| \Delta_{k} u_{j} \|_{L^{\infty} ( \mathbb{R}^{d} )} \lesssim \| u_{j} \|_{L^{\infty} ( \mathbb{R}^{d} )} \lesssim 2^{- \alpha j} c,
   \]
   and that by Bernstein inequalities we have for $k \geqslant 0$
   \[
      \| \Delta_{k} u_{j} \|_{L^{\infty} ( \mathbb{R}^{d} )} \lesssim 2^{-k n} \| \mathD^{n} \Delta_{k} u_{j} \|_{L^{\infty} ( \mathbb{R}^{d} )} \lesssim 2^{-k n} \| \mathD^{n} u_{j} \|_{L^{\infty} ( \mathbb{R}^{d} )} \lesssim 2^{-k n} 2^{ ( n- \alpha ) j} c.
   \]
   Then
   \[
      \sum_{j \geqslant -1} \| \Delta_{-1} u_{j} \|_{L^{\infty}} \lesssim \sum_{j \geqslant -1} 2^{- \alpha j} c \lesssim c.
   \]
   For $k \geqslant 0$ we can bound $\sum_{j \geqslant -1} \| \Delta_{k} u_{j} \|_{L^{\infty}}$ by splitting the sum over $j$ at $k$, and obtain
   \begin{gather*}
      \sum_{j \geqslant -1} \| \Delta_{k} u_{j} \|_{L^{\infty}} = \sum_{-1 \leqslant j \leqslant k} \| \Delta_{k} u_{j} \|_{L^{\infty} (\mathbb{R}^{d} )} + \sum_{j>k} \| \Delta_{k} u_{j} \|_{L^{\infty} (\mathbb{R}^{d})} \lesssim \sum_{-1 \leqslant j \leqslant k} 2^{-k n} 2^{ ( n- \alpha ) j} c+ \sum_{j>k} 2^{- \alpha j} c \\
      \lesssim 2^{-k n} 2^{ ( n- \alpha ) k} c+2^{- \alpha k} c \lesssim c2^{-\alpha k}.
   \end{gather*}
   This bound implies that the series $\Delta_{k} u= \sum_{j \geqslant -1} \Delta_{k} u_{j}$ converges in $L^{\infty} ( \mathbb{R}^{d} )$ for all $k \geqslant -1$, and thus defines a distribution $u= \sum_{k} \Delta_{k} u \in \CS'$. Moreover, we have $\| \Delta_{k} u \|_{L^{\infty}} \lesssim c2^{-k\alpha}$ for all $k \geqslant -1$, which implies that $u \in \CC^{\alpha}$ and that $\| u \|_{\alpha} \lesssim c$.
\end{proof}

\subsection{Linear operators acting on Besov spaces}

Now that we discussed Besov spaces, their properties, and products of tempered distributions, let us study the action of some important linear operators on Besov spaces.

First we consider the scaling of the spatial variable:

\begin{lemma}\label{lemma:scaling}
  For all $\lambda >0$ and $u \in \CS'$ define the scaling transformation $\Lambda_{\lambda} u ( \cdot ) =u ( \lambda \cdot )$. Then we have
  \[
     \| \Lambda_{\lambda} u \|_{\alpha} \lesssim ( 1+ \lambda^{\alpha} ) \| u \|_{\alpha}
  \]
  for all $\alpha \in \mathbb{R} \setminus \{ 0 \}$ and all $u \in \CC^{\alpha}$.
\end{lemma}

\begin{proof}
  Let $u \in \CC^{\alpha}$ and let $\Lambda_{\lambda} u ( x ) =u ( \lambda x )$ for some $\lambda >0$. Note that $\Lambda_{\lambda} \mathD =  \lambda^{-1} \mathD \Lambda_{\lambda}$, and therefore $\Lambda_{\lambda} \Delta_{j} u= \Lambda_{\lambda} \rho ( 2^{-j} \mathD ) u= \rho ( 2^{-j} \lambda^{-1} \mathD ) \Lambda_{\lambda} u$, which implies that the Fourier transform of $\Lambda_{\lambda} \Delta_{j} u$ is supported in the annulus $\lambda 2^{j} \CA$ (where $\CA$ is the annulus in which $\rho$ is supported). In particular, if $k \geqslant 0$, we have $\Delta_{k} \Lambda_{\lambda} \Delta_{j} u \neq 0$ only if $2^{k} \sim \lambda 2^{j}$. Thus, there exist $a,b>0$ such that for $k \geqslant 0$ we have
  \begin{align*}
     \| \Delta_{k} \Lambda_{\lambda} u \|_{L^{\infty}} & \lesssim \sum_{j:a2^{k} \leqslant \lambda 2^{j} \leqslant b2^{k}} \| \Delta_{k} \Lambda_{\lambda} \Delta_{j} u \|_{L^{\infty}} \lesssim \sum_{j:a2^{k} \leqslant \lambda 2^{j} \leqslant b2^{k}} \| \Delta_{j} u \|_{L^{\infty}} \\
     & \lesssim \| u \|_{\alpha} \sum_{j:a2^{k} \leqslant \lambda 2^{j} \leqslant b2^{k}} 2^{- \alpha j} \lesssim \| u \|_{\alpha} \lambda^{\alpha} 2^{- \alpha k} .
  \end{align*}
  For $k=-1$ we can simply bound
  \[
     \| \Delta_{-1} \Lambda_{\lambda} u \|_{L^{\infty}} \lesssim \sum_{j: \lambda 2^{j} \lesssim 1} \| \Delta_{k} \Lambda_{\lambda} \Delta_{j} u \|_{L^{\infty}} \lesssim \| u \|_{\alpha} \sum_{j: \lambda 2^{j} \lesssim 1} 2^{- \alpha j} \lesssim \| u \|_{\alpha} ( 1+ \lambda^{\alpha} ),
  \]
  where we used $\alpha \neq 0$ in the last step. This completes the proof.
\end{proof}

The behavior of Besov spaces under the action of Fourier multipliers is the subject of the next lemma.

\begin{lemma}\label{lemma:fourier multiplier smoothing}
   Let $\alpha \in \mathbb{R}$. Let $\varphi$ be a continuous function, such that $\varphi$ is infinitely differentiable everywhere except possibly at 0, and such that $\varphi$ and all its partial derivatives decay faster than any rational function at infinity. Assume also that $\CF \varphi \in L^{1}$. Then we have for every $\varepsilon \in (0,1]$, for every $\delta \geqslant 0$, and for every $u \in \CS'$ that
  \[
     \| \varphi ( \varepsilon \mathD ) u \|_{\alpha + \delta} \lesssim_{\delta, \varphi} \varepsilon^{- \delta} \| u \|_{\alpha} \hspace{2em} \tmop{and} \hspace{2em} \| \varphi ( \varepsilon \mathD ) u \|_{\delta}  \lesssim_{\delta , \varphi} \varepsilon^{- \delta} \| u \|_{L^{\infty}} .
  \]
\end{lemma}

\begin{proof}
  Let $\psi \in \CD$ with support in an annulus be such that $\psi \rho = \rho$, where $( \chi , \rho )$ is our dyadic partition of unity. Then we have for $j \geqslant 0$ that
  \[
     \varphi ( \varepsilon \mathD ) \Delta_{j} u= \left[ \CF^{-1} ( \varphi (\varepsilon \cdummy ) \psi ( 2^{-j} \cdummy ) ) \right] \ast \Delta_{j}  u,
  \]
  and therefore Young's inequality implies
  \[
     \| \varphi ( \varepsilon \mathD ) \Delta_{j} u \|_{L^{\infty}} \lesssim \left\| \CF^{-1} ( \varphi ( \varepsilon \cdummy ) \psi ( 2^{-j} \cdummy ) ) \right\|_{L^{1}} 2^{-j \alpha} \| u \|_{\alpha} = \left\| \CF^{-1} ( \varphi ( 2^{j} \varepsilon \cdummy ) \psi ) \right\|_{L^{1}} 2^{-j \alpha} \| u \|_{\alpha}.
  \]
  Hence, it suffices to prove $\left\| \CF^{-1} ( \varphi ( 2^{j} \varepsilon \cdummy ) \psi ) \right\|_{L^{1}} \lesssim \varepsilon^{- \delta} 2^{-j \delta}$. But
  \begin{gather*}
     \left\| \CF^{-1} ( \varphi ( 2^{j} \varepsilon \cdummy ) \psi ) \right\|_{L^{1}} \lesssim \left\| ( 1+ | \cdot |^{2} )^{d} \CF^{-1} ( \varphi ( 2^{j} \varepsilon \cdummy ) \psi ) \right\|_{L^{\infty}}     \lesssim \left\| \CF^{-1} ( ( 1+ \Delta )^{d} ( \varphi ( 2^{j} \varepsilon \cdummy ) \psi ) ) \right\|_{L^{\infty}} \\
      \lesssim \| ( 1+ \Delta )^{d} ( \varphi ( 2^{j} \varepsilon \cdummy ) \psi ) \|_{L^{1}} \lesssim ( 1+2^{j} \varepsilon )^{2d} \max_{\mu \in \mathbb{N}^{d} : | \mu | \leqslant 2d} \| \partial^{\mu} \varphi ( 2^{j}  \varepsilon \cdot ) \|_{L^{\infty} ( \tmop{supp} ( \psi ) )} .
  \end{gather*}
  By assumption, $\varphi$ is smooth away from 0, and $\varphi$ and all its partial derivatives decay faster than any rational function at infinity. Thus, for every $\mu \in \mathbb{N}^{d}$ and every $\delta \in \mathbb{R}$ there exists a constant $C_{\mu , \delta} >0$ such that
  \[
     \sup_{x \geqslant 1} ( 1+ | x | )^{\delta +2d} | \partial^{\mu} \varphi ( x ) | \leqslant C_{\mu , \delta}.
  \]
  Since $\tmop{supp} ( \psi )$ is bounded away from $0$, there exists a minimal $j_{0} \in \mathbb{N}$, such that $2^{j_{0}} \varepsilon | x | \geqslant 1$ for all $x \in \tmop{supp} ( \psi )$. Thus, we obtain
  \[
     \left\| \CF^{-1} ( \varphi ( 2^{j} \varepsilon \cdummy ) \psi ) \right\|_{L^{1}} \lesssim \Big(\max_{\mu \in \mathbb{N}^{d} : | \mu | \leqslant 2d} C_{\mu, \delta}\Big) ( 1+2^{j} \varepsilon )^{2d} ( 1+2^{j}\varepsilon )^{- \delta -2d} \lesssim_{\delta , \varphi} ( 1+2^{j} \varepsilon )^{- \delta} \lesssim 2^{-j \delta} \varepsilon^{- \delta}
  \]
  for all $j \geqslant j_{0}$. On the other side, we have for $j \leqslant j_{0}$ that
  \begin{align*}
     \| \varphi ( \varepsilon \mathD ) \Delta_{j} u \|_{L^{\infty}} &\lesssim \| \CF^{-1} ( \varphi ( \varepsilon \cdummy ) ) \|_{L^{1}} \| \Delta_{j} u \|_{L^{\infty}} \lesssim 2^{-j \alpha} \| u \|_{\alpha} = ( \varepsilon 2^{j} )^{\delta} \varepsilon^{- \delta} 2^{-j ( \alpha +  \delta )} \| u \|_{\alpha} \\
     & \leqslant ( \varepsilon 2^{j_{0}} )^{\delta} \varepsilon^{- \delta} 2^{-j ( \alpha + \delta )} \| u \|_{\alpha} \lesssim_{\delta , \varphi} \varepsilon^{- \delta} 2^{-j ( \alpha + \delta )} \| u \|_{\alpha},
  \end{align*}
  where we used that $\delta \geqslant 0$.
  
  The estimate for $v \in L^{\infty}$ follows from the same arguments.
\end{proof}

\begin{remark}
   If the support of $\CF u$ has a ``hole'' at 0, that is if there exists a ball $\CB$ centered at 0 such that $\CF u$ is supported outside of $\CB$, then the estimates of Lemma~\ref{lemma:fourier multiplier smoothing} hold for all $\varepsilon > 0$ and not just for $\varepsilon \in (0,1]$. This is an immediate consequence of the proof above.
\end{remark}

As an application, we derive the smoothing properties of the heat kernel generated by the fractional Laplacian.

\begin{lemma}\label{lemma:heat flow smoothing}
   Let $\sigma \in ( 0,1 ]$, let $- ( - \Delta )^{\sigma}$ be the fractional Laplacian with periodic boundary conditions on $\mathbb{T}^{d} \nosymbol$, and let $( P_{t} )_{t \geqslant 0}$ be the semigroup generated by $- ( - \Delta )^{\sigma}$. Then for all $T>0$, $t \in (0,T]$, $\alpha \in \R$, $\delta \geqslant 0$, and $u \in \CS'$ we have
  \[
     \|P_{t} u\|_{\alpha + \delta} \lesssim_T t^{- \delta / ( 2 \sigma )} \|u\|_{\alpha} \hspace{2em} \tmop{and} \hspace{2em} \| P_{t} v \|_{\delta} \lesssim_T t^{- \delta / ( 2 \sigma )} \| v \|_{L^{\infty}} .
  \]
  If $\CF u$ is supported outside of a ball centered at 0, then these estimates are uniform in $t > 0$ and not just in $t \in (0,T]$.
\end{lemma}

\begin{proof}
  The semigroup is given by $P_{t} = \varphi ( t^{1/ ( 2 \sigma )} \mathD )$ with $\varphi ( z ) =e^{- | z |^{2 \sigma}}$. Now $\varphi$ and its derivatives decay faster than any rational function at $\infty$. Also, for $\sigma \le 1$, $\CF \varphi$ is the density of a symmetric $2 \sigma$-stable random variable, and therefore in $L^{1}$. For $\sigma > 1$ it is easily shown that $(1+|\cdot|^{d+1}) \CF \varphi$ is bounded, and therefore in $L^1$. Thus, the estimates follow from Lemma~\ref{lemma:fourier multiplier smoothing}.
\end{proof}

%\begin{remark}
%   The estimates of Lemma~\ref{lemma:heat flow smoothing} extends of course from $t \in [0,1]$ to $t \in [0,T]$ with $T>0$, only that then the implicit constants on the right hand side will depend on $T$.
%\end{remark}


\begin{lemma}\label{lemma:semigroup hoelder in time}
   Let $\sigma \in ( 0,1 ]$, let $- ( -\Delta )^{\sigma}$ be the fractional Laplacian with periodic boundary conditions on $\mathbb{T}^{d} \nosymbol$, and let $( P_{t} )_{t \geqslant 0}$ be the semigroup generated by $- ( - \Delta )^{\sigma}$. Let $\alpha \in \mathbb{R}$, $\beta \in [ 0,1 ]$, and let $u \in \CC^{\alpha}$. Then for any $t \geqslant 0$ we have
  \[
     \| ( P_{t} - \tmop{Id} ) u \|_{L^{\infty}} \lesssim t^{\beta / ( 2 \sigma )} \| u \|_{\beta}.
  \]
\end{lemma}

\begin{proof}
  For the uniform estimate of $( P_{t} - \tmop{Id} ) u$, we write $P_{t} - \tmop{Id}$ as convolution operator: if $\varphi ( z ) =e^{- | z |^{2\sigma}}$ and $K ( x ) =\mathcal{F}^{-1} \varphi$, then
  \begin{align*}
     | ( P_{t} - \tmop{Id} ) u ( x ) | & = \left| t^{-d/ ( 2 \sigma )} \int K \left( \frac{x-y}{t^{1/ ( 2 \sigma )}} \right) ( u ( y ) -u ( x ) ) \mathd y \right| \\
     & \lesssim t^{-d/ ( 2 \sigma )} \int K \left( \frac{x-y}{t^{1/ ( 2 \sigma )}} \right) | y-x |^{\beta} \| u \|_{\beta} \mathd y \lesssim t^{\beta / ( 2 \sigma )} \| u \|_{\beta},
  \end{align*}
  where we identified $\CC^{\beta}$ with the classical space of H{\"o}lder continuous functions.
\end{proof}

Based on Lemma~\ref{lemma:heat flow smoothing} and Lemma~\ref{lemma:semigroup hoelder in time}, we derive the following Schauder estimates:

\begin{lemma}\label{lemma:schauder}
  Assume that $v \in C_{T} \CC^{\beta}$ for some $\beta \in \mathbb{R}$ and $T>0$. Then letting $V ( t ) = \int_{0}^{t} P_{t-s} v ( s ) \mathd s$, we have
  \begin{equation}\label{eq:schauder-heat}
     t^{\gamma} \| V ( t ) \|_{\beta +2 \sigma} \lesssim \sup_{s \in [ 0,t ]} ( s^{\gamma} \| v ( s ) \|_{\beta} )
  \end{equation}
  for all $\gamma \in [ 0,1 )$ and all $t \in [ 0,T ]$. If $0< \beta +2 \sigma <1$, then we also have
  \begin{equation}\label{eq:holder-est-heat}
     \| V \|_{C_{T}^{( \beta +2 ) / ( 2 \sigma )} L^{\infty}} \lesssim \sup_{s \in [ 0,t ]} \| v ( s ) \|_{\beta} .
  \end{equation}
\end{lemma}

\begin{proof}
  Consider $\Delta_{q} V$ for some $q \geqslant 0$ and let $\delta \in [ 0,t/2 ]$. We decompose the integral into two parts:
  \[
     \Delta_{q} V ( t ) = \int_{0}^{t} P_{t-s} ( \Delta_{q} v ) ( s ) \mathd s= \int_{0}^{\delta} P_{s} ( \Delta_{q} v ) ( t-s ) \mathd s+  \int_{\delta}^{t} P_{s} ( \Delta_{q} v ) ( t-s ) \mathd s.
  \]
  Letting $M= \sup_{s \in [ 0,t ]} ( s^{\gamma} \| v ( s ) \|_{\beta} )$, we estimate the first term by
  \begin{align*}
     \left\| \int_{0}^{\delta} P_{s} ( \Delta_{q} v ) ( t-s ) \mathd s \right\|_{L^{\infty}} & \leqslant \int_{0}^{\delta} 2^{-q \beta} \| v ( t-s ) \|_{\beta} \mathd s \leqslant 2^{-q \beta} M \int_{0}^{\delta} ( t-s)^{- \gamma} \mathd s \\
     & =M2^{-q \beta} t^{1- \gamma} \int_{0}^{\delta /t} \frac{\mathd s}{( 1-s )^{\gamma}} \lesssim M2^{-q \beta} t^{- \gamma} \delta ,
  \end{align*}
  using $| 1- ( 1- \delta /t )^{1- \gamma} | \lesssim \delta /t$ in the last step. On the other side, we can use Lemma~\ref{lemma:heat flow smoothing} to
  estimate the second term for $\varepsilon >0$ by
  \begin{align*}
     &\left\| \int_{\delta}^{t} P_{s} ( \Delta_{q} v ) ( t-s ) \mathd s \right\|_{L^{\infty}} \lesssim \int_{\delta}^{t} s^{-1- \varepsilon} 2^{-q ( \beta +2 \sigma ( 1+ \varepsilon ) )} \| v ( t-s ) \|_{\beta} \mathd s \\
     &\hspace{80pt} \lesssim M2^{-q ( \beta +2 \sigma ( 1+ \varepsilon ) )} \int_{\delta}^{t} \frac{\mathd s}{s^{1+ \varepsilon} ( t-s )^{\gamma}} =M2^{-q ( \beta +2 \sigma ( 1+ \varepsilon ) )} t^{- \varepsilon - \gamma} \int_{\delta /t}^{1} \frac{\mathd s}{s^{1+ \varepsilon} ( 1-s )^{\gamma}} \\
     &\hspace{80pt} \lesssim M2^{-q ( \beta +2 \sigma ( 1+ \varepsilon ) )} t^{- \gamma} \delta^{- \varepsilon} =M2^{-q ( \beta +2 \sigma )} ( 2^{q2 \sigma} \delta  )^{- \varepsilon} t^{- \gamma} .
  \end{align*}
  If $2^{-q2 \sigma} \leqslant t/2$, we can take $\delta =2^{-q2 \sigma}$ to
  obtain $\| \Delta_{q} V ( t ) \|_{L^{\infty}} \lesssim Mt^{- \gamma} 2^{-q (\beta +2 \sigma )}$. If $2^{-q2 \sigma} >t/2$, we have $\| \Delta_{q} V ( t ) \|_{L^{\infty}} \lesssim M2^{-q \beta} t^{1- \gamma} \lesssim Mt^{- \gamma} 2^{-q ( \beta +2 \sigma )}$, and the first claim follows.
  
  As for the second claim, note that for $0 \leqslant s<t \leqslant T$ we have
  \[
      V ( t ) -V ( s ) = ( P_{t-s} - \tmop{Id} ) V ( s ) + \int_{s}^{t} P_{t-r} v ( r ) \mathd r,
  \]
  and therefore we can apply Lemma~\ref{lemma:semigroup hoelder in time} to obtain
  \begin{align*}
     \| V ( t ) &-V ( s ) \|_{L^{\infty}}  \lesssim \| ( P_{t-s} - \tmop{Id} ) V( s ) \|_{L^{\infty}} + \int_{s}^{t} \| P_{t-r} v ( r ) \|_{L^{\infty}} \mathd r \\
     & \lesssim | t-s |^{( \beta +2 ) / ( 2 \sigma )} \| V ( s ) \|_{\beta +2} +  \int_{s}^{t} \| v ( r ) \|_{\beta} \mathd r \lesssim_{T} | t-s |^{( \beta +2 ) / ( 2 \sigma )} \sup_{r \in [ 0,t ]} \| v ( r ) \|_{\beta} ,
  \end{align*}
  where we used that $( \beta +2 ) \in ( 0,1 )$ and that $| t-s | \leqslant T$. This yields the second claim.
\end{proof}

When dealing with \textsc{rde}s, the convolution with the (fractional) heat kernel has a natural correspondence in the integral map.

\begin{lemma}\label{lemma:integral}
  Let $u \in \CC^{\alpha -1} ( \mathbb{R} )$ for some $\alpha \in ( 0,1 )$. Then there exists a unique $U \in \CC^{\alpha}_{\tmop{loc}} ( \mathbb{R} )$ such that $\mathD U=u$ and $U ( 0
  ) =0$. This antiderivative $U$ satisfies
  \begin{equation}\label{eq:integral estimate 1}
     | U ( t ) -U ( s ) | \lesssim | t-s |^{\alpha} \| u \|_{\alpha -1}
  \end{equation}
  for all $s,t \in \mathbb{R}$ with $| s-t | \leqslant 1$.We will use the notation $U ( t ) = \int_{0}^{t} u ( s ) \mathd s$ to denote this map, which is an extension of the usual definite integral.
  
  If the support of $u$ is contained in $[ -T,T ]$ for some $T>0$, then $U \in \CC^{\alpha}$ and
  \[
     \| U \|_{\alpha} \lesssim T \| u \|_{\alpha -1}.
  \]
\end{lemma}

\begin{proof}
  The second statement about compactly supported $u$ follows from the first statement by identifying $\CC^{\alpha}$ with the space of bounded H{\"o}lder continuous functions.
  
  As for the first statement, we define
  \[
     U ( t ) = \sum_{j \geqslant -1} \int_{0}^{t} \Delta_{j} u ( s ) \mathd s.
  \]
  If we can show~{\eqref{eq:integral estimate 1}}, then $U$ is indeed in $\CC^{\alpha}_{\tmop{loc}}$ and therefore in particular in $\CS'$. Since the derivative $\mathD$ is a continuous operator on $\CS'$, we then conclude that $\mathD U= \sum_{j} \Delta_{j} u=u$. Let therefore $s,t \in \mathbb{R}$ with $| s-t | \leqslant 1$. We have
  \[
     \left| \int_{s}^{t} \Delta_{j} u ( r ) \mathd r \right| \lesssim 2^{j (1- \alpha )} \| u \|_{\alpha -1} | t-s | .
  \]
  If $j \geqslant 0$, then $\Delta_{j} u= \mathD \mathD^{-1} ( \Delta_{j} u)$, where $\mathD^{-1}$ is the Fourier multiplier with symbol $1/ ( \iota  z )$, and therefore
  \[
     \left| \int_{s}^{t} \Delta_{j} u ( r ) \mathd r \right| = | \mathD^{-1} \Delta_{j} u ( t ) - \mathD^{-1} \Delta_{j} u ( s ) | \lesssim 2^{-j} \| \Delta_{j} u \|_{L^{\infty}} \lesssim 2^{-j \alpha} \| u \|_{\alpha -1} ,
  \]
  where we used the Bernstein inequality, Lemma~\ref{lemma:Bernstein}. If $j_{0}$ is such that $2^{-j_{0}} \leqslant | t-s | <2^{-j_{0} +1}$, then we use the first estimate for $j \leqslant j_{0}$ and the second estimate for $j>j_{0}$, and obtain
  \begin{align*}
     | U ( t ) -U ( s ) | &\leqslant \sum_{j \geqslant -1} \left| \int_{s}^{t} \Delta_{j} u ( r ) \mathd r \right| \lesssim \sum_{j \leqslant j_{0}} 2^{j ( 1- \alpha )} \| u \|_{\alpha -1} | t-s | + \sum_{j>j_{0}} 2^{-j \alpha} \| u \|_{\alpha -1} \\
     & \lesssim ( 2^{j_{0} ( 1- \alpha )} | t-s | +2^{-j_{0} \alpha} ) \| u \|_{\alpha -1} \simeq | t-s |^{\alpha} \| u \|_{\alpha -1} .
  \end{align*}
  Uniqueness is easy since a distribution with zero derivative is necessarily a constant function.
\end{proof}



\section{Yet another commutator estimate}\label{sec:convolution commutators}

Here we prove that it is possible to ``pull the time integral inside the paraproduct'':

\begin{lemma}\label{lemma:controlled old and new}
  Let $\alpha , \beta \in ( 0,1 )$ with $\alpha + \beta <1$. Let $u \in \CC^\alpha(\R, \R^{d \times n})$ and $v \in \CC^{\beta}(\R, \R^n)$. Then for all $s,t \in \mathbb{R}$ with $| t-s | \leqslant 1$, we have
  \[
     \left| \int_{s}^{t} ( u \prec \partial_{t} v ) ( r ) \mathd r-u ( s ) ( v ( t ) -v ( s ) ) \right| \lesssim | t-s |^{\alpha + \beta} \| u  \|_{\alpha} \| v \|_{\beta} ,
  \]
  where we write $\int_{s}^{t} f ( r ) \mathd r= \int_{0}^{t} f ( r ) \mathd r- \int_{0}^{s} f ( r ) \mathd r$.
\end{lemma}

\begin{proof}
  Fix $s,t \in \mathbb{R}$ with $| s-t | \leqslant 1$. We can rewrite
  \[
     \int_{s}^{t} ( u \prec \partial_{t} v ) ( r ) \mathd r-u ( s ) ( v ( t ) -v ( s ) ) = \sum_{j} \int_{s}^{t} [ S_{j-1} u ( r ) -u ( s ) ] \partial_{r} \Delta_{j} v ( r ) \mathd r.
  \]
  We will use two different estimates, one for large $j$ and one for small $j$. First note that
  \begin{align*}
     \left| \int_{s}^{t} [ S_{j-1} u ( r ) -u ( s ) ] \partial_{r} \Delta_{j} v ( r ) \mathd r \right| & \leqslant \left| \int_{s}^{t} [ S_{j-1} u ( r ) -S_{j-1} u ( s ) ]  \partial_{r} \Delta_{j} v ( r ) \mathd r \right| \\
     &\qquad + \left| \int_{s}^{t} [ S_{j-1} u ( s ) -u ( s ) ] \partial_{r} \Delta_{j} v ( r ) \mathd r  \right| .
  \end{align*}
  Now $| \Delta_{i} u ( r ) - \Delta_{i} u ( s ) | \lesssim 2^{i ( 1- \alpha )} | r-s | \| u \|_{\alpha}$, which follows for example by interpolating the two inequalities $| \Delta_{i} u ( r ) - \Delta_{i} u ( s ) | \lesssim \| u \|_{\alpha} | r-s |^{\alpha}$ and $| \Delta_{i} u ( r ) - \Delta_{i} u ( s ) | \lesssim 2^{-i \alpha} \| u \|_{\alpha}$. This leads to
  \begin{align}\label{eq:controlled old and new pr1} \nonumber
     \left| \int_{s}^{t} [ S_{j-1} u ( r ) -u ( s ) ] \partial_{r} \Delta_{j} v ( r ) \mathd r \right| &\lesssim \left( 2^{j ( 1- \alpha )} \int_{s}^{t}  | r-s | 2^{j ( 1- \beta )} \mathd r+ \int_{s}^{t} 2^{-j \alpha} 2^{j ( 1- \beta )} \mathd r \right)\! \| u \|_{\alpha} \| v \|_{\beta} \\
    & \lesssim ( 2^{j ( 2- \alpha - \beta)} | t-s |^{2} +2^{j ( 1- \alpha - \beta )} | t-s | ) \| u \|_{\alpha} \| v \|_{\beta} .
  \end{align}
  On the other side, it follows from integration by parts that
  \[
     \left| \int_{s}^{t} [ S_{j-1} u ( r ) -S_{j-1} u ( s ) ] \partial_{r}  \Delta_{j} v ( r ) \mathd r \right| \leqslant | S_{j-1} u ( t )  \Delta_{j} v ( t ) | + \left| \int_{s}^{t} \partial_{r} S_{j-1} u ( r )    \Delta_{j} v ( r ) \mathd r \right|.
  \]
  The first term on the right hand side can be bounded by $\| u \|_{\alpha} \| v \|_{\beta} 2^{-j}$. For the second term, we apply
  Lemma~\ref{lemma:integral} to obtain
  \[
     \left| \int_{s}^{t} \partial_{r} S_{j-1} u ( r ) \Delta_{j} v ( r ) \mathd r \right| \lesssim | t-s |^{\alpha + \beta - \varepsilon} \| \partial_{r} S_{j-1} u \Delta_{j} v \|_{\alpha + \beta -1- \varepsilon}
  \]
  for any $\varepsilon \in [ 0, \alpha + \beta )$. Since the Fourier transform
  of $\partial_{r} S_{j-1} u \Delta_{j} v$ is supported in an annulus of the
  form $2^{j} \mathcal{A}$, we get
  \[
     \| \partial_{r} S_{j-1} u \Delta_{j} v \|_{\alpha + \beta -1-  \varepsilon} \lesssim 2^{j ( \alpha + \beta -1- \varepsilon )} \|  \partial_{r} S_{j-1} u \Delta_{j} v \|_{L^{\infty}} \lesssim 2^{-j     \varepsilon} \| u \|_{\alpha} \| v \|_{\beta} .
  \]
  Thus, we have
  \begin{align}\label{eq:controlled old and new pr2} \nonumber
     &\left| \int_{s}^{t} [ S_{j-1} u ( r ) -S_{j-1} u ( s ) ] \partial_{r} \Delta_{j} v ( r ) \mathd r \right| + \left| \int_{s}^{t} [ S_{j-1} u ( s ) -u ( s ) ] \partial_{r} \Delta_{j} v ( r ) \mathd r \right| \\ \nonumber
     &\hspace{70pt} \lesssim ( 2^{-j} + | t-s |^{\alpha + \beta - \varepsilon} 2^{-j  \varepsilon} ) \| u \|_{\alpha} \| v \|_{\beta} + \| S_{j-1} u-u \|_{L^{\infty}} | \Delta_{j} v ( t ) - \Delta_{j} v ( s ) | \\
     &\hspace{70pt} \lesssim ( 2^{-j} + | t-s |^{\alpha + \beta - \varepsilon} 2^{-j \varepsilon} +2^{-j \alpha} 2^{-j \beta} ) \|  u \|_{\alpha} \| v \|_{\beta} .
  \end{align}
  Let now $j_{0} \in \mathbb{N}$ be such that $2^{-j_{0}} \leqslant | t-s | <2^{-j_{0} +1}$. We use estimate~{\eqref{eq:controlled old and new pr1}} for
  $j \leqslant j_{0}$ and~{\eqref{eq:controlled old and new pr2}} for
  $j>j_{0}$ to obtain
  \begin{align*}
     \left| \int_{s}^{t} ( u \prec \partial_{t} v ) ( r ) \mathd r-u ( s ) ( v ( t ) -v ( s ) ) \right| & \lesssim \sum_{j \leqslant j_{0}} ( 2^{j ( 2- \alpha - \beta )} | t-s |^{2} +2^{j ( 1- \alpha - \beta )} | t-s | ) \| u \|_{\alpha} \| v \|_{\beta} \\
     &\quad + \sum_{j>j_{0}} ( 2^{-j} + | t-s |^{\alpha + \beta - \varepsilon} 2^{-j \varepsilon} +2^{-j \alpha} 2^{-j \beta} ) \| u \|_{\alpha} \| v \|_{\beta} \\
%  \[ \lesssim \| u \|_{\alpha} \| v \|_{\beta} ( 2^{j_{0} ( 2- \alpha - \beta
%     )} | t-s |^{2} +2^{j_{0} ( 1- \alpha - \beta )} | t-s | +2^{-j_{0}} + |
%     t-s |^{\alpha + \beta - \varepsilon} 2^{-j_{0} \varepsilon} +2^{-j_{0}
%     \alpha} 2^{-j_{0} \beta} ) \]
     & \simeq \| u \|_{\alpha} \| v \|_{\beta} | t-s |^{\alpha + \beta},
  \end{align*}
  where we used that $\alpha + \beta <1$ and $| t-s | \leqslant 1$.
\end{proof}




%\bibliographystyle{alpha}
%\bibliographystyle{plainnat}
%\bibliographystyle{plain}
%\bibliography{bib/all}


\begin{thebibliography}{FGGR12}

\bibitem[BCD11]{Bahouri2011}
Hajer Bahouri, Jean-Yves Chemin, and Raphael Danchin.
\newblock {\em {Fourier analysis and nonlinear partial differential
  equations}}.
\newblock Springer, 2011.

\bibitem[BGN13]{Brzezniak2010}
Zdzis{\l}aw Brze{\'z}niak, Massimiliano Gubinelli, and Misha Neklyudov.
\newblock {Global evolution of random vortex filament equation}.
\newblock {\em Nonlinearity}, 26(9):2499, 2013.

\bibitem[BGR05]{Bessaih2005}
Hakima Bessaih, Massimiliano Gubinelli, and Francesco Russo.
\newblock {The evolution of a random vortex filament}.
\newblock {\em Ann. Probab.}, 33(5):1825--1855, 2005.

\bibitem[Bon81]{Bony1981}
Jean-Michel Bony.
\newblock {Calcul symbolique et propagation des singularites pour les
  {\'e}quations aux d{\'e}riv{\'e}es partielles non lin{\'e}aires}.
\newblock {\em Ann. Sci. {\'E}c. Norm. Sup{\'e}r. (4)}, 14:209--246, 1981.

\bibitem[CC13]{Catellier2013}
Remi Catellier and Khalil Chouk.
\newblock Stochastic quantization in $\mathbb{T}^3$ and paracontrolled
  distribution.
\newblock preprint, 2013.

\bibitem[CF09]{Caruana2009}
Michael Caruana and Peter Friz.
\newblock {Partial differential equations driven by rough paths}.
\newblock {\em J. Differential Equations}, 247(1):140--173, 2009.

\bibitem[CFO11]{Caruana2011}
Michael Caruana, Peter~K. Friz, and Harald Oberhauser.
\newblock {A (rough) pathwise approach to a class of non-linear stochastic
  partial differential equations}.
\newblock {\em Annales de {l'Institut} Henri Poincare {(C)} Non Linear
  Analysis}, 28(1):27--46, January 2011.

\bibitem[CG06]{Chemin2006}
Jean-Yves Chemin and Isabelle Gallagher.
\newblock {On the global wellposedness of the 3-{D} {N}avier-{S}tokes equations
  with large initial data}.
\newblock {\em Ann. Sci. {\'E}cole Norm. Sup. (4)}, 39(4):679--698, 2006.

\bibitem[CG13]{Chouk2013}
Khalil Chouk and Massimiliano Gubinelli.
\newblock Rough sheets.
\newblock preprint, 2013.

\bibitem[DF12]{Diehl2012}
Joscha Diehl and Peter Friz.
\newblock {Backward stochastic differential equations with rough drivers}.
\newblock {\em Ann. Probab.}, 40(4):1715--1758, 2012.

\bibitem[DGT12]{Deya2012}
Aur{\'e}lien Deya, Massimiliano Gubinelli, and Samy Tindel.
\newblock {Non-linear rough heat equations}.
\newblock {\em Probab. Theory Related Fields}, 153(1-2):97--147, 2012.

\bibitem[FGGR12]{Friz2012}
Peter~K Friz, Benjamin Gess, Archil Gulisashvili, and Sebastian Riedel.
\newblock {Spatial rough path lifts of stochastic convolutions}.
\newblock {\em arXiv preprint arXiv:1211.0046}, 2012.

\bibitem[FO11]{Friz2011b}
Peter Friz and Harald Oberhauser.
\newblock {On the splitting-up method for rough (partial) differential
  equations}.
\newblock {\em J. Differential Equations}, 251(2):316--338, 2011.

\bibitem[FV10]{Friz2010}
Peter Friz and Nicolas Victoir.
\newblock {\em {Multidimensional stochastic processes as rough paths. Theory
  and applications}}.
\newblock Cambridge University Press, 2010.

\bibitem[GIP13]{Gubinelli2013}
Massimiliano Gubinelli, Peter Imkeller, and Nicolas Perkowski.
\newblock A {F}ourier approach to pathwise integration.
\newblock in preparation, 2013.

\bibitem[GLT06]{Gubinelli2006}
Massimiliano Gubinelli, Antoine Lejay, and Samy Tindel.
\newblock {Young integrals and {SPDE}s}.
\newblock {\em Potential Anal.}, 25(4):307--326, 2006.

\bibitem[Gub04]{Gubinelli2004}
Massimiliano Gubinelli.
\newblock {Controlling rough paths}.
\newblock {\em J. Funct. Anal.}, 216(1):86--140, nov 2004.

\bibitem[Gub12]{Gubinelli2012a}
Massimiliano Gubinelli.
\newblock {Rough solutions for the periodic {K}orteweg--de {V}ries equation}.
\newblock {\em Commun. Pure Appl. Anal.}, 11(2):709--733, 2012.

\bibitem[Gub13]{Gubinelli-KPZ-2013}
Massimiliano Gubinelli.
\newblock {KPZ} reloaded.
\newblock in preparation, 2013.

\bibitem[Hai11]{Hairer2011}
Martin Hairer.
\newblock {Rough stochastic {PDE}s}.
\newblock {\em Comm. Pure Appl. Math.}, 64(11):1547--1585, 2011.

\bibitem[Hai13a]{Hairer2013a}
Martin Hairer.
\newblock {A theory of regularity structures}.
\newblock {\em arXiv preprint arXiv:1303.5113}, 2013.

\bibitem[Hai13b]{Hairer2013b}
Martin Hairer.
\newblock {Solving the KPZ equation}.
\newblock {\em Ann. Math.}, 178(2):559--664, 2013.

\bibitem[HMW12]{Hairer2012}
Martin Hairer, Jan Maas, and Hendrik Weber.
\newblock {Approximating rough stochastic PDEs}.
\newblock {\em arXiv preprint arXiv:1202.3094}, 2012.

\bibitem[HW13]{Hairer2013}
Martin Hairer and Hendrik Weber.
\newblock {Rough {B}urgers-like equations with multiplicative noise}.
\newblock {\em Probab. Theory Related Fields}, 155(1-2):71--126, 2013.

\bibitem[Jan97]{Janson1997}
Svante Janson.
\newblock {\em {Gaussian {H}ilbert spaces}}, volume 129 of {\em {Cambridge
  Tracts in Mathematics}}.
\newblock Cambridge University Press, Cambridge, 1997.

\bibitem[KPZ86]{Kardar1986}
Mehran Kardar, Giorgio Parisi, and Yi-Cheng Zhang.
\newblock {Dynamic scaling of growing interfaces}.
\newblock {\em Physical Review Letters}, 56(9):889--892, 1986.

\bibitem[LCL07]{Lyons2007}
Terry~J. Lyons, Michael Caruana, and Thierry L{\'e}vy.
\newblock {\em {Differential equations driven by rough paths}}, volume 1908 of
  {\em {Lecture Notes in Mathematics}}.
\newblock Springer, Berlin, 2007.

\bibitem[LQ02]{Lyons2002}
Terry Lyons and Zhongmin Qian.
\newblock {\em {System control and rough paths}}.
\newblock Oxford University Press, 2002.

\bibitem[Lyo98]{Lyons1998}
Terry~J. Lyons.
\newblock {Differential equations driven by rough signals}.
\newblock {\em Rev. Mat. Iberoam.}, 14(2):215--310, 1998.

\bibitem[NT11]{Nualart2011}
David Nualart and Samy Tindel.
\newblock {A construction of the rough path above fractional {B}rownian motion
  using {V}olterra's representation}.
\newblock {\em Ann. Probab.}, 39(3):1061--1096, 2011.

\bibitem[Per13]{Perkowski2013Thesis}
Nicolas Perkowski.
\newblock {\em Studies of Robustness in Stochastic Analysis and Mathematical
  Finance}.
\newblock PhD thesis, Humboldt-Universit{\"a}t zu Berlin, 2013.

\bibitem[ST87]{Schmeisser1987}
Hans-J{\"u}rgen Schmeisser and Hans Triebel.
\newblock {\em {Topics in {F}ourier analysis and function spaces}}, volume~42.
\newblock Akademische Verlagsgesellschaft Geest \& Portig K.-G., Leipzig, 1987.

\bibitem[Tei11]{Teichmann2011}
Josef Teichmann.
\newblock {Another approach to some rough and stochastic partial differential
  equations}.
\newblock {\em Stoch. Dyn.}, 11(2-3):535--550, 2011.

\bibitem[Tri06]{Triebel2006}
Hans Triebel.
\newblock {\em {Theory of function spaces. {III}}}, volume 100 of {\em
  {Monographs in Mathematics}}.
\newblock Birkh{\"a}user Verlag, Basel, 2006.

\bibitem[Unt10a]{Unterberger2010a}
J{\'e}r{\'e}mie Unterberger.
\newblock {A rough path over multidimensional fractional Brownian motion with
  arbitrary Hurst index by Fourier normal ordering}.
\newblock {\em Stochastic Processes and their Applications}, 120(8):1444--1472,
  2010.

\bibitem[Unt10b]{Unterberger2010}
J{\'e}r{\'e}mie Unterberger.
\newblock {H{\"o}lder-Continuous Rough Paths by Fourier Normal Ordering}.
\newblock {\em Comm. Math. Phys.}, 298(1):1--36, 2010.

\end{thebibliography}


\end{document}
